\makeatletter
\def\input@path{{./}{../../../preamble/}{book/preamble/}}
\makeatother
\input{preamble_exams.tex}
\input{highlights/highlights.tex}
% \GTenHighlightsfalse % Uncomment to disable instructional highlights.
\input{highlights/specific_highlights.tex}
\CoordinatedTermTwoPalette
% Local semantic roles: indicator, success probability, complementary probability,
% monetary consequence, transformation function, and expected monetary value. Shared rendering and toggle remain authoritative.
\newcommand{\Indicator}[1]{\GTenHLValue{#1}}
\newcommand{\SuccessProb}[1]{\HLSuccessProbability{#1}}
\newcommand{\FailureProb}[1]{\HLComplementProbability{#1}}
\newcommand{\Money}[1]{\GTenHLSecond{#1}}
\newcommand{\TransformationFunction}[1]{\HLGtenCThreeFunction{#1}}
\newcommand{\ExpectedValue}[1]{\HLExpectedValue{#1}}
% Dark inks from the marker families keep graph strokes and text readable.
\ifGTenHighlights
  \colorlet{CThreeExpectedInk}{ExpectedValueColor!55!black}
  \colorlet{CThreeProbabilityInk}{MathHLSuccessColor!50!black}
  \newcommand{\CThreeGraphLegendShift}{-\baselineskip}
\else
  \colorlet{CThreeExpectedInk}{black}
  \colorlet{CThreeProbabilityInk}{black}
  \newcommand{\CThreeGraphLegendShift}{0pt}
\fi
\usepackage{needspace}
\hypersetup{pdftitle={Grade 10 Term 2 Discrete Probability}}
\setlength{\emergencystretch}{2em}
\title{Grade 10, Term 2, Lesson 1, Criterion C3\\Bernoulli Models and Individual Break-Even Analysis}
\author{}
\date{}
% Corresponding alternatives share a row; all inherited spacing is retained.
\NewDocumentCommand{\PracticeColumnPair}{+m +m}{%
\par\noindent\begin{minipage}[t]{.48\linewidth}\raggedright #1\end{minipage}\hfill
\begin{minipage}[t]{.48\linewidth}\raggedright #2\end{minipage}\par}
% Local companion graphs. Both axes use the same physical probability unit.
\usetikzlibrary{calc,backgrounds}
\newcommand{\CThreePlain}[1]{\begingroup\GTenHighlightsfalse #1\endgroup}
\ifGTenHighlights
  \colorlet{CThreeMassInk}{GTenHLProbability!65!black}
  \colorlet{CThreeCDFInk}{MathHLCumulativeColor!65!black}
\else
  \colorlet{CThreeMassInk}{black}
  \colorlet{CThreeCDFInk}{black}
\fi
\colorlet{CThreeShadowInk}{CThreeMassInk!32}
\tikzset{c3arrow/.style={gray!75,line width=.35pt,-{Stealth[length=4pt]}},
 c3card/.style={draw=gray!65,rounded corners=2pt,fill=white,inner sep=4pt},
 c3flow/.style={c3arrow,rounded corners=3pt}}
% variable, state-0 value, state-1 value, initial p, units, value role.
\newcommand{\CThreeDistributionPair}[6]{%
\begingroup
\tikzset{gten marker/.append style={inner xsep=.35mm,inner ysep=.15mm}}
\pgfmathsetmacro{\cThreeLow}{min(#2,#3)}%
\pgfmathsetmacro{\cThreeHigh}{max(#2,#3)}%
\pgfmathsetmacro{\cThreeQ}{ifthenelse(#2<#3,1-(#4),#4)}%
\pgfmathsetmacro{\cThreeR}{1-\cThreeQ}%
\pgfmathsetmacro{\cThreeMin}{\cThreeLow-.28*(\cThreeHigh-\cThreeLow)}%
\pgfmathsetmacro{\cThreeMax}{\cThreeHigh+.28*(\cThreeHigh-\cThreeLow)}%
\begin{tikzpicture}
\pgfplotsset{c3pair/.style={width=\dimexpr\linewidth-12mm\relax,height=1.45cm,
 scale only axis,xmin=\cThreeMin,xmax=\cThreeMax,ymin=-.04,ymax=1.14,
 axis lines=left,xtick={\cThreeLow,\cThreeHigh},ytick={0,1},
 tick label style={font=\footnotesize},label style={font=\footnotesize},
 xticklabel={$#6{\pgfmathprintnumber{\tick}}$},clip=false}}
\begin{axis}[c3pair,name=c3top,anchor=south west,
 ylabel={$\HLCumulativeProbability{F_{#1}}$},
 yticklabel={$\HLCumulativeProbability{\pgfmathprintnumber{\tick}}$}]
\addplot[CThreeShadowInk,densely dashed,line width=.9pt] coordinates
 {(\cThreeLow,0)(\cThreeLow,\cThreeQ)};
\addplot[CThreeShadowInk,densely dashed,line width=.9pt] coordinates
 {(\cThreeHigh,\cThreeQ)(\cThreeHigh,1)};
\coordinate (c3shadow0) at (axis cs:\cThreeLow,{\cThreeQ/2});
\coordinate (c3shadow1) at (axis cs:\cThreeHigh,{\cThreeQ+\cThreeR/2});
\addplot[CThreeCDFInk,thick,<-] coordinates {(\cThreeMin,0)(\cThreeLow,0)};
\addplot[CThreeCDFInk,thick] coordinates {(\cThreeLow,\cThreeQ)(\cThreeHigh,\cThreeQ)};
\addplot[CThreeCDFInk,thick,->] coordinates {(\cThreeHigh,1)(\cThreeMax,1)};
\addplot[only marks,mark=*,mark size=1.4pt,CThreeCDFInk,mark options={fill=white}]
 coordinates {(\cThreeLow,0)(\cThreeHigh,\cThreeQ)};
\addplot[only marks,mark=*,mark size=1.4pt,CThreeCDFInk]
 coordinates {(\cThreeLow,\cThreeQ)(\cThreeHigh,1)};
\node[font=\footnotesize,anchor=\ifdim\cThreeQ pt>.7pt north\else south\fi,
 fill=white,inner sep=1pt] at (axis cs:{(\cThreeLow+\cThreeHigh)/2},\cThreeQ)
 {$\HLCumulativeProbability{\pgfmathprintnumber[fixed,precision=2]{\cThreeQ}}$};
\end{axis}
\begin{axis}[c3pair,at={(c3top.south west)},anchor=north west,yshift=-.85cm,
 ylabel={$\HLProbability{m_{#1}}$},xlabel={#5},
 yticklabel={$\HLProbability{\pgfmathprintnumber{\tick}}$}]
\addplot[CThreeMassInk,thick] coordinates {(\cThreeLow,0)(\cThreeLow,\cThreeQ)};
\addplot[CThreeMassInk,thick] coordinates {(\cThreeHigh,0)(\cThreeHigh,\cThreeR)};
\addplot[only marks,mark=*,mark size=1.5pt,CThreeMassInk]
 coordinates {(\cThreeLow,\cThreeQ)(\cThreeHigh,\cThreeR)};
\node[font=\footnotesize,anchor=south,fill=white,inner sep=1pt]
 at (axis cs:\cThreeLow,\cThreeQ) {$\HLProbability{\pgfmathprintnumber[fixed,precision=2]{\cThreeQ}}$};
\node[font=\footnotesize,anchor=south,fill=white,inner sep=1pt]
 at (axis cs:\cThreeHigh,\cThreeR) {$\HLProbability{\pgfmathprintnumber[fixed,precision=2]{\cThreeR}}$};
\coordinate (c3mass0) at (axis cs:\cThreeLow,{\cThreeQ/2});
\coordinate (c3mass1) at (axis cs:\cThreeHigh,{\cThreeR/2});
\end{axis}
\begin{pgfonlayer}{background}
\draw[c3arrow] (c3mass0) .. controls ([xshift=2mm]$(c3mass0)!.33!(c3shadow0)$)
 and ([xshift=2mm]$(c3mass0)!.67!(c3shadow0)$) .. (c3shadow0);
\draw[c3arrow] (c3mass1) .. controls ([xshift=2mm]$(c3mass1)!.33!(c3shadow1)$)
 and ([xshift=2mm]$(c3mass1)!.67!(c3shadow1)$) .. (c3shadow1);
\end{pgfonlayer}
\end{tikzpicture}\endgroup}
% i, state-0 payoff, state-1 payoff, exact root, initial p.
\newcommand{\CThreeExpectationGraph}[5]{%
\begingroup
\tikzset{gten marker/.append style={inner xsep=.35mm,inner ysep=.15mm}}
\pgfmathsetmacro{\cThreeInitial}{#2+(#3-#2)*(#5)}%
\pgfmathsetmacro{\cThreeP}{#5}%
\def\cThreeTickAnchor{east}%
\pgfmathparse{min(abs(#2),abs(#3))/abs(#3-#2)<.15}%
\ifdim\pgfmathresult pt>0pt
 \def\cThreeTickAnchor{\ifdim\tick pt=0pt south east\else\ifdim\tick pt<0pt north east\else east\fi\fi}%
\fi
\begin{tikzpicture}\begin{axis}[
 width=\dimexpr\linewidth-14mm\relax,height=3.6cm,scale only axis,
 xmin=0,xmax=1,ymin={min(#2,#3)-.22*abs(#3-#2)},
 ymax={max(#2,#3)+.25*abs(#3-#2)},axis lines=left,
 xtick={0,1},ytick={#2,0,#3},grid=major,grid style={gray!15},
 yticklabel style={anchor=\cThreeTickAnchor},
 tick label style={font=\footnotesize},label style={font=\footnotesize},
 xlabel={$\SuccessProb{p}$},xlabel style={yshift=-3mm},
 ylabel={$\ExpectedValue{g_{#1}}$ (\ExpectedValue{USD})},
 yticklabel={$\ifdim\tick pt=0pt\ExpectedValue{0}\else\Money{\pgfmathprintnumber{\tick}}\fi$},
 xticklabel={$\SuccessProb{\pgfmathprintnumber{\tick}}$},clip=false]
\addplot[black,dashed,domain=0:1,samples=2]{0};
\addplot[CThreeExpectedInk,thick,domain=0:1,samples=2]{#2+(#3-#2)*x};
\addplot[black,only marks,mark=*,mark size=1.5pt] coordinates {(0,#2)(1,#3)(#4,0)};
\addplot[gray,dotted] coordinates {(#4,0)(#4,{min(#2,#3)-.22*abs(#3-#2)})};
\node[font=\footnotesize,anchor=north,yshift=-2pt] at
 (axis cs:#4,{min(#2,#3)-.22*abs(#3-#2)}) {$\SuccessProb{p_{#1}^*}$};
\addplot[CThreeExpectedInk,only marks,mark=square*,mark size=2pt]
 coordinates {(#5,\cThreeInitial)};
\node[font=\footnotesize,anchor=\ifdim#3pt>#2pt south east\else south west\fi,
 inner sep=2pt,fill=white] at (axis cs:#5,\cThreeInitial) {\ExpectedValue{initial}};
\node[font=\small,fill=white,inner sep=1pt] at
 (axis cs:{(#4)/2},{max(#2,#3)+.13*abs(#3-#2)})
 {$\ExpectedValue{\ifdim#3pt>#2pt -\else +\fi}$};
\node[font=\small,fill=white,inner sep=1pt] at
 (axis cs:{(1+(#4))/2},{max(#2,#3)+.13*abs(#3-#2)})
 {$\ExpectedValue{\ifdim#3pt>#2pt +\else -\fi}$};
\end{axis}\end{tikzpicture}\endgroup}
\newcommand{\CThreeGraphReading}{%
\small\textbf{Same distribution: CDF above, PMF below.}
Solid \HLProbability{stems} have filled endpoints. Pale dashed shadows translate
each \HLProbability{mass} to its \HLCumulativeProbability{CDF jump}; thin arrows link their midpoints.
Within a pair, horizontal limits and physical probability scales are identical.
Open circles exclude an endpoint; filled circles include it, making each
\HLCumulativeProbability{CDF} right-continuous. Between variables, arrows show the
\TransformationFunction{induced distribution}: apply \TransformationFunction{$f_i$}
to state values and retain their probabilities, rather than applying it to CDF heights.
At $\SuccessProb{p}=\SuccessProb{0}$ or $\SuccessProb{1}$, zero-mass jumps disappear.}

\newcommand{\CThreeTaskOne}{%
\textbf{Task 1 --- Establish the binary model and monetary transformations.}
\textbf{1.} \textbf{Binary model and monetary transformations.} Present the shared observation, period, \HLMathEvent{event} $\HLMathEvent{E}$, \HLMathEvent{complement} $\HLMathEvent{E^c}$, \Indicator{indicator} $\Indicator{X}$, \Indicator{Bernoulli model} and initial $\SuccessProb{p}=P(\HLMathEvent{E})$. Then give each \Money{monetary variable}'s \Money{state payoffs} in \Money{USD} and derive $\Money{Y_i}=\TransformationFunction{f_i}(\Indicator{X})=\Money{f_i(0)}+[\Money{f_i(1)}-\Money{f_i(0)}]\Indicator{X}$.
\textbf{2.} \textbf{Supports, \HLProbability{PMFs}, and ordered \HLCumulativeProbability{CDFs}.} Give the \Indicator{indicator}'s set-builder support, complete \HLProbability{PMF} and piecewise \HLCumulativeProbability{CDF}, including support cases at $\SuccessProb{p}=0,1$. Then give the \Money{monetary supports}, complete \HLProbability{PMFs} (zero elsewhere) and numerically ordered piecewise \HLCumulativeProbability{CDFs}, with endpoint support cases.
\textbf{3.} \textbf{Initial probabilities and loss \HLMathEvent{events}.} State the initial \Indicator{indicator} probabilities. For each $\Money{Y_i}$, give both \HLProbability{initial masses}, identify $\HLMathEvent{\{Y_i<0\}}$ in terms of $\Indicator{X}$ and $\HLMathEvent{E}$ or $\HLMathEvent{E^c}$, calculate its probability and interpret it in context. All alternatives share $\Indicator{X}$; no independence is assumed.
}
\newcommand{\CThreeTaskTwo}{%
\textbf{Task 2 --- Individual \ExpectedValue{expected values}.}
\textbf{1.} \textbf{Initial \ExpectedValue{expected values}.} Calculate and interpret every \ExpectedValue{monetary expectation} at the supplied probability before sensitivity analysis.
\textbf{2.} \textbf{\ExpectedValue{Expectation formulas} and sign intervals.} Hold \Money{payoffs} fixed and vary only $\SuccessProb{p}\in\left[0,1\right]$. For each $\Money{Y_i}$, derive $\ExpectedValue{\mathbb E}[\Money{Y_i}]=\SuccessProb{p} \Money{f_i(1)}+(\FailureProb{1-p})\Money{f_i(0)}$, solve its exact \SuccessProb{break-even equation}, and solve both strict sign inequalities with correct directions and endpoints. Explain \ExpectedValue{positive, negative and zero expectation}.
\textbf{3.} \textbf{Graphs.} Give a separate labelled graph for each variable, showing its \ExpectedValue{expectation line}, zero line, endpoint \Money{payoffs} and exact \SuccessProb{break-even probability}. Finish with endpoint and root checks, and distinguish \ExpectedValue{expected profit} from \ExpectedValue{realized profit}.
}
\newcommand{\CThreeIntroTaskOne}{\textbf{Task 1 --- Modelling and distributions.} Model the shared \Indicator{indicator} and both monetary transformations; give their supports, \HLProbability{PMFs}, ordered \HLCumulativeProbability{CDFs} and loss events.}
\newcommand{\CThreeIntroTaskTwo}{\textbf{Task 2 --- Expectations.} Calculate and interpret both initial \ExpectedValue{expected values}, derive their \ExpectedValue{expectation formulas}, solve their exact break-even equations and sign intervals, and graph each \ExpectedValue{expectation}.}
\newcommand{\CThreeStatementIntroOne}{%
Observe a student's art-print stall over one weekend. \HLMathEvent{High demand} means campus-market
attendance meets the organizer's fixed target; \HLMathEvent{low demand} means it falls below that target.
These states cover every outcome and apply to either hypothetical print-batch choice.
Let $\HLMathEvent{E}$ be \HLMathEvent{high demand} and
$\HLMathEvent{E^c}$ \HLMathEvent{low demand}. Define $\Indicator{X}=\Indicator{1}$ on $\HLMathEvent{E}$ and
$\Indicator{X}=\Indicator{0}$ on $\HLMathEvent{E^c}$, with $\SuccessProb{p}=P(\HLMathEvent{E})=\SuccessProb{0.4}$
and $\FailureProb{1-p}=\FailureProb{0.6}$. The \Indicator{indicator} is unitless.
After printing and stall costs, a small print batch gives $\Money{f_1(0)}=\Money{-20}$, $\Money{f_1(1)}=\Money{60}$ \Money{USD};
a large print batch gives $\Money{f_2(0)}=\Money{-40}$, $\Money{f_2(1)}=\Money{120}$ \Money{USD}.

\begin{center}\begin{tabular}{lrr}\toprule Alternative & \HLMathEvent{low demand} & \HLMathEvent{high demand}\\\midrule
$A_1$: small print batch & \Money{-20} & \Money{60}\\
$A_2$: large print batch & \Money{-40} & \Money{120}\\
\bottomrule\end{tabular}\end{center}
}
\newcommand{\CThreeModelIntroOne}{%
\clearpage\noindent\textbf{Task 1 summary: shared states and induced distributions.}\par\medskip
\textbf{Numerical illustration: introductory example 1.}\par
\begingroup\small\setlength{\parskip}{2pt}\tikzset{gten marker/.append style={inner xsep=.4mm,inner ysep=.2mm}}
\begin{center}\begin{tikzpicture}[x=1cm,y=1cm]
\node[c3card,anchor=north,text width=7cm,align=center] (ev1) at (4,0) {$\HLMathEvent{E}$: \HLMathEvent{high demand}, $\SuccessProb{P}(\HLMathEvent{E})=\SuccessProb{p}$};

\node[c3card,anchor=north,text width=7cm,align=center] (ev0) at (15,0) {$\HLMathEvent{E^c}$: \HLMathEvent{low demand}, $\FailureProb{P}(\HLMathEvent{E^c})=\FailureProb{1-p}$};

\node[c3card,anchor=north,text width=6cm,align=center] (ind1) at (4,-1.2) {$\Indicator{X(\omega)}=\Indicator{1}$, for $\HLMathEvent{\omega\in E}$};

\node[c3card,anchor=north,text width=6cm,align=center] (ind0) at (15,-1.2) {$\Indicator{X(\omega)}=\Indicator{0}$, for $\HLMathEvent{\omega\in E^c}$};

\draw[c3flow] (ev1) -- (ind1);\draw[c3flow] (ev0) -- (ind0);
\node[c3card,anchor=north,text width=16.2cm,align=center] (map1) at (9.5,-2.65) {\textbf{\Money{$Y_1$}: small print batch}\par 
$\Money{Y_1(\omega)}=\TransformationFunction{f_1}(\Indicator{X(\omega)})=\Money{f_1(1)}=\Money{60}\ \Money{\mathrm{USD}}$, for $\HLMathEvent{\omega\in E}$\par 
$\Money{Y_1(\omega)}=\TransformationFunction{f_1}(\Indicator{X(\omega)})=\Money{f_1(0)}=\Money{-20}\ \Money{\mathrm{USD}}$, for $\HLMathEvent{\omega\in E^c}$
};

\draw[c3flow] (ind1.west) -- (.15,-1.65) |- ([yshift=-7mm]map1.north west);
\draw[c3flow,dashed] (ind0.east) -- (18.85,-1.65) |- ([yshift=-11.5mm]map1.north east);
\node[c3card,anchor=north,text width=16.2cm,align=center] (map2) at (9.5,-4.45) {\textbf{\Money{$Y_2$}: large print batch}\par 
$\Money{Y_2(\omega)}=\TransformationFunction{f_2}(\Indicator{X(\omega)})=\Money{f_2(1)}=\Money{120}\ \Money{\mathrm{USD}}$, for $\HLMathEvent{\omega\in E}$\par 
$\Money{Y_2(\omega)}=\TransformationFunction{f_2}(\Indicator{X(\omega)})=\Money{f_2(0)}=\Money{-40}\ \Money{\mathrm{USD}}$, for $\HLMathEvent{\omega\in E^c}$
};

\draw[c3flow] (ind1.west) -- (.15,-1.65) |- ([yshift=-7mm]map2.north west);
\draw[c3flow,dashed] (ind0.east) -- (18.85,-1.65) |- ([yshift=-11.5mm]map2.north east);
\node[inner sep=1pt,anchor=north,text width=5.913333333333333cm,align=center] (head0) at (3.0666666666666664,-6.4) {$\Indicator{X}$};

\node[c3card,anchor=north,text width=5.913333333333333cm,align=center] (cdf0) at (3.0666666666666664,-7.050000000000001) {$\HLCumulativeProbability{F_{X}(t)}$\par $\begin{cases}\HLCumulativeProbability{0},&t<\Indicator{0},\\\FailureProb{1-p},&\Indicator{0}\leq t<\Indicator{1},\\\HLCumulativeProbability{1},&t\geq \Indicator{1}.\end{cases}$};

\node[c3card,anchor=north,text width=5.913333333333333cm,align=center] (pmf0) at (3.0666666666666664,-9.8) {$\HLProbability{m_{X}}$\par $\begin{array}{ll}\Indicator{0}:&\FailureProb{1-p}\\\Indicator{1}:&\SuccessProb{p}\end{array}$\par \HLProbability{zero elsewhere}};

\node[inner sep=0pt,anchor=north,text width=5.913333333333333cm,align=center] (graph0) at (3.0666666666666664,-12.100000000000001) {\CThreeDistributionPair{X}{0}{1}{2/5}{unitless}{\Indicator}};

\node[inner sep=1pt,anchor=north,text width=5.913333333333333cm,align=center] (head1) at (9.4,-6.4) {$\Money{Y_1}$ $\leftarrow\TransformationFunction{f_1}$};

\node[c3card,anchor=north,text width=5.913333333333333cm,align=center] (cdf1) at (9.4,-7.050000000000001) {$\HLCumulativeProbability{F_{Y_1}(t)}$\par $\begin{cases}\HLCumulativeProbability{0},&t<\Money{-20},\\\FailureProb{1-p},&\Money{-20}\leq t<\Money{60},\\\HLCumulativeProbability{1},&t\geq \Money{60}.\end{cases}$};

\node[c3card,anchor=north,text width=5.913333333333333cm,align=center] (pmf1) at (9.4,-9.8) {$\HLProbability{m_{Y_1}}$\par $\begin{array}{ll}\Money{-20}:&\FailureProb{1-p}\\\Money{60}:&\SuccessProb{p}\end{array}$\par \HLProbability{zero elsewhere}};

\node[inner sep=0pt,anchor=north,text width=5.913333333333333cm,align=center] (graph1) at (9.4,-12.100000000000001) {\CThreeDistributionPair{Y_1}{-20}{60}{2/5}{\Money{USD}}{\Money}};

\draw[c3flow] (cdf0.north) -- ++(0,.27) -| (cdf1.north);
\draw[c3flow] (pmf0.south) -- ++(0,-.30) -| (pmf1.south);
\draw[c3flow] (map1.east) -- (18.28,-3.3) -- (18.28,-6.1000000000000005) -| (head1.north);
\node[inner sep=1pt,anchor=north,text width=5.913333333333333cm,align=center] (head2) at (15.733333333333333,-6.4) {$\Money{Y_2}$ $\leftarrow\TransformationFunction{f_2}$};

\node[c3card,anchor=north,text width=5.913333333333333cm,align=center] (cdf2) at (15.733333333333333,-7.050000000000001) {$\HLCumulativeProbability{F_{Y_2}(t)}$\par $\begin{cases}\HLCumulativeProbability{0},&t<\Money{-40},\\\FailureProb{1-p},&\Money{-40}\leq t<\Money{120},\\\HLCumulativeProbability{1},&t\geq \Money{120}.\end{cases}$};

\node[c3card,anchor=north,text width=5.913333333333333cm,align=center] (pmf2) at (15.733333333333333,-9.8) {$\HLProbability{m_{Y_2}}$\par $\begin{array}{ll}\Money{-40}:&\FailureProb{1-p}\\\Money{120}:&\SuccessProb{p}\end{array}$\par \HLProbability{zero elsewhere}};

\node[inner sep=0pt,anchor=north,text width=5.913333333333333cm,align=center] (graph2) at (15.733333333333333,-12.100000000000001) {\CThreeDistributionPair{Y_2}{-40}{120}{2/5}{\Money{USD}}{\Money}};

\draw[c3flow] (cdf0.north) -- ++(0,.27) -| (cdf2.north);
\draw[c3flow] (pmf0.south) -- ++(0,-.30) -| (pmf2.south);
\draw[c3flow] (map2.east) -- (18.1,-5.1) -- (18.1,-5.98) -| (head2.north);
\end{tikzpicture}\end{center}
\textbf{Numerical graph illustration:} $\SuccessProb{p_0}=\SuccessProb{\frac{2}{5}}$. All displayed distribution formulas remain symbolic in $\SuccessProb{p}$.
\par\CThreeGraphReading\par
For $0<\SuccessProb{p}<1$, $\Indicator{S_X}=\{\Indicator{0},\Indicator{1}\}$; at $\SuccessProb{p}=0$ only $\Indicator{0}$ remains, and at $\SuccessProb{p}=1$ only $\Indicator{1}$ remains. Each monetary support retains $\Money{f_i(0)}$ at $\SuccessProb{p}=0$ and $\Money{f_i(1)}$ at $\SuccessProb{p}=1$.
\par\textbf{Loss events.} $\HLMathEvent{\{Y_1<0\}}=\HLMathEvent{\{X=0\}}=\HLMathEvent{E^c}$, probability $\FailureProb{1-p}$; $\HLMathEvent{\{Y_2<0\}}=\HLMathEvent{\{X=0\}}=\HLMathEvent{E^c}$, probability $\FailureProb{1-p}$.
\endgroup
}
\newcommand{\CThreeExpectIntroOne}{%
\clearpage\noindent\textbf{Task 2 summary: parallel expectation analyses.}\par\medskip
\textbf{Numerical illustration: introductory example 1.}\par
\begingroup\small\tikzset{gten marker/.append style={inner xsep=.4mm,inner ysep=.2mm}}
\begin{center}\begin{tikzpicture}[x=1cm,y=1cm]
\node[c3card,anchor=north,text width=17cm,align=center] (inputs) at (9.5,0) {Shared probability inputs: $\SuccessProb{p}$ and $\FailureProb{1-p}$; hold each alternative\textquotesingle s \Money{payoffs} fixed.};

\node[c3card,anchor=north,text width=17cm,align=center] (formula) at (9.5,-1.2) {$\ExpectedValue{g_i}(\SuccessProb{p})=\ExpectedValue{\mathbb E}[\Money{Y_i}]=\SuccessProb{p}\Money{f_i(1)}+(\FailureProb{1-p})\Money{f_i(0)}$};

\draw[c3flow] (inputs) -- (formula);
\node[c3card,anchor=north,text width=8.8cm,align=center] (pay1) at (4.75,-2.7) {\textbf{\Money{$Y_1$}: small print batch}\par $\Money{f_1(0)}=\Money{-20},\quad\Money{f_1(1)}=\Money{60}\ \Money{\mathrm{USD}}$\par $\ExpectedValue{g_1}(\SuccessProb{p})=\Money{-20}+\Money{80}\SuccessProb{p}$};

\draw[c3flow] (formula.south) -- ++(0,-.25) -| (pay1.north);
\node[c3card,anchor=north,text width=8.8cm,align=center] (initial1) at (4.75,-5.0) {\ExpectedValue{Initial expectation}\par $\ExpectedValue{g_1}\!\left(\SuccessProb{\frac{2}{5}}\right)=\ExpectedValue{12\ \mathrm{USD}}$\par \ExpectedValue{expected net gain} per stated period};

\node[c3card,anchor=north,text width=8.8cm,align=center] (root1) at (4.75,-7.15) {\SuccessProb{Exact break-even equation}\par $\Money{-20}+\Money{80}\SuccessProb{p}=\ExpectedValue{0}$\par $\SuccessProb{p_1^*}=\SuccessProb{\frac{1}{4}}$};

\node[c3card,anchor=north,text width=8.8cm,align=center] (sign1) at (4.75,-9.45) {$\ExpectedValue{g_1}>\ExpectedValue{0}:\ \SuccessProb{p}\in \left(\SuccessProb{\frac{1}{4}},\SuccessProb{1}\right]$\par $\ExpectedValue{g_1}=\ExpectedValue{0}:\ \SuccessProb{p}=\SuccessProb{\frac{1}{4}}$\par $\ExpectedValue{g_1}<\ExpectedValue{0}:\ \SuccessProb{p}\in \left[\SuccessProb{0},\SuccessProb{\frac{1}{4}}\right)$};

\node[inner sep=0pt,anchor=north,text width=8.8cm,align=center] (eg1) at (4.75,-12.0) {\CThreeExpectationGraph{1}{-20}{60}{1/4}{2/5}};

\draw[c3flow] (pay1.south) -- (initial1.north);
\draw[c3flow] (initial1.south) -- (root1.north);
\draw[c3flow] (root1.south) -- (sign1.north);
\draw[c3flow] (sign1.south) -- (eg1.north);
\node[c3card,anchor=north,text width=8.8cm,align=center] (pay2) at (14.25,-2.7) {\textbf{\Money{$Y_2$}: large print batch}\par $\Money{f_2(0)}=\Money{-40},\quad\Money{f_2(1)}=\Money{120}\ \Money{\mathrm{USD}}$\par $\ExpectedValue{g_2}(\SuccessProb{p})=\Money{-40}+\Money{160}\SuccessProb{p}$};

\draw[c3flow] (formula.south) -- ++(0,-.25) -| (pay2.north);
\node[c3card,anchor=north,text width=8.8cm,align=center] (initial2) at (14.25,-5.0) {\ExpectedValue{Initial expectation}\par $\ExpectedValue{g_2}\!\left(\SuccessProb{\frac{2}{5}}\right)=\ExpectedValue{24\ \mathrm{USD}}$\par \ExpectedValue{expected net gain} per stated period};

\node[c3card,anchor=north,text width=8.8cm,align=center] (root2) at (14.25,-7.15) {\SuccessProb{Exact break-even equation}\par $\Money{-40}+\Money{160}\SuccessProb{p}=\ExpectedValue{0}$\par $\SuccessProb{p_2^*}=\SuccessProb{\frac{1}{4}}$};

\node[c3card,anchor=north,text width=8.8cm,align=center] (sign2) at (14.25,-9.45) {$\ExpectedValue{g_2}>\ExpectedValue{0}:\ \SuccessProb{p}\in \left(\SuccessProb{\frac{1}{4}},\SuccessProb{1}\right]$\par $\ExpectedValue{g_2}=\ExpectedValue{0}:\ \SuccessProb{p}=\SuccessProb{\frac{1}{4}}$\par $\ExpectedValue{g_2}<\ExpectedValue{0}:\ \SuccessProb{p}\in \left[\SuccessProb{0},\SuccessProb{\frac{1}{4}}\right)$};

\node[inner sep=0pt,anchor=north,text width=8.8cm,align=center] (eg2) at (14.25,-12.0) {\CThreeExpectationGraph{2}{-40}{120}{1/4}{2/5}};

\draw[c3flow] (pay2.south) -- (initial2.north);
\draw[c3flow] (initial2.south) -- (root2.north);
\draw[c3flow] (root2.south) -- (sign2.north);
\draw[c3flow] (sign2.south) -- (eg2.north);
\end{tikzpicture}\end{center}
The dashed horizontal line is \ExpectedValue{zero}; filled circles show endpoint payoffs and the exact root. The square marks the \ExpectedValue{initial expectation}. The $+$ and $-$ labels connect the graph to its strict sign intervals.
\par \ExpectedValue{Positive expectation} means \ExpectedValue{expected net gain}; \ExpectedValue{negative expectation} means \ExpectedValue{expected net loss}; \ExpectedValue{zero expectation} means break-even. These are long-run averages per stated period, not guarantees of realized profit. Each branch uses its own payoffs and root.
\endgroup
}
\newcommand{\CThreeStatementIntroTwo}{%
Observe a freelancer hosting paid livestream sessions over one day. Let $\HLMathEvent{E}$ be \HLMathEvent{a home-internet disruption}
and $\HLMathEvent{E^c}$ \HLMathEvent{no disruption} that day. These are the only states. Compare hypothetical choices:
use home internet alone or buy mobile data as backup. All \Money{payoffs} are \Money{net gains or losses}
after platform fees, data charges and refunds. Set $\Indicator{X}=\Indicator{1}$ on $\HLMathEvent{E}$ and
$\Indicator{X}=\Indicator{0}$ on $\HLMathEvent{E^c}$. Here $\SuccessProb{p}=P(\HLMathEvent{E})=\SuccessProb{0.1}$,
$\FailureProb{1-p}=\FailureProb{0.9}$. Mathematical success can be financially undesirable.
With home internet alone, $\Money{f_1(0)}=\Money{100}$, $\Money{f_1(1)}=\Money{-20}$ \Money{USD}.
With mobile-data backup, $\Money{f_2(0)}=\Money{60}$, $\Money{f_2(1)}=\Money{-10}$ \Money{USD}:
\begin{center}\begin{tabular}{lrr}\toprule Alternative & \HLMathEvent{no disruption} & \HLMathEvent{disruption}\\\midrule
$A_1$: home internet alone & \Money{100} & \Money{-20}\\
$A_2$: mobile-data backup & \Money{60} & \Money{-10}\\
\bottomrule\end{tabular}\end{center}
}
\newcommand{\CThreeModelIntroTwo}{%
\clearpage\noindent\textbf{Task 1 summary: shared states and induced distributions.}\par\medskip

\begingroup\small\setlength{\parskip}{2pt}\tikzset{gten marker/.append style={inner xsep=.4mm,inner ysep=.2mm}}
\begin{center}\begin{tikzpicture}[x=1cm,y=1cm]
\node[c3card,anchor=north,text width=7cm,align=center] (ev1) at (4,0) {$\HLMathEvent{E}$: \HLMathEvent{disruption}, $\SuccessProb{P}(\HLMathEvent{E})=\SuccessProb{p}$};

\node[c3card,anchor=north,text width=7cm,align=center] (ev0) at (15,0) {$\HLMathEvent{E^c}$: \HLMathEvent{no disruption}, $\FailureProb{P}(\HLMathEvent{E^c})=\FailureProb{1-p}$};

\node[c3card,anchor=north,text width=6cm,align=center] (ind1) at (4,-1.2) {$\Indicator{X(\omega)}=\Indicator{1}$, for $\HLMathEvent{\omega\in E}$};

\node[c3card,anchor=north,text width=6cm,align=center] (ind0) at (15,-1.2) {$\Indicator{X(\omega)}=\Indicator{0}$, for $\HLMathEvent{\omega\in E^c}$};

\draw[c3flow] (ev1) -- (ind1);\draw[c3flow] (ev0) -- (ind0);
\node[c3card,anchor=north,text width=16.2cm,align=center] (map1) at (9.5,-2.65) {\textbf{\Money{$Y_1$}: home internet alone}\par 
$\Money{Y_1(\omega)}=\TransformationFunction{f_1}(\Indicator{X(\omega)})=\Money{f_1(1)}=\Money{-20}\ \Money{\mathrm{USD}}$, for $\HLMathEvent{\omega\in E}$\par 
$\Money{Y_1(\omega)}=\TransformationFunction{f_1}(\Indicator{X(\omega)})=\Money{f_1(0)}=\Money{100}\ \Money{\mathrm{USD}}$, for $\HLMathEvent{\omega\in E^c}$
};

\draw[c3flow] (ind1.west) -- (.15,-1.65) |- ([yshift=-7mm]map1.north west);
\draw[c3flow,dashed] (ind0.east) -- (18.85,-1.65) |- ([yshift=-11.5mm]map1.north east);
\node[c3card,anchor=north,text width=16.2cm,align=center] (map2) at (9.5,-4.45) {\textbf{\Money{$Y_2$}: mobile-data backup}\par 
$\Money{Y_2(\omega)}=\TransformationFunction{f_2}(\Indicator{X(\omega)})=\Money{f_2(1)}=\Money{-10}\ \Money{\mathrm{USD}}$, for $\HLMathEvent{\omega\in E}$\par 
$\Money{Y_2(\omega)}=\TransformationFunction{f_2}(\Indicator{X(\omega)})=\Money{f_2(0)}=\Money{60}\ \Money{\mathrm{USD}}$, for $\HLMathEvent{\omega\in E^c}$
};

\draw[c3flow] (ind1.west) -- (.15,-1.65) |- ([yshift=-7mm]map2.north west);
\draw[c3flow,dashed] (ind0.east) -- (18.85,-1.65) |- ([yshift=-11.5mm]map2.north east);
\node[inner sep=1pt,anchor=north,text width=5.913333333333333cm,align=center] (head0) at (3.0666666666666664,-6.4) {$\Indicator{X}$};

\node[c3card,anchor=north,text width=5.913333333333333cm,align=center] (cdf0) at (3.0666666666666664,-7.050000000000001) {$\HLCumulativeProbability{F_{X}(t)}$\par $\begin{cases}\HLCumulativeProbability{0},&t<\Indicator{0},\\\FailureProb{1-p},&\Indicator{0}\leq t<\Indicator{1},\\\HLCumulativeProbability{1},&t\geq \Indicator{1}.\end{cases}$};

\node[c3card,anchor=north,text width=5.913333333333333cm,align=center] (pmf0) at (3.0666666666666664,-9.8) {$\HLProbability{m_{X}}$\par $\begin{array}{ll}\Indicator{0}:&\FailureProb{1-p}\\\Indicator{1}:&\SuccessProb{p}\end{array}$\par \HLProbability{zero elsewhere}};

\node[inner sep=0pt,anchor=north,text width=5.913333333333333cm,align=center] (graph0) at (3.0666666666666664,-12.100000000000001) {\CThreeDistributionPair{X}{0}{1}{1/10}{unitless}{\Indicator}};

\node[inner sep=1pt,anchor=north,text width=5.913333333333333cm,align=center] (head1) at (9.4,-6.4) {$\Money{Y_1}$ $\leftarrow\TransformationFunction{f_1}$};

\node[c3card,anchor=north,text width=5.913333333333333cm,align=center] (cdf1) at (9.4,-7.050000000000001) {$\HLCumulativeProbability{F_{Y_1}(t)}$\par $\begin{cases}\HLCumulativeProbability{0},&t<\Money{-20},\\\SuccessProb{p},&\Money{-20}\leq t<\Money{100},\\\HLCumulativeProbability{1},&t\geq \Money{100}.\end{cases}$};

\node[c3card,anchor=north,text width=5.913333333333333cm,align=center] (pmf1) at (9.4,-9.8) {$\HLProbability{m_{Y_1}}$\par $\begin{array}{ll}\Money{100}:&\FailureProb{1-p}\\\Money{-20}:&\SuccessProb{p}\end{array}$\par \HLProbability{zero elsewhere}};

\node[inner sep=0pt,anchor=north,text width=5.913333333333333cm,align=center] (graph1) at (9.4,-12.100000000000001) {\CThreeDistributionPair{Y_1}{100}{-20}{1/10}{\Money{USD}}{\Money}};

\draw[c3flow] (cdf0.north) -- ++(0,.27) -| (cdf1.north);
\draw[c3flow] (pmf0.south) -- ++(0,-.30) -| (pmf1.south);
\draw[c3flow] (map1.east) -- (18.28,-3.3) -- (18.28,-6.1000000000000005) -| (head1.north);
\node[inner sep=1pt,anchor=north,text width=5.913333333333333cm,align=center] (head2) at (15.733333333333333,-6.4) {$\Money{Y_2}$ $\leftarrow\TransformationFunction{f_2}$};

\node[c3card,anchor=north,text width=5.913333333333333cm,align=center] (cdf2) at (15.733333333333333,-7.050000000000001) {$\HLCumulativeProbability{F_{Y_2}(t)}$\par $\begin{cases}\HLCumulativeProbability{0},&t<\Money{-10},\\\SuccessProb{p},&\Money{-10}\leq t<\Money{60},\\\HLCumulativeProbability{1},&t\geq \Money{60}.\end{cases}$};

\node[c3card,anchor=north,text width=5.913333333333333cm,align=center] (pmf2) at (15.733333333333333,-9.8) {$\HLProbability{m_{Y_2}}$\par $\begin{array}{ll}\Money{60}:&\FailureProb{1-p}\\\Money{-10}:&\SuccessProb{p}\end{array}$\par \HLProbability{zero elsewhere}};

\node[inner sep=0pt,anchor=north,text width=5.913333333333333cm,align=center] (graph2) at (15.733333333333333,-12.100000000000001) {\CThreeDistributionPair{Y_2}{60}{-10}{1/10}{\Money{USD}}{\Money}};

\draw[c3flow] (cdf0.north) -- ++(0,.27) -| (cdf2.north);
\draw[c3flow] (pmf0.south) -- ++(0,-.30) -| (pmf2.south);
\draw[c3flow] (map2.east) -- (18.1,-5.1) -- (18.1,-5.98) -| (head2.north);
\end{tikzpicture}\end{center}
\textbf{Numerical graph illustration:} $\SuccessProb{p_0}=\SuccessProb{\frac{1}{10}}$. All displayed distribution formulas remain symbolic in $\SuccessProb{p}$.
\par\textbf{CDF above, PMF below.} Filled stems translate to pale dashed jump shadows; thin arrows link equal physical lengths. Open and filled circles show excluded and included endpoints. Distribution arrows retain state probabilities while mapping values, rather than applying $\TransformationFunction{f_i}$ to CDF heights. See \hyperref[ex:intro-one]{Example 1}.\par
For $0<\SuccessProb{p}<1$, $\Indicator{S_X}=\{\Indicator{0},\Indicator{1}\}$; at $\SuccessProb{p}=0$ only $\Indicator{0}$ remains, and at $\SuccessProb{p}=1$ only $\Indicator{1}$ remains. Each monetary support retains $\Money{f_i(0)}$ at $\SuccessProb{p}=0$ and $\Money{f_i(1)}$ at $\SuccessProb{p}=1$.
\par\textbf{Loss events.} $\HLMathEvent{\{Y_1<0\}}=\HLMathEvent{\{X=1\}}=\HLMathEvent{E}$, probability $\SuccessProb{p}$; $\HLMathEvent{\{Y_2<0\}}=\HLMathEvent{\{X=1\}}=\HLMathEvent{E}$, probability $\SuccessProb{p}$.
\endgroup
}
\newcommand{\CThreeExpectIntroTwo}{%
\clearpage\noindent\textbf{Task 2 summary: parallel expectation analyses.}\par\medskip

\begingroup\small\tikzset{gten marker/.append style={inner xsep=.4mm,inner ysep=.2mm}}
\begin{center}\begin{tikzpicture}[x=1cm,y=1cm]
\node[c3card,anchor=north,text width=17cm,align=center] (inputs) at (9.5,0) {Shared probability inputs: $\SuccessProb{p}$ and $\FailureProb{1-p}$; hold each alternative\textquotesingle s \Money{payoffs} fixed.};

\node[c3card,anchor=north,text width=17cm,align=center] (formula) at (9.5,-1.2) {$\ExpectedValue{g_i}(\SuccessProb{p})=\ExpectedValue{\mathbb E}[\Money{Y_i}]=\SuccessProb{p}\Money{f_i(1)}+(\FailureProb{1-p})\Money{f_i(0)}$};

\draw[c3flow] (inputs) -- (formula);
\node[c3card,anchor=north,text width=8.8cm,align=center] (pay1) at (4.75,-2.7) {\textbf{\Money{$Y_1$}: home internet alone}\par $\Money{f_1(0)}=\Money{100},\quad\Money{f_1(1)}=\Money{-20}\ \Money{\mathrm{USD}}$\par $\ExpectedValue{g_1}(\SuccessProb{p})=\Money{100}-\Money{120}\SuccessProb{p}$};

\draw[c3flow] (formula.south) -- ++(0,-.25) -| (pay1.north);
\node[c3card,anchor=north,text width=8.8cm,align=center] (initial1) at (4.75,-5.0) {\ExpectedValue{Initial expectation}\par $\ExpectedValue{g_1}\!\left(\SuccessProb{\frac{1}{10}}\right)=\ExpectedValue{88\ \mathrm{USD}}$\par \ExpectedValue{expected net gain} per stated period};

\node[c3card,anchor=north,text width=8.8cm,align=center] (root1) at (4.75,-7.15) {\SuccessProb{Exact break-even equation}\par $\Money{100}-\Money{120}\SuccessProb{p}=\ExpectedValue{0}$\par $\SuccessProb{p_1^*}=\SuccessProb{\frac{5}{6}}$};

\node[c3card,anchor=north,text width=8.8cm,align=center] (sign1) at (4.75,-9.45) {$\ExpectedValue{g_1}>\ExpectedValue{0}:\ \SuccessProb{p}\in \left[\SuccessProb{0},\SuccessProb{\frac{5}{6}}\right)$\par $\ExpectedValue{g_1}=\ExpectedValue{0}:\ \SuccessProb{p}=\SuccessProb{\frac{5}{6}}$\par $\ExpectedValue{g_1}<\ExpectedValue{0}:\ \SuccessProb{p}\in \left(\SuccessProb{\frac{5}{6}},\SuccessProb{1}\right]$};

\node[inner sep=0pt,anchor=north,text width=8.8cm,align=center] (eg1) at (4.75,-12.0) {\CThreeExpectationGraph{1}{100}{-20}{5/6}{1/10}};

\draw[c3flow] (pay1.south) -- (initial1.north);
\draw[c3flow] (initial1.south) -- (root1.north);
\draw[c3flow] (root1.south) -- (sign1.north);
\draw[c3flow] (sign1.south) -- (eg1.north);
\node[c3card,anchor=north,text width=8.8cm,align=center] (pay2) at (14.25,-2.7) {\textbf{\Money{$Y_2$}: mobile-data backup}\par $\Money{f_2(0)}=\Money{60},\quad\Money{f_2(1)}=\Money{-10}\ \Money{\mathrm{USD}}$\par $\ExpectedValue{g_2}(\SuccessProb{p})=\Money{60}-\Money{70}\SuccessProb{p}$};

\draw[c3flow] (formula.south) -- ++(0,-.25) -| (pay2.north);
\node[c3card,anchor=north,text width=8.8cm,align=center] (initial2) at (14.25,-5.0) {\ExpectedValue{Initial expectation}\par $\ExpectedValue{g_2}\!\left(\SuccessProb{\frac{1}{10}}\right)=\ExpectedValue{53\ \mathrm{USD}}$\par \ExpectedValue{expected net gain} per stated period};

\node[c3card,anchor=north,text width=8.8cm,align=center] (root2) at (14.25,-7.15) {\SuccessProb{Exact break-even equation}\par $\Money{60}-\Money{70}\SuccessProb{p}=\ExpectedValue{0}$\par $\SuccessProb{p_2^*}=\SuccessProb{\frac{6}{7}}$};

\node[c3card,anchor=north,text width=8.8cm,align=center] (sign2) at (14.25,-9.45) {$\ExpectedValue{g_2}>\ExpectedValue{0}:\ \SuccessProb{p}\in \left[\SuccessProb{0},\SuccessProb{\frac{6}{7}}\right)$\par $\ExpectedValue{g_2}=\ExpectedValue{0}:\ \SuccessProb{p}=\SuccessProb{\frac{6}{7}}$\par $\ExpectedValue{g_2}<\ExpectedValue{0}:\ \SuccessProb{p}\in \left(\SuccessProb{\frac{6}{7}},\SuccessProb{1}\right]$};

\node[inner sep=0pt,anchor=north,text width=8.8cm,align=center] (eg2) at (14.25,-12.0) {\CThreeExpectationGraph{2}{60}{-10}{6/7}{1/10}};

\draw[c3flow] (pay2.south) -- (initial2.north);
\draw[c3flow] (initial2.south) -- (root2.north);
\draw[c3flow] (root2.south) -- (sign2.north);
\draw[c3flow] (sign2.south) -- (eg2.north);
\end{tikzpicture}\end{center}
The dashed horizontal line is \ExpectedValue{zero}; filled circles show endpoint payoffs and the exact root. The square marks the \ExpectedValue{initial expectation}. The $+$ and $-$ labels connect the graph to its strict sign intervals.
\par \ExpectedValue{Positive expectation} means \ExpectedValue{expected net gain}; \ExpectedValue{negative expectation} means \ExpectedValue{expected net loss}; \ExpectedValue{zero expectation} means break-even. These are long-run averages per stated period, not guarantees of realized profit. Each branch uses its own payoffs and root.
\endgroup
}
\newcommand{\CThreeStatementInventory}{%
Observe a student's campus sticker sales over one weekend. \HLMathEvent{High demand} means campus-market attendance meets the organizer's fixed target; \HLMathEvent{low demand} means it falls below that target. Compare a basic sticker pack with a premium pack printed more cheaply per sticker through a shared bulk order. These are alternative stall plans after printing and stall costs. The two possible states are \HLMathEvent{high demand} and \HLMathEvent{low demand}. The \SuccessProb{initial probability} of \HLMathEvent{high demand} is $\SuccessProb{p}=\SuccessProb{\frac{2}{5}}$. All entries below are \Money{net profits} in \Money{USD} for the entire stated period, including relevant costs.
\begin{center}\begin{tabular}{lrr}\toprule Alternative & \HLMathEvent{low demand} & \HLMathEvent{high demand}\\\midrule
$A_1$: basic sticker pack & \Money{-40} & \Money{80}\\
$A_2$: premium sticker pack & \Money{-20} & \Money{180}\\
\bottomrule\end{tabular}\end{center}

}
\newcommand{\CThreeModelInventory}{%
\clearpage\noindent\textbf{Task 1 summary: shared states and induced distributions.}\par\medskip

\begingroup\small\setlength{\parskip}{2pt}\tikzset{gten marker/.append style={inner xsep=.4mm,inner ysep=.2mm}}
\begin{center}\begin{tikzpicture}[x=1cm,y=1cm]
\node[c3card,anchor=north,text width=7cm,align=center] (ev1) at (4,0) {$\HLMathEvent{E}$: \HLMathEvent{high demand}, $\SuccessProb{P}(\HLMathEvent{E})=\SuccessProb{p}$};

\node[c3card,anchor=north,text width=7cm,align=center] (ev0) at (15,0) {$\HLMathEvent{E^c}$: \HLMathEvent{low demand}, $\FailureProb{P}(\HLMathEvent{E^c})=\FailureProb{1-p}$};

\node[c3card,anchor=north,text width=6cm,align=center] (ind1) at (4,-1.2) {$\Indicator{X(\omega)}=\Indicator{1}$, for $\HLMathEvent{\omega\in E}$};

\node[c3card,anchor=north,text width=6cm,align=center] (ind0) at (15,-1.2) {$\Indicator{X(\omega)}=\Indicator{0}$, for $\HLMathEvent{\omega\in E^c}$};

\draw[c3flow] (ev1) -- (ind1);\draw[c3flow] (ev0) -- (ind0);
\node[c3card,anchor=north,text width=16.2cm,align=center] (map1) at (9.5,-2.65) {\textbf{\Money{$Y_1$}: basic sticker pack}\par 
$\Money{Y_1(\omega)}=\TransformationFunction{f_1}(\Indicator{X(\omega)})=\Money{f_1(1)}=\Money{80}\ \Money{\mathrm{USD}}$, for $\HLMathEvent{\omega\in E}$\par 
$\Money{Y_1(\omega)}=\TransformationFunction{f_1}(\Indicator{X(\omega)})=\Money{f_1(0)}=\Money{-40}\ \Money{\mathrm{USD}}$, for $\HLMathEvent{\omega\in E^c}$
};

\draw[c3flow] (ind1.west) -- (.15,-1.65) |- ([yshift=-7mm]map1.north west);
\draw[c3flow,dashed] (ind0.east) -- (18.85,-1.65) |- ([yshift=-11.5mm]map1.north east);
\node[c3card,anchor=north,text width=16.2cm,align=center] (map2) at (9.5,-4.45) {\textbf{\Money{$Y_2$}: premium sticker pack}\par 
$\Money{Y_2(\omega)}=\TransformationFunction{f_2}(\Indicator{X(\omega)})=\Money{f_2(1)}=\Money{180}\ \Money{\mathrm{USD}}$, for $\HLMathEvent{\omega\in E}$\par 
$\Money{Y_2(\omega)}=\TransformationFunction{f_2}(\Indicator{X(\omega)})=\Money{f_2(0)}=\Money{-20}\ \Money{\mathrm{USD}}$, for $\HLMathEvent{\omega\in E^c}$
};

\draw[c3flow] (ind1.west) -- (.15,-1.65) |- ([yshift=-7mm]map2.north west);
\draw[c3flow,dashed] (ind0.east) -- (18.85,-1.65) |- ([yshift=-11.5mm]map2.north east);
\node[inner sep=1pt,anchor=north,text width=5.913333333333333cm,align=center] (head0) at (3.0666666666666664,-6.4) {$\Indicator{X}$};

\node[c3card,anchor=north,text width=5.913333333333333cm,align=center] (cdf0) at (3.0666666666666664,-7.050000000000001) {$\HLCumulativeProbability{F_{X}(t)}$\par $\begin{cases}\HLCumulativeProbability{0},&t<\Indicator{0},\\\FailureProb{1-p},&\Indicator{0}\leq t<\Indicator{1},\\\HLCumulativeProbability{1},&t\geq \Indicator{1}.\end{cases}$};

\node[c3card,anchor=north,text width=5.913333333333333cm,align=center] (pmf0) at (3.0666666666666664,-9.8) {$\HLProbability{m_{X}}$\par $\begin{array}{ll}\Indicator{0}:&\FailureProb{1-p}\\\Indicator{1}:&\SuccessProb{p}\end{array}$\par \HLProbability{zero elsewhere}};

\node[inner sep=0pt,anchor=north,text width=5.913333333333333cm,align=center] (graph0) at (3.0666666666666664,-12.100000000000001) {\CThreeDistributionPair{X}{0}{1}{2/5}{unitless}{\Indicator}};

\node[inner sep=1pt,anchor=north,text width=5.913333333333333cm,align=center] (head1) at (9.4,-6.4) {$\Money{Y_1}$ $\leftarrow\TransformationFunction{f_1}$};

\node[c3card,anchor=north,text width=5.913333333333333cm,align=center] (cdf1) at (9.4,-7.050000000000001) {$\HLCumulativeProbability{F_{Y_1}(t)}$\par $\begin{cases}\HLCumulativeProbability{0},&t<\Money{-40},\\\FailureProb{1-p},&\Money{-40}\leq t<\Money{80},\\\HLCumulativeProbability{1},&t\geq \Money{80}.\end{cases}$};

\node[c3card,anchor=north,text width=5.913333333333333cm,align=center] (pmf1) at (9.4,-9.8) {$\HLProbability{m_{Y_1}}$\par $\begin{array}{ll}\Money{-40}:&\FailureProb{1-p}\\\Money{80}:&\SuccessProb{p}\end{array}$\par \HLProbability{zero elsewhere}};

\node[inner sep=0pt,anchor=north,text width=5.913333333333333cm,align=center] (graph1) at (9.4,-12.100000000000001) {\CThreeDistributionPair{Y_1}{-40}{80}{2/5}{\Money{USD}}{\Money}};

\draw[c3flow] (cdf0.north) -- ++(0,.27) -| (cdf1.north);
\draw[c3flow] (pmf0.south) -- ++(0,-.30) -| (pmf1.south);
\draw[c3flow] (map1.east) -- (18.28,-3.3) -- (18.28,-6.1000000000000005) -| (head1.north);
\node[inner sep=1pt,anchor=north,text width=5.913333333333333cm,align=center] (head2) at (15.733333333333333,-6.4) {$\Money{Y_2}$ $\leftarrow\TransformationFunction{f_2}$};

\node[c3card,anchor=north,text width=5.913333333333333cm,align=center] (cdf2) at (15.733333333333333,-7.050000000000001) {$\HLCumulativeProbability{F_{Y_2}(t)}$\par $\begin{cases}\HLCumulativeProbability{0},&t<\Money{-20},\\\FailureProb{1-p},&\Money{-20}\leq t<\Money{180},\\\HLCumulativeProbability{1},&t\geq \Money{180}.\end{cases}$};

\node[c3card,anchor=north,text width=5.913333333333333cm,align=center] (pmf2) at (15.733333333333333,-9.8) {$\HLProbability{m_{Y_2}}$\par $\begin{array}{ll}\Money{-20}:&\FailureProb{1-p}\\\Money{180}:&\SuccessProb{p}\end{array}$\par \HLProbability{zero elsewhere}};

\node[inner sep=0pt,anchor=north,text width=5.913333333333333cm,align=center] (graph2) at (15.733333333333333,-12.100000000000001) {\CThreeDistributionPair{Y_2}{-20}{180}{2/5}{\Money{USD}}{\Money}};

\draw[c3flow] (cdf0.north) -- ++(0,.27) -| (cdf2.north);
\draw[c3flow] (pmf0.south) -- ++(0,-.30) -| (pmf2.south);
\draw[c3flow] (map2.east) -- (18.1,-5.1) -- (18.1,-5.98) -| (head2.north);
\end{tikzpicture}\end{center}
\textbf{Numerical graph illustration:} $\SuccessProb{p_0}=\SuccessProb{\frac{2}{5}}$. All displayed distribution formulas remain symbolic in $\SuccessProb{p}$.
\par\textbf{CDF above, PMF below.} Filled stems translate to pale dashed jump shadows; thin arrows link equal physical lengths. Open and filled circles show excluded and included endpoints. Distribution arrows retain state probabilities while mapping values, rather than applying $\TransformationFunction{f_i}$ to CDF heights. See \hyperref[ex:intro-one]{Example 1}.\par
For $0<\SuccessProb{p}<1$, $\Indicator{S_X}=\{\Indicator{0},\Indicator{1}\}$; at $\SuccessProb{p}=0$ only $\Indicator{0}$ remains, and at $\SuccessProb{p}=1$ only $\Indicator{1}$ remains. Each monetary support retains $\Money{f_i(0)}$ at $\SuccessProb{p}=0$ and $\Money{f_i(1)}$ at $\SuccessProb{p}=1$.
\par\textbf{Loss events.} $\HLMathEvent{\{Y_1<0\}}=\HLMathEvent{\{X=0\}}=\HLMathEvent{E^c}$, probability $\FailureProb{1-p}$; $\HLMathEvent{\{Y_2<0\}}=\HLMathEvent{\{X=0\}}=\HLMathEvent{E^c}$, probability $\FailureProb{1-p}$.
\endgroup
}
\newcommand{\CThreeExpectInventory}{%
\clearpage\noindent\textbf{Task 2 summary: parallel expectation analyses.}\par\medskip

\begingroup\small\tikzset{gten marker/.append style={inner xsep=.4mm,inner ysep=.2mm}}
\begin{center}\begin{tikzpicture}[x=1cm,y=1cm]
\node[c3card,anchor=north,text width=17cm,align=center] (inputs) at (9.5,0) {Shared probability inputs: $\SuccessProb{p}$ and $\FailureProb{1-p}$; hold each alternative\textquotesingle s \Money{payoffs} fixed.};

\node[c3card,anchor=north,text width=17cm,align=center] (formula) at (9.5,-1.2) {$\ExpectedValue{g_i}(\SuccessProb{p})=\ExpectedValue{\mathbb E}[\Money{Y_i}]=\SuccessProb{p}\Money{f_i(1)}+(\FailureProb{1-p})\Money{f_i(0)}$};

\draw[c3flow] (inputs) -- (formula);
\node[c3card,anchor=north,text width=8.8cm,align=center] (pay1) at (4.75,-2.7) {\textbf{\Money{$Y_1$}: basic sticker pack}\par $\Money{f_1(0)}=\Money{-40},\quad\Money{f_1(1)}=\Money{80}\ \Money{\mathrm{USD}}$\par $\ExpectedValue{g_1}(\SuccessProb{p})=\Money{-40}+\Money{120}\SuccessProb{p}$};

\draw[c3flow] (formula.south) -- ++(0,-.25) -| (pay1.north);
\node[c3card,anchor=north,text width=8.8cm,align=center] (initial1) at (4.75,-5.0) {\ExpectedValue{Initial expectation}\par $\ExpectedValue{g_1}\!\left(\SuccessProb{\frac{2}{5}}\right)=\ExpectedValue{8\ \mathrm{USD}}$\par \ExpectedValue{expected net gain} per stated period};

\node[c3card,anchor=north,text width=8.8cm,align=center] (root1) at (4.75,-7.15) {\SuccessProb{Exact break-even equation}\par $\Money{-40}+\Money{120}\SuccessProb{p}=\ExpectedValue{0}$\par $\SuccessProb{p_1^*}=\SuccessProb{\frac{1}{3}}$};

\node[c3card,anchor=north,text width=8.8cm,align=center] (sign1) at (4.75,-9.45) {$\ExpectedValue{g_1}>\ExpectedValue{0}:\ \SuccessProb{p}\in \left(\SuccessProb{\frac{1}{3}},\SuccessProb{1}\right]$\par $\ExpectedValue{g_1}=\ExpectedValue{0}:\ \SuccessProb{p}=\SuccessProb{\frac{1}{3}}$\par $\ExpectedValue{g_1}<\ExpectedValue{0}:\ \SuccessProb{p}\in \left[\SuccessProb{0},\SuccessProb{\frac{1}{3}}\right)$};

\node[inner sep=0pt,anchor=north,text width=8.8cm,align=center] (eg1) at (4.75,-12.0) {\CThreeExpectationGraph{1}{-40}{80}{1/3}{2/5}};

\draw[c3flow] (pay1.south) -- (initial1.north);
\draw[c3flow] (initial1.south) -- (root1.north);
\draw[c3flow] (root1.south) -- (sign1.north);
\draw[c3flow] (sign1.south) -- (eg1.north);
\node[c3card,anchor=north,text width=8.8cm,align=center] (pay2) at (14.25,-2.7) {\textbf{\Money{$Y_2$}: premium sticker pack}\par $\Money{f_2(0)}=\Money{-20},\quad\Money{f_2(1)}=\Money{180}\ \Money{\mathrm{USD}}$\par $\ExpectedValue{g_2}(\SuccessProb{p})=\Money{-20}+\Money{200}\SuccessProb{p}$};

\draw[c3flow] (formula.south) -- ++(0,-.25) -| (pay2.north);
\node[c3card,anchor=north,text width=8.8cm,align=center] (initial2) at (14.25,-5.0) {\ExpectedValue{Initial expectation}\par $\ExpectedValue{g_2}\!\left(\SuccessProb{\frac{2}{5}}\right)=\ExpectedValue{60\ \mathrm{USD}}$\par \ExpectedValue{expected net gain} per stated period};

\node[c3card,anchor=north,text width=8.8cm,align=center] (root2) at (14.25,-7.15) {\SuccessProb{Exact break-even equation}\par $\Money{-20}+\Money{200}\SuccessProb{p}=\ExpectedValue{0}$\par $\SuccessProb{p_2^*}=\SuccessProb{\frac{1}{10}}$};

\node[c3card,anchor=north,text width=8.8cm,align=center] (sign2) at (14.25,-9.45) {$\ExpectedValue{g_2}>\ExpectedValue{0}:\ \SuccessProb{p}\in \left(\SuccessProb{\frac{1}{10}},\SuccessProb{1}\right]$\par $\ExpectedValue{g_2}=\ExpectedValue{0}:\ \SuccessProb{p}=\SuccessProb{\frac{1}{10}}$\par $\ExpectedValue{g_2}<\ExpectedValue{0}:\ \SuccessProb{p}\in \left[\SuccessProb{0},\SuccessProb{\frac{1}{10}}\right)$};

\node[inner sep=0pt,anchor=north,text width=8.8cm,align=center] (eg2) at (14.25,-12.0) {\CThreeExpectationGraph{2}{-20}{180}{1/10}{2/5}};

\draw[c3flow] (pay2.south) -- (initial2.north);
\draw[c3flow] (initial2.south) -- (root2.north);
\draw[c3flow] (root2.south) -- (sign2.north);
\draw[c3flow] (sign2.south) -- (eg2.north);
\end{tikzpicture}\end{center}
The dashed horizontal line is \ExpectedValue{zero}; filled circles show endpoint payoffs and the exact root. The square marks the \ExpectedValue{initial expectation}. The $+$ and $-$ labels connect the graph to its strict sign intervals.
\par \ExpectedValue{Positive expectation} means \ExpectedValue{expected net gain}; \ExpectedValue{negative expectation} means \ExpectedValue{expected net loss}; \ExpectedValue{zero expectation} means break-even. These are long-run averages per stated period, not guarantees of realized profit. Each branch uses its own payoffs and root.
\endgroup
}
\newcommand{\CThreeStatementCoffee}{%
Observe a graduate's paid sessions over one day. A campus careers fair either goes ahead or is cancelled. Compare paid portraits at the fair with an online CV workshop. The fair brings portrait customers but draws students away from the workshop; cancellation reverses this effect. \Money{Payoffs} include venue, platform, travel and refund costs. The two possible states are \HLMathEvent{the fair goes ahead} and \HLMathEvent{the fair is cancelled}. The \SuccessProb{initial probability} that \HLMathEvent{the fair goes ahead} is $\SuccessProb{p}=\SuccessProb{\frac{1}{2}}$. All entries below are \Money{net profits} in \Money{USD} for the entire stated period, including relevant costs.
\begin{center}\begin{tabular}{lrr}\toprule Alternative & \HLMathEvent{the fair is cancelled} & \HLMathEvent{the fair goes ahead}\\\midrule
$A_1$: portrait sessions & \Money{-40} & \Money{100}\\
$A_2$: online CV workshop & \Money{100} & \Money{-20}\\
\bottomrule\end{tabular}\end{center}

}
\newcommand{\CThreeModelCoffee}{%
\clearpage\noindent\textbf{Task 1 summary: shared states and induced distributions.}\par\medskip

\begingroup\small\setlength{\parskip}{2pt}\tikzset{gten marker/.append style={inner xsep=.4mm,inner ysep=.2mm}}
\begin{center}\begin{tikzpicture}[x=1cm,y=1cm]
\node[c3card,anchor=north,text width=7cm,align=center] (ev1) at (4,0) {$\HLMathEvent{E}$: \HLMathEvent{the fair goes ahead}, $\SuccessProb{P}(\HLMathEvent{E})=\SuccessProb{p}$};

\node[c3card,anchor=north,text width=7cm,align=center] (ev0) at (15,0) {$\HLMathEvent{E^c}$: \HLMathEvent{the fair is cancelled}, $\FailureProb{P}(\HLMathEvent{E^c})=\FailureProb{1-p}$};

\node[c3card,anchor=north,text width=6cm,align=center] (ind1) at (4,-1.2) {$\Indicator{X(\omega)}=\Indicator{1}$, for $\HLMathEvent{\omega\in E}$};

\node[c3card,anchor=north,text width=6cm,align=center] (ind0) at (15,-1.2) {$\Indicator{X(\omega)}=\Indicator{0}$, for $\HLMathEvent{\omega\in E^c}$};

\draw[c3flow] (ev1) -- (ind1);\draw[c3flow] (ev0) -- (ind0);
\node[c3card,anchor=north,text width=16.2cm,align=center] (map1) at (9.5,-2.65) {\textbf{\Money{$Y_1$}: portrait sessions}\par 
$\Money{Y_1(\omega)}=\TransformationFunction{f_1}(\Indicator{X(\omega)})=\Money{f_1(1)}=\Money{100}\ \Money{\mathrm{USD}}$, for $\HLMathEvent{\omega\in E}$\par 
$\Money{Y_1(\omega)}=\TransformationFunction{f_1}(\Indicator{X(\omega)})=\Money{f_1(0)}=\Money{-40}\ \Money{\mathrm{USD}}$, for $\HLMathEvent{\omega\in E^c}$
};

\draw[c3flow] (ind1.west) -- (.15,-1.65) |- ([yshift=-7mm]map1.north west);
\draw[c3flow,dashed] (ind0.east) -- (18.85,-1.65) |- ([yshift=-11.5mm]map1.north east);
\node[c3card,anchor=north,text width=16.2cm,align=center] (map2) at (9.5,-4.45) {\textbf{\Money{$Y_2$}: online CV workshop}\par 
$\Money{Y_2(\omega)}=\TransformationFunction{f_2}(\Indicator{X(\omega)})=\Money{f_2(1)}=\Money{-20}\ \Money{\mathrm{USD}}$, for $\HLMathEvent{\omega\in E}$\par 
$\Money{Y_2(\omega)}=\TransformationFunction{f_2}(\Indicator{X(\omega)})=\Money{f_2(0)}=\Money{100}\ \Money{\mathrm{USD}}$, for $\HLMathEvent{\omega\in E^c}$
};

\draw[c3flow] (ind1.west) -- (.15,-1.65) |- ([yshift=-7mm]map2.north west);
\draw[c3flow,dashed] (ind0.east) -- (18.85,-1.65) |- ([yshift=-11.5mm]map2.north east);
\node[inner sep=1pt,anchor=north,text width=5.913333333333333cm,align=center] (head0) at (3.0666666666666664,-6.4) {$\Indicator{X}$};

\node[c3card,anchor=north,text width=5.913333333333333cm,align=center] (cdf0) at (3.0666666666666664,-7.050000000000001) {$\HLCumulativeProbability{F_{X}(t)}$\par $\begin{cases}\HLCumulativeProbability{0},&t<\Indicator{0},\\\FailureProb{1-p},&\Indicator{0}\leq t<\Indicator{1},\\\HLCumulativeProbability{1},&t\geq \Indicator{1}.\end{cases}$};

\node[c3card,anchor=north,text width=5.913333333333333cm,align=center] (pmf0) at (3.0666666666666664,-9.8) {$\HLProbability{m_{X}}$\par $\begin{array}{ll}\Indicator{0}:&\FailureProb{1-p}\\\Indicator{1}:&\SuccessProb{p}\end{array}$\par \HLProbability{zero elsewhere}};

\node[inner sep=0pt,anchor=north,text width=5.913333333333333cm,align=center] (graph0) at (3.0666666666666664,-12.100000000000001) {\CThreeDistributionPair{X}{0}{1}{1/2}{unitless}{\Indicator}};

\node[inner sep=1pt,anchor=north,text width=5.913333333333333cm,align=center] (head1) at (9.4,-6.4) {$\Money{Y_1}$ $\leftarrow\TransformationFunction{f_1}$};

\node[c3card,anchor=north,text width=5.913333333333333cm,align=center] (cdf1) at (9.4,-7.050000000000001) {$\HLCumulativeProbability{F_{Y_1}(t)}$\par $\begin{cases}\HLCumulativeProbability{0},&t<\Money{-40},\\\FailureProb{1-p},&\Money{-40}\leq t<\Money{100},\\\HLCumulativeProbability{1},&t\geq \Money{100}.\end{cases}$};

\node[c3card,anchor=north,text width=5.913333333333333cm,align=center] (pmf1) at (9.4,-9.8) {$\HLProbability{m_{Y_1}}$\par $\begin{array}{ll}\Money{-40}:&\FailureProb{1-p}\\\Money{100}:&\SuccessProb{p}\end{array}$\par \HLProbability{zero elsewhere}};

\node[inner sep=0pt,anchor=north,text width=5.913333333333333cm,align=center] (graph1) at (9.4,-12.100000000000001) {\CThreeDistributionPair{Y_1}{-40}{100}{1/2}{\Money{USD}}{\Money}};

\draw[c3flow] (cdf0.north) -- ++(0,.27) -| (cdf1.north);
\draw[c3flow] (pmf0.south) -- ++(0,-.30) -| (pmf1.south);
\draw[c3flow] (map1.east) -- (18.28,-3.3) -- (18.28,-6.1000000000000005) -| (head1.north);
\node[inner sep=1pt,anchor=north,text width=5.913333333333333cm,align=center] (head2) at (15.733333333333333,-6.4) {$\Money{Y_2}$ $\leftarrow\TransformationFunction{f_2}$};

\node[c3card,anchor=north,text width=5.913333333333333cm,align=center] (cdf2) at (15.733333333333333,-7.050000000000001) {$\HLCumulativeProbability{F_{Y_2}(t)}$\par $\begin{cases}\HLCumulativeProbability{0},&t<\Money{-20},\\\SuccessProb{p},&\Money{-20}\leq t<\Money{100},\\\HLCumulativeProbability{1},&t\geq \Money{100}.\end{cases}$};

\node[c3card,anchor=north,text width=5.913333333333333cm,align=center] (pmf2) at (15.733333333333333,-9.8) {$\HLProbability{m_{Y_2}}$\par $\begin{array}{ll}\Money{100}:&\FailureProb{1-p}\\\Money{-20}:&\SuccessProb{p}\end{array}$\par \HLProbability{zero elsewhere}};

\node[inner sep=0pt,anchor=north,text width=5.913333333333333cm,align=center] (graph2) at (15.733333333333333,-12.100000000000001) {\CThreeDistributionPair{Y_2}{100}{-20}{1/2}{\Money{USD}}{\Money}};

\draw[c3flow] (cdf0.north) -- ++(0,.27) -| (cdf2.north);
\draw[c3flow] (pmf0.south) -- ++(0,-.30) -| (pmf2.south);
\draw[c3flow] (map2.east) -- (18.1,-5.1) -- (18.1,-5.98) -| (head2.north);
\end{tikzpicture}\end{center}
\textbf{Numerical graph illustration:} $\SuccessProb{p_0}=\SuccessProb{\frac{1}{2}}$. All displayed distribution formulas remain symbolic in $\SuccessProb{p}$.
\par\textbf{CDF above, PMF below.} Filled stems translate to pale dashed jump shadows; thin arrows link equal physical lengths. Open and filled circles show excluded and included endpoints. Distribution arrows retain state probabilities while mapping values, rather than applying $\TransformationFunction{f_i}$ to CDF heights. See \hyperref[ex:intro-one]{Example 1}.\par
For $0<\SuccessProb{p}<1$, $\Indicator{S_X}=\{\Indicator{0},\Indicator{1}\}$; at $\SuccessProb{p}=0$ only $\Indicator{0}$ remains, and at $\SuccessProb{p}=1$ only $\Indicator{1}$ remains. Each monetary support retains $\Money{f_i(0)}$ at $\SuccessProb{p}=0$ and $\Money{f_i(1)}$ at $\SuccessProb{p}=1$.
\par\textbf{Loss events.} $\HLMathEvent{\{Y_1<0\}}=\HLMathEvent{\{X=0\}}=\HLMathEvent{E^c}$, probability $\FailureProb{1-p}$; $\HLMathEvent{\{Y_2<0\}}=\HLMathEvent{\{X=1\}}=\HLMathEvent{E}$, probability $\SuccessProb{p}$.
\endgroup
}
\newcommand{\CThreeExpectCoffee}{%
\clearpage\noindent\textbf{Task 2 summary: parallel expectation analyses.}\par\medskip

\begingroup\small\tikzset{gten marker/.append style={inner xsep=.4mm,inner ysep=.2mm}}
\begin{center}\begin{tikzpicture}[x=1cm,y=1cm]
\node[c3card,anchor=north,text width=17cm,align=center] (inputs) at (9.5,0) {Shared probability inputs: $\SuccessProb{p}$ and $\FailureProb{1-p}$; hold each alternative\textquotesingle s \Money{payoffs} fixed.};

\node[c3card,anchor=north,text width=17cm,align=center] (formula) at (9.5,-1.2) {$\ExpectedValue{g_i}(\SuccessProb{p})=\ExpectedValue{\mathbb E}[\Money{Y_i}]=\SuccessProb{p}\Money{f_i(1)}+(\FailureProb{1-p})\Money{f_i(0)}$};

\draw[c3flow] (inputs) -- (formula);
\node[c3card,anchor=north,text width=8.8cm,align=center] (pay1) at (4.75,-2.7) {\textbf{\Money{$Y_1$}: portrait sessions}\par $\Money{f_1(0)}=\Money{-40},\quad\Money{f_1(1)}=\Money{100}\ \Money{\mathrm{USD}}$\par $\ExpectedValue{g_1}(\SuccessProb{p})=\Money{-40}+\Money{140}\SuccessProb{p}$};

\draw[c3flow] (formula.south) -- ++(0,-.25) -| (pay1.north);
\node[c3card,anchor=north,text width=8.8cm,align=center] (initial1) at (4.75,-5.0) {\ExpectedValue{Initial expectation}\par $\ExpectedValue{g_1}\!\left(\SuccessProb{\frac{1}{2}}\right)=\ExpectedValue{30\ \mathrm{USD}}$\par \ExpectedValue{expected net gain} per stated period};

\node[c3card,anchor=north,text width=8.8cm,align=center] (root1) at (4.75,-7.15) {\SuccessProb{Exact break-even equation}\par $\Money{-40}+\Money{140}\SuccessProb{p}=\ExpectedValue{0}$\par $\SuccessProb{p_1^*}=\SuccessProb{\frac{2}{7}}$};

\node[c3card,anchor=north,text width=8.8cm,align=center] (sign1) at (4.75,-9.45) {$\ExpectedValue{g_1}>\ExpectedValue{0}:\ \SuccessProb{p}\in \left(\SuccessProb{\frac{2}{7}},\SuccessProb{1}\right]$\par $\ExpectedValue{g_1}=\ExpectedValue{0}:\ \SuccessProb{p}=\SuccessProb{\frac{2}{7}}$\par $\ExpectedValue{g_1}<\ExpectedValue{0}:\ \SuccessProb{p}\in \left[\SuccessProb{0},\SuccessProb{\frac{2}{7}}\right)$};

\node[inner sep=0pt,anchor=north,text width=8.8cm,align=center] (eg1) at (4.75,-12.0) {\CThreeExpectationGraph{1}{-40}{100}{2/7}{1/2}};

\draw[c3flow] (pay1.south) -- (initial1.north);
\draw[c3flow] (initial1.south) -- (root1.north);
\draw[c3flow] (root1.south) -- (sign1.north);
\draw[c3flow] (sign1.south) -- (eg1.north);
\node[c3card,anchor=north,text width=8.8cm,align=center] (pay2) at (14.25,-2.7) {\textbf{\Money{$Y_2$}: online CV workshop}\par $\Money{f_2(0)}=\Money{100},\quad\Money{f_2(1)}=\Money{-20}\ \Money{\mathrm{USD}}$\par $\ExpectedValue{g_2}(\SuccessProb{p})=\Money{100}-\Money{120}\SuccessProb{p}$};

\draw[c3flow] (formula.south) -- ++(0,-.25) -| (pay2.north);
\node[c3card,anchor=north,text width=8.8cm,align=center] (initial2) at (14.25,-5.0) {\ExpectedValue{Initial expectation}\par $\ExpectedValue{g_2}\!\left(\SuccessProb{\frac{1}{2}}\right)=\ExpectedValue{40\ \mathrm{USD}}$\par \ExpectedValue{expected net gain} per stated period};

\node[c3card,anchor=north,text width=8.8cm,align=center] (root2) at (14.25,-7.15) {\SuccessProb{Exact break-even equation}\par $\Money{100}-\Money{120}\SuccessProb{p}=\ExpectedValue{0}$\par $\SuccessProb{p_2^*}=\SuccessProb{\frac{5}{6}}$};

\node[c3card,anchor=north,text width=8.8cm,align=center] (sign2) at (14.25,-9.45) {$\ExpectedValue{g_2}>\ExpectedValue{0}:\ \SuccessProb{p}\in \left[\SuccessProb{0},\SuccessProb{\frac{5}{6}}\right)$\par $\ExpectedValue{g_2}=\ExpectedValue{0}:\ \SuccessProb{p}=\SuccessProb{\frac{5}{6}}$\par $\ExpectedValue{g_2}<\ExpectedValue{0}:\ \SuccessProb{p}\in \left(\SuccessProb{\frac{5}{6}},\SuccessProb{1}\right]$};

\node[inner sep=0pt,anchor=north,text width=8.8cm,align=center] (eg2) at (14.25,-12.0) {\CThreeExpectationGraph{2}{100}{-20}{5/6}{1/2}};

\draw[c3flow] (pay2.south) -- (initial2.north);
\draw[c3flow] (initial2.south) -- (root2.north);
\draw[c3flow] (root2.south) -- (sign2.north);
\draw[c3flow] (sign2.south) -- (eg2.north);
\end{tikzpicture}\end{center}
The dashed horizontal line is \ExpectedValue{zero}; filled circles show endpoint payoffs and the exact root. The square marks the \ExpectedValue{initial expectation}. The $+$ and $-$ labels connect the graph to its strict sign intervals.
\par \ExpectedValue{Positive expectation} means \ExpectedValue{expected net gain}; \ExpectedValue{negative expectation} means \ExpectedValue{expected net loss}; \ExpectedValue{zero expectation} means break-even. These are long-run averages per stated period, not guarantees of realized profit. Each branch uses its own payoffs and root.
\endgroup
}
\newcommand{\CThreeStatementEnergy}{%
Observe a freelancer's digital projects over one month. \HLMathEvent{High inquiry volume} means messages from potential clients meet a fixed monthly target; \HLMathEvent{low inquiry volume} means fewer messages. Compare offering social-media graphics with offering short video edits to the same client pool. \Money{Payoffs} include software, platform and promotion costs. The two possible states are \HLMathEvent{high inquiry volume} and \HLMathEvent{low inquiry volume}. The \SuccessProb{initial probability} of \HLMathEvent{high inquiry volume} is $\SuccessProb{p}=\SuccessProb{\frac{1}{4}}$. All entries below are \Money{net profits} in \Money{USD} for the entire stated period, including relevant costs.
\begin{center}\begin{tabular}{lrr}\toprule Alternative & \HLMathEvent{low inquiry volume} & \HLMathEvent{high inquiry volume}\\\midrule
$A_1$: social-media graphics & \Money{-30} & \Money{90}\\
$A_2$: short video edits & \Money{-60} & \Money{60}\\
\bottomrule\end{tabular}\end{center}

}
\newcommand{\CThreeModelEnergy}{%
\clearpage\noindent\textbf{Task 1 summary: shared states and induced distributions.}\par\medskip

\begingroup\small\setlength{\parskip}{2pt}\tikzset{gten marker/.append style={inner xsep=.4mm,inner ysep=.2mm}}
\begin{center}\begin{tikzpicture}[x=1cm,y=1cm]
\node[c3card,anchor=north,text width=7cm,align=center] (ev1) at (4,0) {$\HLMathEvent{E}$: \HLMathEvent{high inquiry volume}, $\SuccessProb{P}(\HLMathEvent{E})=\SuccessProb{p}$};

\node[c3card,anchor=north,text width=7cm,align=center] (ev0) at (15,0) {$\HLMathEvent{E^c}$: \HLMathEvent{low inquiry volume}, $\FailureProb{P}(\HLMathEvent{E^c})=\FailureProb{1-p}$};

\node[c3card,anchor=north,text width=6cm,align=center] (ind1) at (4,-1.2) {$\Indicator{X(\omega)}=\Indicator{1}$, for $\HLMathEvent{\omega\in E}$};

\node[c3card,anchor=north,text width=6cm,align=center] (ind0) at (15,-1.2) {$\Indicator{X(\omega)}=\Indicator{0}$, for $\HLMathEvent{\omega\in E^c}$};

\draw[c3flow] (ev1) -- (ind1);\draw[c3flow] (ev0) -- (ind0);
\node[c3card,anchor=north,text width=16.2cm,align=center] (map1) at (9.5,-2.65) {\textbf{\Money{$Y_1$}: social-media graphics}\par 
$\Money{Y_1(\omega)}=\TransformationFunction{f_1}(\Indicator{X(\omega)})=\Money{f_1(1)}=\Money{90}\ \Money{\mathrm{USD}}$, for $\HLMathEvent{\omega\in E}$\par 
$\Money{Y_1(\omega)}=\TransformationFunction{f_1}(\Indicator{X(\omega)})=\Money{f_1(0)}=\Money{-30}\ \Money{\mathrm{USD}}$, for $\HLMathEvent{\omega\in E^c}$
};

\draw[c3flow] (ind1.west) -- (.15,-1.65) |- ([yshift=-7mm]map1.north west);
\draw[c3flow,dashed] (ind0.east) -- (18.85,-1.65) |- ([yshift=-11.5mm]map1.north east);
\node[c3card,anchor=north,text width=16.2cm,align=center] (map2) at (9.5,-4.45) {\textbf{\Money{$Y_2$}: short video edits}\par 
$\Money{Y_2(\omega)}=\TransformationFunction{f_2}(\Indicator{X(\omega)})=\Money{f_2(1)}=\Money{60}\ \Money{\mathrm{USD}}$, for $\HLMathEvent{\omega\in E}$\par 
$\Money{Y_2(\omega)}=\TransformationFunction{f_2}(\Indicator{X(\omega)})=\Money{f_2(0)}=\Money{-60}\ \Money{\mathrm{USD}}$, for $\HLMathEvent{\omega\in E^c}$
};

\draw[c3flow] (ind1.west) -- (.15,-1.65) |- ([yshift=-7mm]map2.north west);
\draw[c3flow,dashed] (ind0.east) -- (18.85,-1.65) |- ([yshift=-11.5mm]map2.north east);
\node[inner sep=1pt,anchor=north,text width=5.913333333333333cm,align=center] (head0) at (3.0666666666666664,-6.4) {$\Indicator{X}$};

\node[c3card,anchor=north,text width=5.913333333333333cm,align=center] (cdf0) at (3.0666666666666664,-7.050000000000001) {$\HLCumulativeProbability{F_{X}(t)}$\par $\begin{cases}\HLCumulativeProbability{0},&t<\Indicator{0},\\\FailureProb{1-p},&\Indicator{0}\leq t<\Indicator{1},\\\HLCumulativeProbability{1},&t\geq \Indicator{1}.\end{cases}$};

\node[c3card,anchor=north,text width=5.913333333333333cm,align=center] (pmf0) at (3.0666666666666664,-9.8) {$\HLProbability{m_{X}}$\par $\begin{array}{ll}\Indicator{0}:&\FailureProb{1-p}\\\Indicator{1}:&\SuccessProb{p}\end{array}$\par \HLProbability{zero elsewhere}};

\node[inner sep=0pt,anchor=north,text width=5.913333333333333cm,align=center] (graph0) at (3.0666666666666664,-12.100000000000001) {\CThreeDistributionPair{X}{0}{1}{1/4}{unitless}{\Indicator}};

\node[inner sep=1pt,anchor=north,text width=5.913333333333333cm,align=center] (head1) at (9.4,-6.4) {$\Money{Y_1}$ $\leftarrow\TransformationFunction{f_1}$};

\node[c3card,anchor=north,text width=5.913333333333333cm,align=center] (cdf1) at (9.4,-7.050000000000001) {$\HLCumulativeProbability{F_{Y_1}(t)}$\par $\begin{cases}\HLCumulativeProbability{0},&t<\Money{-30},\\\FailureProb{1-p},&\Money{-30}\leq t<\Money{90},\\\HLCumulativeProbability{1},&t\geq \Money{90}.\end{cases}$};

\node[c3card,anchor=north,text width=5.913333333333333cm,align=center] (pmf1) at (9.4,-9.8) {$\HLProbability{m_{Y_1}}$\par $\begin{array}{ll}\Money{-30}:&\FailureProb{1-p}\\\Money{90}:&\SuccessProb{p}\end{array}$\par \HLProbability{zero elsewhere}};

\node[inner sep=0pt,anchor=north,text width=5.913333333333333cm,align=center] (graph1) at (9.4,-12.100000000000001) {\CThreeDistributionPair{Y_1}{-30}{90}{1/4}{\Money{USD}}{\Money}};

\draw[c3flow] (cdf0.north) -- ++(0,.27) -| (cdf1.north);
\draw[c3flow] (pmf0.south) -- ++(0,-.30) -| (pmf1.south);
\draw[c3flow] (map1.east) -- (18.28,-3.3) -- (18.28,-6.1000000000000005) -| (head1.north);
\node[inner sep=1pt,anchor=north,text width=5.913333333333333cm,align=center] (head2) at (15.733333333333333,-6.4) {$\Money{Y_2}$ $\leftarrow\TransformationFunction{f_2}$};

\node[c3card,anchor=north,text width=5.913333333333333cm,align=center] (cdf2) at (15.733333333333333,-7.050000000000001) {$\HLCumulativeProbability{F_{Y_2}(t)}$\par $\begin{cases}\HLCumulativeProbability{0},&t<\Money{-60},\\\FailureProb{1-p},&\Money{-60}\leq t<\Money{60},\\\HLCumulativeProbability{1},&t\geq \Money{60}.\end{cases}$};

\node[c3card,anchor=north,text width=5.913333333333333cm,align=center] (pmf2) at (15.733333333333333,-9.8) {$\HLProbability{m_{Y_2}}$\par $\begin{array}{ll}\Money{-60}:&\FailureProb{1-p}\\\Money{60}:&\SuccessProb{p}\end{array}$\par \HLProbability{zero elsewhere}};

\node[inner sep=0pt,anchor=north,text width=5.913333333333333cm,align=center] (graph2) at (15.733333333333333,-12.100000000000001) {\CThreeDistributionPair{Y_2}{-60}{60}{1/4}{\Money{USD}}{\Money}};

\draw[c3flow] (cdf0.north) -- ++(0,.27) -| (cdf2.north);
\draw[c3flow] (pmf0.south) -- ++(0,-.30) -| (pmf2.south);
\draw[c3flow] (map2.east) -- (18.1,-5.1) -- (18.1,-5.98) -| (head2.north);
\end{tikzpicture}\end{center}
\textbf{Numerical graph illustration:} $\SuccessProb{p_0}=\SuccessProb{\frac{1}{4}}$. All displayed distribution formulas remain symbolic in $\SuccessProb{p}$.
\par\textbf{CDF above, PMF below.} Filled stems translate to pale dashed jump shadows; thin arrows link equal physical lengths. Open and filled circles show excluded and included endpoints. Distribution arrows retain state probabilities while mapping values, rather than applying $\TransformationFunction{f_i}$ to CDF heights. See \hyperref[ex:intro-one]{Example 1}.\par
For $0<\SuccessProb{p}<1$, $\Indicator{S_X}=\{\Indicator{0},\Indicator{1}\}$; at $\SuccessProb{p}=0$ only $\Indicator{0}$ remains, and at $\SuccessProb{p}=1$ only $\Indicator{1}$ remains. Each monetary support retains $\Money{f_i(0)}$ at $\SuccessProb{p}=0$ and $\Money{f_i(1)}$ at $\SuccessProb{p}=1$.
\par\textbf{Loss events.} $\HLMathEvent{\{Y_1<0\}}=\HLMathEvent{\{X=0\}}=\HLMathEvent{E^c}$, probability $\FailureProb{1-p}$; $\HLMathEvent{\{Y_2<0\}}=\HLMathEvent{\{X=0\}}=\HLMathEvent{E^c}$, probability $\FailureProb{1-p}$.
\endgroup
}
\newcommand{\CThreeExpectEnergy}{%
\clearpage\noindent\textbf{Task 2 summary: parallel expectation analyses.}\par\medskip

\begingroup\small\tikzset{gten marker/.append style={inner xsep=.4mm,inner ysep=.2mm}}
\begin{center}\begin{tikzpicture}[x=1cm,y=1cm]
\node[c3card,anchor=north,text width=17cm,align=center] (inputs) at (9.5,0) {Shared probability inputs: $\SuccessProb{p}$ and $\FailureProb{1-p}$; hold each alternative\textquotesingle s \Money{payoffs} fixed.};

\node[c3card,anchor=north,text width=17cm,align=center] (formula) at (9.5,-1.2) {$\ExpectedValue{g_i}(\SuccessProb{p})=\ExpectedValue{\mathbb E}[\Money{Y_i}]=\SuccessProb{p}\Money{f_i(1)}+(\FailureProb{1-p})\Money{f_i(0)}$};

\draw[c3flow] (inputs) -- (formula);
\node[c3card,anchor=north,text width=8.8cm,align=center] (pay1) at (4.75,-2.7) {\textbf{\Money{$Y_1$}: social-media graphics}\par $\Money{f_1(0)}=\Money{-30},\quad\Money{f_1(1)}=\Money{90}\ \Money{\mathrm{USD}}$\par $\ExpectedValue{g_1}(\SuccessProb{p})=\Money{-30}+\Money{120}\SuccessProb{p}$};

\draw[c3flow] (formula.south) -- ++(0,-.25) -| (pay1.north);
\node[c3card,anchor=north,text width=8.8cm,align=center] (initial1) at (4.75,-5.0) {\ExpectedValue{Initial expectation}\par $\ExpectedValue{g_1}\!\left(\SuccessProb{\frac{1}{4}}\right)=\ExpectedValue{0\ \mathrm{USD}}$\par \ExpectedValue{break-even} per stated period};

\node[c3card,anchor=north,text width=8.8cm,align=center] (root1) at (4.75,-7.15) {\SuccessProb{Exact break-even equation}\par $\Money{-30}+\Money{120}\SuccessProb{p}=\ExpectedValue{0}$\par $\SuccessProb{p_1^*}=\SuccessProb{\frac{1}{4}}$};

\node[c3card,anchor=north,text width=8.8cm,align=center] (sign1) at (4.75,-9.45) {$\ExpectedValue{g_1}>\ExpectedValue{0}:\ \SuccessProb{p}\in \left(\SuccessProb{\frac{1}{4}},\SuccessProb{1}\right]$\par $\ExpectedValue{g_1}=\ExpectedValue{0}:\ \SuccessProb{p}=\SuccessProb{\frac{1}{4}}$\par $\ExpectedValue{g_1}<\ExpectedValue{0}:\ \SuccessProb{p}\in \left[\SuccessProb{0},\SuccessProb{\frac{1}{4}}\right)$};

\node[inner sep=0pt,anchor=north,text width=8.8cm,align=center] (eg1) at (4.75,-12.0) {\CThreeExpectationGraph{1}{-30}{90}{1/4}{1/4}};

\draw[c3flow] (pay1.south) -- (initial1.north);
\draw[c3flow] (initial1.south) -- (root1.north);
\draw[c3flow] (root1.south) -- (sign1.north);
\draw[c3flow] (sign1.south) -- (eg1.north);
\node[c3card,anchor=north,text width=8.8cm,align=center] (pay2) at (14.25,-2.7) {\textbf{\Money{$Y_2$}: short video edits}\par $\Money{f_2(0)}=\Money{-60},\quad\Money{f_2(1)}=\Money{60}\ \Money{\mathrm{USD}}$\par $\ExpectedValue{g_2}(\SuccessProb{p})=\Money{-60}+\Money{120}\SuccessProb{p}$};

\draw[c3flow] (formula.south) -- ++(0,-.25) -| (pay2.north);
\node[c3card,anchor=north,text width=8.8cm,align=center] (initial2) at (14.25,-5.0) {\ExpectedValue{Initial expectation}\par $\ExpectedValue{g_2}\!\left(\SuccessProb{\frac{1}{4}}\right)=\ExpectedValue{-30\ \mathrm{USD}}$\par \ExpectedValue{expected net loss} per stated period};

\node[c3card,anchor=north,text width=8.8cm,align=center] (root2) at (14.25,-7.15) {\SuccessProb{Exact break-even equation}\par $\Money{-60}+\Money{120}\SuccessProb{p}=\ExpectedValue{0}$\par $\SuccessProb{p_2^*}=\SuccessProb{\frac{1}{2}}$};

\node[c3card,anchor=north,text width=8.8cm,align=center] (sign2) at (14.25,-9.45) {$\ExpectedValue{g_2}>\ExpectedValue{0}:\ \SuccessProb{p}\in \left(\SuccessProb{\frac{1}{2}},\SuccessProb{1}\right]$\par $\ExpectedValue{g_2}=\ExpectedValue{0}:\ \SuccessProb{p}=\SuccessProb{\frac{1}{2}}$\par $\ExpectedValue{g_2}<\ExpectedValue{0}:\ \SuccessProb{p}\in \left[\SuccessProb{0},\SuccessProb{\frac{1}{2}}\right)$};

\node[inner sep=0pt,anchor=north,text width=8.8cm,align=center] (eg2) at (14.25,-12.0) {\CThreeExpectationGraph{2}{-60}{60}{1/2}{1/4}};

\draw[c3flow] (pay2.south) -- (initial2.north);
\draw[c3flow] (initial2.south) -- (root2.north);
\draw[c3flow] (root2.south) -- (sign2.north);
\draw[c3flow] (sign2.south) -- (eg2.north);
\end{tikzpicture}\end{center}
The dashed horizontal line is \ExpectedValue{zero}; filled circles show endpoint payoffs and the exact root. The square marks the \ExpectedValue{initial expectation}. The $+$ and $-$ labels connect the graph to its strict sign intervals.
\par \ExpectedValue{Positive expectation} means \ExpectedValue{expected net gain}; \ExpectedValue{negative expectation} means \ExpectedValue{expected net loss}; \ExpectedValue{zero expectation} means break-even. These are long-run averages per stated period, not guarantees of realized profit. Each branch uses its own payoffs and root.
\endgroup
}
\newcommand{\CThreeStatementManufacturing}{%
Observe a graduate's photography work over one quarter. The university graduation ceremony either goes ahead or is cancelled during that quarter. Compare a basic photo package with a package including extra editing. Both involve the same advance preparation costs; the edited package earns more when the ceremony takes place. \Money{Payoffs} include travel, software, preparation and refunds. The two possible states are \HLMathEvent{the ceremony goes ahead} and \HLMathEvent{the ceremony is cancelled}. The \SuccessProb{initial probability} that \HLMathEvent{the ceremony goes ahead} is $\SuccessProb{p}=\SuccessProb{\frac{3}{5}}$. All entries below are \Money{net profits} in \Money{USD} for the entire stated period, including relevant costs.
\begin{center}\begin{tabular}{lrr}\toprule Alternative & \HLMathEvent{the ceremony is cancelled} & \HLMathEvent{the ceremony goes ahead}\\\midrule
$A_1$: basic photo package & \Money{-100} & \Money{200}\\
$A_2$: edited photo package & \Money{-100} & \Money{500}\\
\bottomrule\end{tabular}\end{center}

}
\newcommand{\CThreeModelManufacturing}{%
\clearpage\noindent\textbf{Task 1 summary: shared states and induced distributions.}\par\medskip

\begingroup\small\setlength{\parskip}{2pt}\tikzset{gten marker/.append style={inner xsep=.4mm,inner ysep=.2mm}}
\begin{center}\begin{tikzpicture}[x=1cm,y=1cm]
\node[c3card,anchor=north,text width=7cm,align=center] (ev1) at (4,0) {$\HLMathEvent{E}$: \HLMathEvent{the ceremony goes ahead}, $\SuccessProb{P}(\HLMathEvent{E})=\SuccessProb{p}$};

\node[c3card,anchor=north,text width=7cm,align=center] (ev0) at (15,0) {$\HLMathEvent{E^c}$: \HLMathEvent{the ceremony is cancelled}, $\FailureProb{P}(\HLMathEvent{E^c})=\FailureProb{1-p}$};

\node[c3card,anchor=north,text width=6cm,align=center] (ind1) at (4,-1.2) {$\Indicator{X(\omega)}=\Indicator{1}$, for $\HLMathEvent{\omega\in E}$};

\node[c3card,anchor=north,text width=6cm,align=center] (ind0) at (15,-1.2) {$\Indicator{X(\omega)}=\Indicator{0}$, for $\HLMathEvent{\omega\in E^c}$};

\draw[c3flow] (ev1) -- (ind1);\draw[c3flow] (ev0) -- (ind0);
\node[c3card,anchor=north,text width=16.2cm,align=center] (map1) at (9.5,-2.65) {\textbf{\Money{$Y_1$}: basic photo package}\par 
$\Money{Y_1(\omega)}=\TransformationFunction{f_1}(\Indicator{X(\omega)})=\Money{f_1(1)}=\Money{200}\ \Money{\mathrm{USD}}$, for $\HLMathEvent{\omega\in E}$\par 
$\Money{Y_1(\omega)}=\TransformationFunction{f_1}(\Indicator{X(\omega)})=\Money{f_1(0)}=\Money{-100}\ \Money{\mathrm{USD}}$, for $\HLMathEvent{\omega\in E^c}$
};

\draw[c3flow] (ind1.west) -- (.15,-1.65) |- ([yshift=-7mm]map1.north west);
\draw[c3flow,dashed] (ind0.east) -- (18.85,-1.65) |- ([yshift=-11.5mm]map1.north east);
\node[c3card,anchor=north,text width=16.2cm,align=center] (map2) at (9.5,-4.45) {\textbf{\Money{$Y_2$}: edited photo package}\par 
$\Money{Y_2(\omega)}=\TransformationFunction{f_2}(\Indicator{X(\omega)})=\Money{f_2(1)}=\Money{500}\ \Money{\mathrm{USD}}$, for $\HLMathEvent{\omega\in E}$\par 
$\Money{Y_2(\omega)}=\TransformationFunction{f_2}(\Indicator{X(\omega)})=\Money{f_2(0)}=\Money{-100}\ \Money{\mathrm{USD}}$, for $\HLMathEvent{\omega\in E^c}$
};

\draw[c3flow] (ind1.west) -- (.15,-1.65) |- ([yshift=-7mm]map2.north west);
\draw[c3flow,dashed] (ind0.east) -- (18.85,-1.65) |- ([yshift=-11.5mm]map2.north east);
\node[inner sep=1pt,anchor=north,text width=5.913333333333333cm,align=center] (head0) at (3.0666666666666664,-6.4) {$\Indicator{X}$};

\node[c3card,anchor=north,text width=5.913333333333333cm,align=center] (cdf0) at (3.0666666666666664,-7.050000000000001) {$\HLCumulativeProbability{F_{X}(t)}$\par $\begin{cases}\HLCumulativeProbability{0},&t<\Indicator{0},\\\FailureProb{1-p},&\Indicator{0}\leq t<\Indicator{1},\\\HLCumulativeProbability{1},&t\geq \Indicator{1}.\end{cases}$};

\node[c3card,anchor=north,text width=5.913333333333333cm,align=center] (pmf0) at (3.0666666666666664,-9.8) {$\HLProbability{m_{X}}$\par $\begin{array}{ll}\Indicator{0}:&\FailureProb{1-p}\\\Indicator{1}:&\SuccessProb{p}\end{array}$\par \HLProbability{zero elsewhere}};

\node[inner sep=0pt,anchor=north,text width=5.913333333333333cm,align=center] (graph0) at (3.0666666666666664,-12.100000000000001) {\CThreeDistributionPair{X}{0}{1}{3/5}{unitless}{\Indicator}};

\node[inner sep=1pt,anchor=north,text width=5.913333333333333cm,align=center] (head1) at (9.4,-6.4) {$\Money{Y_1}$ $\leftarrow\TransformationFunction{f_1}$};

\node[c3card,anchor=north,text width=5.913333333333333cm,align=center] (cdf1) at (9.4,-7.050000000000001) {$\HLCumulativeProbability{F_{Y_1}(t)}$\par $\begin{cases}\HLCumulativeProbability{0},&t<\Money{-100},\\\FailureProb{1-p},&\Money{-100}\leq t<\Money{200},\\\HLCumulativeProbability{1},&t\geq \Money{200}.\end{cases}$};

\node[c3card,anchor=north,text width=5.913333333333333cm,align=center] (pmf1) at (9.4,-9.8) {$\HLProbability{m_{Y_1}}$\par $\begin{array}{ll}\Money{-100}:&\FailureProb{1-p}\\\Money{200}:&\SuccessProb{p}\end{array}$\par \HLProbability{zero elsewhere}};

\node[inner sep=0pt,anchor=north,text width=5.913333333333333cm,align=center] (graph1) at (9.4,-12.100000000000001) {\CThreeDistributionPair{Y_1}{-100}{200}{3/5}{\Money{USD}}{\Money}};

\draw[c3flow] (cdf0.north) -- ++(0,.27) -| (cdf1.north);
\draw[c3flow] (pmf0.south) -- ++(0,-.30) -| (pmf1.south);
\draw[c3flow] (map1.east) -- (18.28,-3.3) -- (18.28,-6.1000000000000005) -| (head1.north);
\node[inner sep=1pt,anchor=north,text width=5.913333333333333cm,align=center] (head2) at (15.733333333333333,-6.4) {$\Money{Y_2}$ $\leftarrow\TransformationFunction{f_2}$};

\node[c3card,anchor=north,text width=5.913333333333333cm,align=center] (cdf2) at (15.733333333333333,-7.050000000000001) {$\HLCumulativeProbability{F_{Y_2}(t)}$\par $\begin{cases}\HLCumulativeProbability{0},&t<\Money{-100},\\\FailureProb{1-p},&\Money{-100}\leq t<\Money{500},\\\HLCumulativeProbability{1},&t\geq \Money{500}.\end{cases}$};

\node[c3card,anchor=north,text width=5.913333333333333cm,align=center] (pmf2) at (15.733333333333333,-9.8) {$\HLProbability{m_{Y_2}}$\par $\begin{array}{ll}\Money{-100}:&\FailureProb{1-p}\\\Money{500}:&\SuccessProb{p}\end{array}$\par \HLProbability{zero elsewhere}};

\node[inner sep=0pt,anchor=north,text width=5.913333333333333cm,align=center] (graph2) at (15.733333333333333,-12.100000000000001) {\CThreeDistributionPair{Y_2}{-100}{500}{3/5}{\Money{USD}}{\Money}};

\draw[c3flow] (cdf0.north) -- ++(0,.27) -| (cdf2.north);
\draw[c3flow] (pmf0.south) -- ++(0,-.30) -| (pmf2.south);
\draw[c3flow] (map2.east) -- (18.1,-5.1) -- (18.1,-5.98) -| (head2.north);
\end{tikzpicture}\end{center}
\textbf{Numerical graph illustration:} $\SuccessProb{p_0}=\SuccessProb{\frac{3}{5}}$. All displayed distribution formulas remain symbolic in $\SuccessProb{p}$.
\par\textbf{CDF above, PMF below.} Filled stems translate to pale dashed jump shadows; thin arrows link equal physical lengths. Open and filled circles show excluded and included endpoints. Distribution arrows retain state probabilities while mapping values, rather than applying $\TransformationFunction{f_i}$ to CDF heights. See \hyperref[ex:intro-one]{Example 1}.\par
For $0<\SuccessProb{p}<1$, $\Indicator{S_X}=\{\Indicator{0},\Indicator{1}\}$; at $\SuccessProb{p}=0$ only $\Indicator{0}$ remains, and at $\SuccessProb{p}=1$ only $\Indicator{1}$ remains. Each monetary support retains $\Money{f_i(0)}$ at $\SuccessProb{p}=0$ and $\Money{f_i(1)}$ at $\SuccessProb{p}=1$.
\par\textbf{Loss events.} $\HLMathEvent{\{Y_1<0\}}=\HLMathEvent{\{X=0\}}=\HLMathEvent{E^c}$, probability $\FailureProb{1-p}$; $\HLMathEvent{\{Y_2<0\}}=\HLMathEvent{\{X=0\}}=\HLMathEvent{E^c}$, probability $\FailureProb{1-p}$.
\endgroup
}
\newcommand{\CThreeExpectManufacturing}{%
\clearpage\noindent\textbf{Task 2 summary: parallel expectation analyses.}\par\medskip

\begingroup\small\tikzset{gten marker/.append style={inner xsep=.4mm,inner ysep=.2mm}}
\begin{center}\begin{tikzpicture}[x=1cm,y=1cm]
\node[c3card,anchor=north,text width=17cm,align=center] (inputs) at (9.5,0) {Shared probability inputs: $\SuccessProb{p}$ and $\FailureProb{1-p}$; hold each alternative\textquotesingle s \Money{payoffs} fixed.};

\node[c3card,anchor=north,text width=17cm,align=center] (formula) at (9.5,-1.2) {$\ExpectedValue{g_i}(\SuccessProb{p})=\ExpectedValue{\mathbb E}[\Money{Y_i}]=\SuccessProb{p}\Money{f_i(1)}+(\FailureProb{1-p})\Money{f_i(0)}$};

\draw[c3flow] (inputs) -- (formula);
\node[c3card,anchor=north,text width=8.8cm,align=center] (pay1) at (4.75,-2.7) {\textbf{\Money{$Y_1$}: basic photo package}\par $\Money{f_1(0)}=\Money{-100},\quad\Money{f_1(1)}=\Money{200}\ \Money{\mathrm{USD}}$\par $\ExpectedValue{g_1}(\SuccessProb{p})=\Money{-100}+\Money{300}\SuccessProb{p}$};

\draw[c3flow] (formula.south) -- ++(0,-.25) -| (pay1.north);
\node[c3card,anchor=north,text width=8.8cm,align=center] (initial1) at (4.75,-5.0) {\ExpectedValue{Initial expectation}\par $\ExpectedValue{g_1}\!\left(\SuccessProb{\frac{3}{5}}\right)=\ExpectedValue{80\ \mathrm{USD}}$\par \ExpectedValue{expected net gain} per stated period};

\node[c3card,anchor=north,text width=8.8cm,align=center] (root1) at (4.75,-7.15) {\SuccessProb{Exact break-even equation}\par $\Money{-100}+\Money{300}\SuccessProb{p}=\ExpectedValue{0}$\par $\SuccessProb{p_1^*}=\SuccessProb{\frac{1}{3}}$};

\node[c3card,anchor=north,text width=8.8cm,align=center] (sign1) at (4.75,-9.45) {$\ExpectedValue{g_1}>\ExpectedValue{0}:\ \SuccessProb{p}\in \left(\SuccessProb{\frac{1}{3}},\SuccessProb{1}\right]$\par $\ExpectedValue{g_1}=\ExpectedValue{0}:\ \SuccessProb{p}=\SuccessProb{\frac{1}{3}}$\par $\ExpectedValue{g_1}<\ExpectedValue{0}:\ \SuccessProb{p}\in \left[\SuccessProb{0},\SuccessProb{\frac{1}{3}}\right)$};

\node[inner sep=0pt,anchor=north,text width=8.8cm,align=center] (eg1) at (4.75,-12.0) {\CThreeExpectationGraph{1}{-100}{200}{1/3}{3/5}};

\draw[c3flow] (pay1.south) -- (initial1.north);
\draw[c3flow] (initial1.south) -- (root1.north);
\draw[c3flow] (root1.south) -- (sign1.north);
\draw[c3flow] (sign1.south) -- (eg1.north);
\node[c3card,anchor=north,text width=8.8cm,align=center] (pay2) at (14.25,-2.7) {\textbf{\Money{$Y_2$}: edited photo package}\par $\Money{f_2(0)}=\Money{-100},\quad\Money{f_2(1)}=\Money{500}\ \Money{\mathrm{USD}}$\par $\ExpectedValue{g_2}(\SuccessProb{p})=\Money{-100}+\Money{600}\SuccessProb{p}$};

\draw[c3flow] (formula.south) -- ++(0,-.25) -| (pay2.north);
\node[c3card,anchor=north,text width=8.8cm,align=center] (initial2) at (14.25,-5.0) {\ExpectedValue{Initial expectation}\par $\ExpectedValue{g_2}\!\left(\SuccessProb{\frac{3}{5}}\right)=\ExpectedValue{260\ \mathrm{USD}}$\par \ExpectedValue{expected net gain} per stated period};

\node[c3card,anchor=north,text width=8.8cm,align=center] (root2) at (14.25,-7.15) {\SuccessProb{Exact break-even equation}\par $\Money{-100}+\Money{600}\SuccessProb{p}=\ExpectedValue{0}$\par $\SuccessProb{p_2^*}=\SuccessProb{\frac{1}{6}}$};

\node[c3card,anchor=north,text width=8.8cm,align=center] (sign2) at (14.25,-9.45) {$\ExpectedValue{g_2}>\ExpectedValue{0}:\ \SuccessProb{p}\in \left(\SuccessProb{\frac{1}{6}},\SuccessProb{1}\right]$\par $\ExpectedValue{g_2}=\ExpectedValue{0}:\ \SuccessProb{p}=\SuccessProb{\frac{1}{6}}$\par $\ExpectedValue{g_2}<\ExpectedValue{0}:\ \SuccessProb{p}\in \left[\SuccessProb{0},\SuccessProb{\frac{1}{6}}\right)$};

\node[inner sep=0pt,anchor=north,text width=8.8cm,align=center] (eg2) at (14.25,-12.0) {\CThreeExpectationGraph{2}{-100}{500}{1/6}{3/5}};

\draw[c3flow] (pay2.south) -- (initial2.north);
\draw[c3flow] (initial2.south) -- (root2.north);
\draw[c3flow] (root2.south) -- (sign2.north);
\draw[c3flow] (sign2.south) -- (eg2.north);
\end{tikzpicture}\end{center}
The dashed horizontal line is \ExpectedValue{zero}; filled circles show endpoint payoffs and the exact root. The square marks the \ExpectedValue{initial expectation}. The $+$ and $-$ labels connect the graph to its strict sign intervals.
\par \ExpectedValue{Positive expectation} means \ExpectedValue{expected net gain}; \ExpectedValue{negative expectation} means \ExpectedValue{expected net loss}; \ExpectedValue{zero expectation} means break-even. These are long-run averages per stated period, not guarantees of realized profit. Each branch uses its own payoffs and root.
\endgroup
}
\newcommand{\CThreeStatementRestaurant}{%
Observe a student group's ticketed concert over one evening. \HLMathEvent{High demand} means advance expressions of interest meet a fixed target; \HLMathEvent{low demand} means they fall below it. Compare hiring a small room, a discounted student-union room or a large room for the same concert. \Money{Payoffs} include venue hire, performers, ticket fees and refunds. The two possible states are \HLMathEvent{high demand} and \HLMathEvent{low demand}. The \SuccessProb{initial probability} of \HLMathEvent{high demand} is $\SuccessProb{p}=\SuccessProb{\frac{9}{25}}$. All entries below are \Money{net profits} in \Money{USD} for the entire stated period, including relevant costs.
\begin{center}\begin{tabular}{lrr}\toprule Alternative & \HLMathEvent{low demand} & \HLMathEvent{high demand}\\\midrule
$A_1$: small hired room & \Money{-80} & \Money{100}\\
$A_2$: student-union room & \Money{-20} & \Money{220}\\
$A_3$: large hired room & \Money{-100} & \Money{400}\\
\bottomrule\end{tabular}\end{center}

}
\newcommand{\CThreeModelRestaurant}{%
\clearpage\noindent\textbf{Task 1 summary: shared states and induced distributions.}\par\medskip

\begingroup\small\setlength{\parskip}{2pt}\tikzset{gten marker/.append style={inner xsep=.4mm,inner ysep=.2mm}}
\begin{center}\begin{tikzpicture}[x=1cm,y=1cm]
\node[c3card,anchor=north,text width=7cm,align=center] (ev1) at (4,0) {$\HLMathEvent{E}$: \HLMathEvent{high demand}, $\SuccessProb{P}(\HLMathEvent{E})=\SuccessProb{p}$};

\node[c3card,anchor=north,text width=7cm,align=center] (ev0) at (15,0) {$\HLMathEvent{E^c}$: \HLMathEvent{low demand}, $\FailureProb{P}(\HLMathEvent{E^c})=\FailureProb{1-p}$};

\node[c3card,anchor=north,text width=6cm,align=center] (ind1) at (4,-1.2) {$\Indicator{X(\omega)}=\Indicator{1}$, for $\HLMathEvent{\omega\in E}$};

\node[c3card,anchor=north,text width=6cm,align=center] (ind0) at (15,-1.2) {$\Indicator{X(\omega)}=\Indicator{0}$, for $\HLMathEvent{\omega\in E^c}$};

\draw[c3flow] (ev1) -- (ind1);\draw[c3flow] (ev0) -- (ind0);
\node[c3card,anchor=north,text width=16.2cm,align=center] (map1) at (9.5,-2.65) {\textbf{\Money{$Y_1$}: small hired room}\par 
$\Money{Y_1(\omega)}=\TransformationFunction{f_1}(\Indicator{X(\omega)})=\Money{f_1(1)}=\Money{100}\ \Money{\mathrm{USD}}$, for $\HLMathEvent{\omega\in E}$\par 
$\Money{Y_1(\omega)}=\TransformationFunction{f_1}(\Indicator{X(\omega)})=\Money{f_1(0)}=\Money{-80}\ \Money{\mathrm{USD}}$, for $\HLMathEvent{\omega\in E^c}$
};

\draw[c3flow] (ind1.west) -- (.15,-1.65) |- ([yshift=-7mm]map1.north west);
\draw[c3flow,dashed] (ind0.east) -- (18.85,-1.65) |- ([yshift=-11.5mm]map1.north east);
\node[c3card,anchor=north,text width=16.2cm,align=center] (map2) at (9.5,-4.45) {\textbf{\Money{$Y_2$}: student-union room}\par 
$\Money{Y_2(\omega)}=\TransformationFunction{f_2}(\Indicator{X(\omega)})=\Money{f_2(1)}=\Money{220}\ \Money{\mathrm{USD}}$, for $\HLMathEvent{\omega\in E}$\par 
$\Money{Y_2(\omega)}=\TransformationFunction{f_2}(\Indicator{X(\omega)})=\Money{f_2(0)}=\Money{-20}\ \Money{\mathrm{USD}}$, for $\HLMathEvent{\omega\in E^c}$
};

\draw[c3flow] (ind1.west) -- (.15,-1.65) |- ([yshift=-7mm]map2.north west);
\draw[c3flow,dashed] (ind0.east) -- (18.85,-1.65) |- ([yshift=-11.5mm]map2.north east);
\node[c3card,anchor=north,text width=16.2cm,align=center] (map3) at (9.5,-6.25) {\textbf{\Money{$Y_3$}: large hired room}\par 
$\Money{Y_3(\omega)}=\TransformationFunction{f_3}(\Indicator{X(\omega)})=\Money{f_3(1)}=\Money{400}\ \Money{\mathrm{USD}}$, for $\HLMathEvent{\omega\in E}$\par 
$\Money{Y_3(\omega)}=\TransformationFunction{f_3}(\Indicator{X(\omega)})=\Money{f_3(0)}=\Money{-100}\ \Money{\mathrm{USD}}$, for $\HLMathEvent{\omega\in E^c}$
};

\draw[c3flow] (ind1.west) -- (.15,-1.65) |- ([yshift=-7mm]map3.north west);
\draw[c3flow,dashed] (ind0.east) -- (18.85,-1.65) |- ([yshift=-11.5mm]map3.north east);
\node[inner sep=1pt,anchor=north,text width=4.33cm,align=center] (head0) at (2.275,-8.2) {$\Indicator{X}$};

\node[c3card,anchor=north,text width=4.33cm,align=center] (cdf0) at (2.275,-8.85) {$\HLCumulativeProbability{F_{X}(t)}$\par $\begin{cases}\HLCumulativeProbability{0},&t<\Indicator{0},\\\FailureProb{1-p},&\Indicator{0}\leq t<\Indicator{1},\\\HLCumulativeProbability{1},&t\geq \Indicator{1}.\end{cases}$};

\node[c3card,anchor=north,text width=4.33cm,align=center] (pmf0) at (2.275,-11.6) {$\HLProbability{m_{X}}$\par $\begin{array}{ll}\Indicator{0}:&\FailureProb{1-p}\\\Indicator{1}:&\SuccessProb{p}\end{array}$\par \HLProbability{zero elsewhere}};

\node[inner sep=0pt,anchor=north,text width=4.33cm,align=center] (graph0) at (2.275,-13.899999999999999) {\CThreeDistributionPair{X}{0}{1}{9/25}{unitless}{\Indicator}};

\node[inner sep=1pt,anchor=north,text width=4.33cm,align=center] (head1) at (7.025,-8.2) {$\Money{Y_1}$ $\leftarrow\TransformationFunction{f_1}$};

\node[c3card,anchor=north,text width=4.33cm,align=center] (cdf1) at (7.025,-8.85) {$\HLCumulativeProbability{F_{Y_1}(t)}$\par $\begin{cases}\HLCumulativeProbability{0},&t<\Money{-80},\\\FailureProb{1-p},&\Money{-80}\leq t<\Money{100},\\\HLCumulativeProbability{1},&t\geq \Money{100}.\end{cases}$};

\node[c3card,anchor=north,text width=4.33cm,align=center] (pmf1) at (7.025,-11.6) {$\HLProbability{m_{Y_1}}$\par $\begin{array}{ll}\Money{-80}:&\FailureProb{1-p}\\\Money{100}:&\SuccessProb{p}\end{array}$\par \HLProbability{zero elsewhere}};

\node[inner sep=0pt,anchor=north,text width=4.33cm,align=center] (graph1) at (7.025,-13.899999999999999) {\CThreeDistributionPair{Y_1}{-80}{100}{9/25}{\Money{USD}}{\Money}};

\draw[c3flow] (cdf0.north) -- ++(0,.27) -| (cdf1.north);
\draw[c3flow] (pmf0.south) -- ++(0,-.30) -| (pmf1.south);
\draw[c3flow] (map1.east) -- (18.46,-3.3) -- (18.46,-7.8999999999999995) -| (head1.north);
\node[inner sep=1pt,anchor=north,text width=4.33cm,align=center] (head2) at (11.775,-8.2) {$\Money{Y_2}$ $\leftarrow\TransformationFunction{f_2}$};

\node[c3card,anchor=north,text width=4.33cm,align=center] (cdf2) at (11.775,-8.85) {$\HLCumulativeProbability{F_{Y_2}(t)}$\par $\begin{cases}\HLCumulativeProbability{0},&t<\Money{-20},\\\FailureProb{1-p},&\Money{-20}\leq t<\Money{220},\\\HLCumulativeProbability{1},&t\geq \Money{220}.\end{cases}$};

\node[c3card,anchor=north,text width=4.33cm,align=center] (pmf2) at (11.775,-11.6) {$\HLProbability{m_{Y_2}}$\par $\begin{array}{ll}\Money{-20}:&\FailureProb{1-p}\\\Money{220}:&\SuccessProb{p}\end{array}$\par \HLProbability{zero elsewhere}};

\node[inner sep=0pt,anchor=north,text width=4.33cm,align=center] (graph2) at (11.775,-13.899999999999999) {\CThreeDistributionPair{Y_2}{-20}{220}{9/25}{\Money{USD}}{\Money}};

\draw[c3flow] (cdf0.north) -- ++(0,.27) -| (cdf2.north);
\draw[c3flow] (pmf0.south) -- ++(0,-.30) -| (pmf2.south);
\draw[c3flow] (map2.east) -- (18.28,-5.1) -- (18.28,-7.779999999999999) -| (head2.north);
\node[inner sep=1pt,anchor=north,text width=4.33cm,align=center] (head3) at (16.525,-8.2) {$\Money{Y_3}$ $\leftarrow\TransformationFunction{f_3}$};

\node[c3card,anchor=north,text width=4.33cm,align=center] (cdf3) at (16.525,-8.85) {$\HLCumulativeProbability{F_{Y_3}(t)}$\par $\begin{cases}\HLCumulativeProbability{0},&t<\Money{-100},\\\FailureProb{1-p},&\Money{-100}\leq t<\Money{400},\\\HLCumulativeProbability{1},&t\geq \Money{400}.\end{cases}$};

\node[c3card,anchor=north,text width=4.33cm,align=center] (pmf3) at (16.525,-11.6) {$\HLProbability{m_{Y_3}}$\par $\begin{array}{ll}\Money{-100}:&\FailureProb{1-p}\\\Money{400}:&\SuccessProb{p}\end{array}$\par \HLProbability{zero elsewhere}};

\node[inner sep=0pt,anchor=north,text width=4.33cm,align=center] (graph3) at (16.525,-13.899999999999999) {\CThreeDistributionPair{Y_3}{-100}{400}{9/25}{\Money{USD}}{\Money}};

\draw[c3flow] (cdf0.north) -- ++(0,.27) -| (cdf3.north);
\draw[c3flow] (pmf0.south) -- ++(0,-.30) -| (pmf3.south);
\draw[c3flow] (map3.east) -- (18.1,-6.9) -- (18.1,-7.659999999999999) -| (head3.north);
\end{tikzpicture}\end{center}
\textbf{Numerical graph illustration:} $\SuccessProb{p_0}=\SuccessProb{\frac{9}{25}}$. All displayed distribution formulas remain symbolic in $\SuccessProb{p}$.
\par\textbf{CDF above, PMF below.} Filled stems translate to pale dashed jump shadows; thin arrows link equal physical lengths. Open and filled circles show excluded and included endpoints. Distribution arrows retain state probabilities while mapping values, rather than applying $\TransformationFunction{f_i}$ to CDF heights. See \hyperref[ex:intro-one]{Example 1}.\par
For $0<\SuccessProb{p}<1$, $\Indicator{S_X}=\{\Indicator{0},\Indicator{1}\}$; at $\SuccessProb{p}=0$ only $\Indicator{0}$ remains, and at $\SuccessProb{p}=1$ only $\Indicator{1}$ remains. Each monetary support retains $\Money{f_i(0)}$ at $\SuccessProb{p}=0$ and $\Money{f_i(1)}$ at $\SuccessProb{p}=1$.
\par\textbf{Loss events.} $\HLMathEvent{\{Y_1<0\}}=\HLMathEvent{\{X=0\}}=\HLMathEvent{E^c}$, probability $\FailureProb{1-p}$; $\HLMathEvent{\{Y_2<0\}}=\HLMathEvent{\{X=0\}}=\HLMathEvent{E^c}$, probability $\FailureProb{1-p}$; $\HLMathEvent{\{Y_3<0\}}=\HLMathEvent{\{X=0\}}=\HLMathEvent{E^c}$, probability $\FailureProb{1-p}$.
\endgroup
}
\newcommand{\CThreeExpectRestaurant}{%
\clearpage\noindent\textbf{Task 2 summary: parallel expectation analyses.}\par\medskip

\begingroup\small\tikzset{gten marker/.append style={inner xsep=.4mm,inner ysep=.2mm}}
\begin{center}\begin{tikzpicture}[x=1cm,y=1cm]
\node[c3card,anchor=north,text width=17cm,align=center] (inputs) at (9.5,0) {Shared probability inputs: $\SuccessProb{p}$ and $\FailureProb{1-p}$; hold each alternative\textquotesingle s \Money{payoffs} fixed.};

\node[c3card,anchor=north,text width=17cm,align=center] (formula) at (9.5,-1.2) {$\ExpectedValue{g_i}(\SuccessProb{p})=\ExpectedValue{\mathbb E}[\Money{Y_i}]=\SuccessProb{p}\Money{f_i(1)}+(\FailureProb{1-p})\Money{f_i(0)}$};

\draw[c3flow] (inputs) -- (formula);
\node[c3card,anchor=north,text width=5.633333333333333cm,align=center] (pay1) at (3.1666666666666665,-2.7) {\textbf{\Money{$Y_1$}: small hired room}\par $\Money{f_1(0)}=\Money{-80},\quad\Money{f_1(1)}=\Money{100}\ \Money{\mathrm{USD}}$\par $\ExpectedValue{g_1}(\SuccessProb{p})=\Money{-80}+\Money{180}\SuccessProb{p}$};

\draw[c3flow] (formula.south) -- ++(0,-.25) -| (pay1.north);
\node[c3card,anchor=north,text width=5.633333333333333cm,align=center] (initial1) at (3.1666666666666665,-5.0) {\ExpectedValue{Initial expectation}\par $\ExpectedValue{g_1}\!\left(\SuccessProb{\frac{9}{25}}\right)=\ExpectedValue{\frac{-76}{5}\ \mathrm{USD}}$\par \ExpectedValue{expected net loss} per stated period};

\node[c3card,anchor=north,text width=5.633333333333333cm,align=center] (root1) at (3.1666666666666665,-7.15) {\SuccessProb{Exact break-even equation}\par $\Money{-80}+\Money{180}\SuccessProb{p}=\ExpectedValue{0}$\par $\SuccessProb{p_1^*}=\SuccessProb{\frac{4}{9}}$};

\node[c3card,anchor=north,text width=5.633333333333333cm,align=center] (sign1) at (3.1666666666666665,-9.45) {$\ExpectedValue{g_1}>\ExpectedValue{0}:\ \SuccessProb{p}\in \left(\SuccessProb{\frac{4}{9}},\SuccessProb{1}\right]$\par $\ExpectedValue{g_1}=\ExpectedValue{0}:\ \SuccessProb{p}=\SuccessProb{\frac{4}{9}}$\par $\ExpectedValue{g_1}<\ExpectedValue{0}:\ \SuccessProb{p}\in \left[\SuccessProb{0},\SuccessProb{\frac{4}{9}}\right)$};

\node[inner sep=0pt,anchor=north,text width=5.633333333333333cm,align=center] (eg1) at (3.1666666666666665,-12.0) {\CThreeExpectationGraph{1}{-80}{100}{4/9}{9/25}};

\draw[c3flow] (pay1.south) -- (initial1.north);
\draw[c3flow] (initial1.south) -- (root1.north);
\draw[c3flow] (root1.south) -- (sign1.north);
\draw[c3flow] (sign1.south) -- (eg1.north);
\node[c3card,anchor=north,text width=5.633333333333333cm,align=center] (pay2) at (9.5,-2.7) {\textbf{\Money{$Y_2$}: student-union room}\par $\Money{f_2(0)}=\Money{-20},\quad\Money{f_2(1)}=\Money{220}\ \Money{\mathrm{USD}}$\par $\ExpectedValue{g_2}(\SuccessProb{p})=\Money{-20}+\Money{240}\SuccessProb{p}$};

\draw[c3flow] (formula.south) -- ++(0,-.25) -| (pay2.north);
\node[c3card,anchor=north,text width=5.633333333333333cm,align=center] (initial2) at (9.5,-5.0) {\ExpectedValue{Initial expectation}\par $\ExpectedValue{g_2}\!\left(\SuccessProb{\frac{9}{25}}\right)=\ExpectedValue{\frac{332}{5}\ \mathrm{USD}}$\par \ExpectedValue{expected net gain} per stated period};

\node[c3card,anchor=north,text width=5.633333333333333cm,align=center] (root2) at (9.5,-7.15) {\SuccessProb{Exact break-even equation}\par $\Money{-20}+\Money{240}\SuccessProb{p}=\ExpectedValue{0}$\par $\SuccessProb{p_2^*}=\SuccessProb{\frac{1}{12}}$};

\node[c3card,anchor=north,text width=5.633333333333333cm,align=center] (sign2) at (9.5,-9.45) {$\ExpectedValue{g_2}>\ExpectedValue{0}:\ \SuccessProb{p}\in \left(\SuccessProb{\frac{1}{12}},\SuccessProb{1}\right]$\par $\ExpectedValue{g_2}=\ExpectedValue{0}:\ \SuccessProb{p}=\SuccessProb{\frac{1}{12}}$\par $\ExpectedValue{g_2}<\ExpectedValue{0}:\ \SuccessProb{p}\in \left[\SuccessProb{0},\SuccessProb{\frac{1}{12}}\right)$};

\node[inner sep=0pt,anchor=north,text width=5.633333333333333cm,align=center] (eg2) at (9.5,-12.0) {\CThreeExpectationGraph{2}{-20}{220}{1/12}{9/25}};

\draw[c3flow] (pay2.south) -- (initial2.north);
\draw[c3flow] (initial2.south) -- (root2.north);
\draw[c3flow] (root2.south) -- (sign2.north);
\draw[c3flow] (sign2.south) -- (eg2.north);
\node[c3card,anchor=north,text width=5.633333333333333cm,align=center] (pay3) at (15.833333333333332,-2.7) {\textbf{\Money{$Y_3$}: large hired room}\par $\Money{f_3(0)}=\Money{-100},\quad\Money{f_3(1)}=\Money{400}\ \Money{\mathrm{USD}}$\par $\ExpectedValue{g_3}(\SuccessProb{p})=\Money{-100}+\Money{500}\SuccessProb{p}$};

\draw[c3flow] (formula.south) -- ++(0,-.25) -| (pay3.north);
\node[c3card,anchor=north,text width=5.633333333333333cm,align=center] (initial3) at (15.833333333333332,-5.0) {\ExpectedValue{Initial expectation}\par $\ExpectedValue{g_3}\!\left(\SuccessProb{\frac{9}{25}}\right)=\ExpectedValue{80\ \mathrm{USD}}$\par \ExpectedValue{expected net gain} per stated period};

\node[c3card,anchor=north,text width=5.633333333333333cm,align=center] (root3) at (15.833333333333332,-7.15) {\SuccessProb{Exact break-even equation}\par $\Money{-100}+\Money{500}\SuccessProb{p}=\ExpectedValue{0}$\par $\SuccessProb{p_3^*}=\SuccessProb{\frac{1}{5}}$};

\node[c3card,anchor=north,text width=5.633333333333333cm,align=center] (sign3) at (15.833333333333332,-9.45) {$\ExpectedValue{g_3}>\ExpectedValue{0}:\ \SuccessProb{p}\in \left(\SuccessProb{\frac{1}{5}},\SuccessProb{1}\right]$\par $\ExpectedValue{g_3}=\ExpectedValue{0}:\ \SuccessProb{p}=\SuccessProb{\frac{1}{5}}$\par $\ExpectedValue{g_3}<\ExpectedValue{0}:\ \SuccessProb{p}\in \left[\SuccessProb{0},\SuccessProb{\frac{1}{5}}\right)$};

\node[inner sep=0pt,anchor=north,text width=5.633333333333333cm,align=center] (eg3) at (15.833333333333332,-12.0) {\CThreeExpectationGraph{3}{-100}{400}{1/5}{9/25}};

\draw[c3flow] (pay3.south) -- (initial3.north);
\draw[c3flow] (initial3.south) -- (root3.north);
\draw[c3flow] (root3.south) -- (sign3.north);
\draw[c3flow] (sign3.south) -- (eg3.north);
\end{tikzpicture}\end{center}
The dashed horizontal line is \ExpectedValue{zero}; filled circles show endpoint payoffs and the exact root. The square marks the \ExpectedValue{initial expectation}. The $+$ and $-$ labels connect the graph to its strict sign intervals.
\par \ExpectedValue{Positive expectation} means \ExpectedValue{expected net gain}; \ExpectedValue{negative expectation} means \ExpectedValue{expected net loss}; \ExpectedValue{zero expectation} means break-even. These are long-run averages per stated period, not guarantees of realized profit. Each branch uses its own payoffs and root.
\endgroup
}
\newcommand{\CThreeStatementAgriculture}{%
Observe a student's paid fitness sessions over one university term. \HLMathEvent{High enrollment} means responses to a campus fitness-interest survey meet a fixed target; \HLMathEvent{low enrollment} means they fall below it. Compare a walking group, dance sessions or circuit sessions, using the same survey outcome for each hypothetical format. \Money{Payoffs} include space hire, supplies, promotion and refunds. The two possible states are \HLMathEvent{high enrollment} and \HLMathEvent{low enrollment}. The \SuccessProb{initial probability} of \HLMathEvent{high enrollment} is $\SuccessProb{p}=\SuccessProb{\frac{1}{2}}$. All entries below are \Money{net profits} in \Money{USD} for the entire stated period, including relevant costs.
\begin{center}\begin{tabular}{lrr}\toprule Alternative & \HLMathEvent{low enrollment} & \HLMathEvent{high enrollment}\\\midrule
$A_1$: walking group & \Money{-70} & \Money{90}\\
$A_2$: dance sessions & \Money{-30} & \Money{190}\\
$A_3$: circuit sessions & \Money{-150} & \Money{150}\\
\bottomrule\end{tabular}\end{center}

}
\newcommand{\CThreeModelAgriculture}{%
\clearpage\noindent\textbf{Task 1 summary: shared states and induced distributions.}\par\medskip

\begingroup\small\setlength{\parskip}{2pt}\tikzset{gten marker/.append style={inner xsep=.4mm,inner ysep=.2mm}}
\begin{center}\begin{tikzpicture}[x=1cm,y=1cm]
\node[c3card,anchor=north,text width=7cm,align=center] (ev1) at (4,0) {$\HLMathEvent{E}$: \HLMathEvent{high enrollment}, $\SuccessProb{P}(\HLMathEvent{E})=\SuccessProb{p}$};

\node[c3card,anchor=north,text width=7cm,align=center] (ev0) at (15,0) {$\HLMathEvent{E^c}$: \HLMathEvent{low enrollment}, $\FailureProb{P}(\HLMathEvent{E^c})=\FailureProb{1-p}$};

\node[c3card,anchor=north,text width=6cm,align=center] (ind1) at (4,-1.2) {$\Indicator{X(\omega)}=\Indicator{1}$, for $\HLMathEvent{\omega\in E}$};

\node[c3card,anchor=north,text width=6cm,align=center] (ind0) at (15,-1.2) {$\Indicator{X(\omega)}=\Indicator{0}$, for $\HLMathEvent{\omega\in E^c}$};

\draw[c3flow] (ev1) -- (ind1);\draw[c3flow] (ev0) -- (ind0);
\node[c3card,anchor=north,text width=16.2cm,align=center] (map1) at (9.5,-2.65) {\textbf{\Money{$Y_1$}: walking group}\par 
$\Money{Y_1(\omega)}=\TransformationFunction{f_1}(\Indicator{X(\omega)})=\Money{f_1(1)}=\Money{90}\ \Money{\mathrm{USD}}$, for $\HLMathEvent{\omega\in E}$\par 
$\Money{Y_1(\omega)}=\TransformationFunction{f_1}(\Indicator{X(\omega)})=\Money{f_1(0)}=\Money{-70}\ \Money{\mathrm{USD}}$, for $\HLMathEvent{\omega\in E^c}$
};

\draw[c3flow] (ind1.west) -- (.15,-1.65) |- ([yshift=-7mm]map1.north west);
\draw[c3flow,dashed] (ind0.east) -- (18.85,-1.65) |- ([yshift=-11.5mm]map1.north east);
\node[c3card,anchor=north,text width=16.2cm,align=center] (map2) at (9.5,-4.45) {\textbf{\Money{$Y_2$}: dance sessions}\par 
$\Money{Y_2(\omega)}=\TransformationFunction{f_2}(\Indicator{X(\omega)})=\Money{f_2(1)}=\Money{190}\ \Money{\mathrm{USD}}$, for $\HLMathEvent{\omega\in E}$\par 
$\Money{Y_2(\omega)}=\TransformationFunction{f_2}(\Indicator{X(\omega)})=\Money{f_2(0)}=\Money{-30}\ \Money{\mathrm{USD}}$, for $\HLMathEvent{\omega\in E^c}$
};

\draw[c3flow] (ind1.west) -- (.15,-1.65) |- ([yshift=-7mm]map2.north west);
\draw[c3flow,dashed] (ind0.east) -- (18.85,-1.65) |- ([yshift=-11.5mm]map2.north east);
\node[c3card,anchor=north,text width=16.2cm,align=center] (map3) at (9.5,-6.25) {\textbf{\Money{$Y_3$}: circuit sessions}\par 
$\Money{Y_3(\omega)}=\TransformationFunction{f_3}(\Indicator{X(\omega)})=\Money{f_3(1)}=\Money{150}\ \Money{\mathrm{USD}}$, for $\HLMathEvent{\omega\in E}$\par 
$\Money{Y_3(\omega)}=\TransformationFunction{f_3}(\Indicator{X(\omega)})=\Money{f_3(0)}=\Money{-150}\ \Money{\mathrm{USD}}$, for $\HLMathEvent{\omega\in E^c}$
};

\draw[c3flow] (ind1.west) -- (.15,-1.65) |- ([yshift=-7mm]map3.north west);
\draw[c3flow,dashed] (ind0.east) -- (18.85,-1.65) |- ([yshift=-11.5mm]map3.north east);
\node[inner sep=1pt,anchor=north,text width=4.33cm,align=center] (head0) at (2.275,-8.2) {$\Indicator{X}$};

\node[c3card,anchor=north,text width=4.33cm,align=center] (cdf0) at (2.275,-8.85) {$\HLCumulativeProbability{F_{X}(t)}$\par $\begin{cases}\HLCumulativeProbability{0},&t<\Indicator{0},\\\FailureProb{1-p},&\Indicator{0}\leq t<\Indicator{1},\\\HLCumulativeProbability{1},&t\geq \Indicator{1}.\end{cases}$};

\node[c3card,anchor=north,text width=4.33cm,align=center] (pmf0) at (2.275,-11.6) {$\HLProbability{m_{X}}$\par $\begin{array}{ll}\Indicator{0}:&\FailureProb{1-p}\\\Indicator{1}:&\SuccessProb{p}\end{array}$\par \HLProbability{zero elsewhere}};

\node[inner sep=0pt,anchor=north,text width=4.33cm,align=center] (graph0) at (2.275,-13.899999999999999) {\CThreeDistributionPair{X}{0}{1}{1/2}{unitless}{\Indicator}};

\node[inner sep=1pt,anchor=north,text width=4.33cm,align=center] (head1) at (7.025,-8.2) {$\Money{Y_1}$ $\leftarrow\TransformationFunction{f_1}$};

\node[c3card,anchor=north,text width=4.33cm,align=center] (cdf1) at (7.025,-8.85) {$\HLCumulativeProbability{F_{Y_1}(t)}$\par $\begin{cases}\HLCumulativeProbability{0},&t<\Money{-70},\\\FailureProb{1-p},&\Money{-70}\leq t<\Money{90},\\\HLCumulativeProbability{1},&t\geq \Money{90}.\end{cases}$};

\node[c3card,anchor=north,text width=4.33cm,align=center] (pmf1) at (7.025,-11.6) {$\HLProbability{m_{Y_1}}$\par $\begin{array}{ll}\Money{-70}:&\FailureProb{1-p}\\\Money{90}:&\SuccessProb{p}\end{array}$\par \HLProbability{zero elsewhere}};

\node[inner sep=0pt,anchor=north,text width=4.33cm,align=center] (graph1) at (7.025,-13.899999999999999) {\CThreeDistributionPair{Y_1}{-70}{90}{1/2}{\Money{USD}}{\Money}};

\draw[c3flow] (cdf0.north) -- ++(0,.27) -| (cdf1.north);
\draw[c3flow] (pmf0.south) -- ++(0,-.30) -| (pmf1.south);
\draw[c3flow] (map1.east) -- (18.46,-3.3) -- (18.46,-7.8999999999999995) -| (head1.north);
\node[inner sep=1pt,anchor=north,text width=4.33cm,align=center] (head2) at (11.775,-8.2) {$\Money{Y_2}$ $\leftarrow\TransformationFunction{f_2}$};

\node[c3card,anchor=north,text width=4.33cm,align=center] (cdf2) at (11.775,-8.85) {$\HLCumulativeProbability{F_{Y_2}(t)}$\par $\begin{cases}\HLCumulativeProbability{0},&t<\Money{-30},\\\FailureProb{1-p},&\Money{-30}\leq t<\Money{190},\\\HLCumulativeProbability{1},&t\geq \Money{190}.\end{cases}$};

\node[c3card,anchor=north,text width=4.33cm,align=center] (pmf2) at (11.775,-11.6) {$\HLProbability{m_{Y_2}}$\par $\begin{array}{ll}\Money{-30}:&\FailureProb{1-p}\\\Money{190}:&\SuccessProb{p}\end{array}$\par \HLProbability{zero elsewhere}};

\node[inner sep=0pt,anchor=north,text width=4.33cm,align=center] (graph2) at (11.775,-13.899999999999999) {\CThreeDistributionPair{Y_2}{-30}{190}{1/2}{\Money{USD}}{\Money}};

\draw[c3flow] (cdf0.north) -- ++(0,.27) -| (cdf2.north);
\draw[c3flow] (pmf0.south) -- ++(0,-.30) -| (pmf2.south);
\draw[c3flow] (map2.east) -- (18.28,-5.1) -- (18.28,-7.779999999999999) -| (head2.north);
\node[inner sep=1pt,anchor=north,text width=4.33cm,align=center] (head3) at (16.525,-8.2) {$\Money{Y_3}$ $\leftarrow\TransformationFunction{f_3}$};

\node[c3card,anchor=north,text width=4.33cm,align=center] (cdf3) at (16.525,-8.85) {$\HLCumulativeProbability{F_{Y_3}(t)}$\par $\begin{cases}\HLCumulativeProbability{0},&t<\Money{-150},\\\FailureProb{1-p},&\Money{-150}\leq t<\Money{150},\\\HLCumulativeProbability{1},&t\geq \Money{150}.\end{cases}$};

\node[c3card,anchor=north,text width=4.33cm,align=center] (pmf3) at (16.525,-11.6) {$\HLProbability{m_{Y_3}}$\par $\begin{array}{ll}\Money{-150}:&\FailureProb{1-p}\\\Money{150}:&\SuccessProb{p}\end{array}$\par \HLProbability{zero elsewhere}};

\node[inner sep=0pt,anchor=north,text width=4.33cm,align=center] (graph3) at (16.525,-13.899999999999999) {\CThreeDistributionPair{Y_3}{-150}{150}{1/2}{\Money{USD}}{\Money}};

\draw[c3flow] (cdf0.north) -- ++(0,.27) -| (cdf3.north);
\draw[c3flow] (pmf0.south) -- ++(0,-.30) -| (pmf3.south);
\draw[c3flow] (map3.east) -- (18.1,-6.9) -- (18.1,-7.659999999999999) -| (head3.north);
\end{tikzpicture}\end{center}
\textbf{Numerical graph illustration:} $\SuccessProb{p_0}=\SuccessProb{\frac{1}{2}}$. All displayed distribution formulas remain symbolic in $\SuccessProb{p}$.
\par\textbf{CDF above, PMF below.} Filled stems translate to pale dashed jump shadows; thin arrows link equal physical lengths. Open and filled circles show excluded and included endpoints. Distribution arrows retain state probabilities while mapping values, rather than applying $\TransformationFunction{f_i}$ to CDF heights. See \hyperref[ex:intro-one]{Example 1}.\par
For $0<\SuccessProb{p}<1$, $\Indicator{S_X}=\{\Indicator{0},\Indicator{1}\}$; at $\SuccessProb{p}=0$ only $\Indicator{0}$ remains, and at $\SuccessProb{p}=1$ only $\Indicator{1}$ remains. Each monetary support retains $\Money{f_i(0)}$ at $\SuccessProb{p}=0$ and $\Money{f_i(1)}$ at $\SuccessProb{p}=1$.
\par\textbf{Loss events.} $\HLMathEvent{\{Y_1<0\}}=\HLMathEvent{\{X=0\}}=\HLMathEvent{E^c}$, probability $\FailureProb{1-p}$; $\HLMathEvent{\{Y_2<0\}}=\HLMathEvent{\{X=0\}}=\HLMathEvent{E^c}$, probability $\FailureProb{1-p}$; $\HLMathEvent{\{Y_3<0\}}=\HLMathEvent{\{X=0\}}=\HLMathEvent{E^c}$, probability $\FailureProb{1-p}$.
\endgroup
}
\newcommand{\CThreeExpectAgriculture}{%
\clearpage\noindent\textbf{Task 2 summary: parallel expectation analyses.}\par\medskip

\begingroup\small\tikzset{gten marker/.append style={inner xsep=.4mm,inner ysep=.2mm}}
\begin{center}\begin{tikzpicture}[x=1cm,y=1cm]
\node[c3card,anchor=north,text width=17cm,align=center] (inputs) at (9.5,0) {Shared probability inputs: $\SuccessProb{p}$ and $\FailureProb{1-p}$; hold each alternative\textquotesingle s \Money{payoffs} fixed.};

\node[c3card,anchor=north,text width=17cm,align=center] (formula) at (9.5,-1.2) {$\ExpectedValue{g_i}(\SuccessProb{p})=\ExpectedValue{\mathbb E}[\Money{Y_i}]=\SuccessProb{p}\Money{f_i(1)}+(\FailureProb{1-p})\Money{f_i(0)}$};

\draw[c3flow] (inputs) -- (formula);
\node[c3card,anchor=north,text width=5.633333333333333cm,align=center] (pay1) at (3.1666666666666665,-2.7) {\textbf{\Money{$Y_1$}: walking group}\par $\Money{f_1(0)}=\Money{-70},\quad\Money{f_1(1)}=\Money{90}\ \Money{\mathrm{USD}}$\par $\ExpectedValue{g_1}(\SuccessProb{p})=\Money{-70}+\Money{160}\SuccessProb{p}$};

\draw[c3flow] (formula.south) -- ++(0,-.25) -| (pay1.north);
\node[c3card,anchor=north,text width=5.633333333333333cm,align=center] (initial1) at (3.1666666666666665,-5.0) {\ExpectedValue{Initial expectation}\par $\ExpectedValue{g_1}\!\left(\SuccessProb{\frac{1}{2}}\right)=\ExpectedValue{10\ \mathrm{USD}}$\par \ExpectedValue{expected net gain} per stated period};

\node[c3card,anchor=north,text width=5.633333333333333cm,align=center] (root1) at (3.1666666666666665,-7.15) {\SuccessProb{Exact break-even equation}\par $\Money{-70}+\Money{160}\SuccessProb{p}=\ExpectedValue{0}$\par $\SuccessProb{p_1^*}=\SuccessProb{\frac{7}{16}}$};

\node[c3card,anchor=north,text width=5.633333333333333cm,align=center] (sign1) at (3.1666666666666665,-9.45) {$\ExpectedValue{g_1}>\ExpectedValue{0}:\ \SuccessProb{p}\in \left(\SuccessProb{\frac{7}{16}},\SuccessProb{1}\right]$\par $\ExpectedValue{g_1}=\ExpectedValue{0}:\ \SuccessProb{p}=\SuccessProb{\frac{7}{16}}$\par $\ExpectedValue{g_1}<\ExpectedValue{0}:\ \SuccessProb{p}\in \left[\SuccessProb{0},\SuccessProb{\frac{7}{16}}\right)$};

\node[inner sep=0pt,anchor=north,text width=5.633333333333333cm,align=center] (eg1) at (3.1666666666666665,-12.0) {\CThreeExpectationGraph{1}{-70}{90}{7/16}{1/2}};

\draw[c3flow] (pay1.south) -- (initial1.north);
\draw[c3flow] (initial1.south) -- (root1.north);
\draw[c3flow] (root1.south) -- (sign1.north);
\draw[c3flow] (sign1.south) -- (eg1.north);
\node[c3card,anchor=north,text width=5.633333333333333cm,align=center] (pay2) at (9.5,-2.7) {\textbf{\Money{$Y_2$}: dance sessions}\par $\Money{f_2(0)}=\Money{-30},\quad\Money{f_2(1)}=\Money{190}\ \Money{\mathrm{USD}}$\par $\ExpectedValue{g_2}(\SuccessProb{p})=\Money{-30}+\Money{220}\SuccessProb{p}$};

\draw[c3flow] (formula.south) -- ++(0,-.25) -| (pay2.north);
\node[c3card,anchor=north,text width=5.633333333333333cm,align=center] (initial2) at (9.5,-5.0) {\ExpectedValue{Initial expectation}\par $\ExpectedValue{g_2}\!\left(\SuccessProb{\frac{1}{2}}\right)=\ExpectedValue{80\ \mathrm{USD}}$\par \ExpectedValue{expected net gain} per stated period};

\node[c3card,anchor=north,text width=5.633333333333333cm,align=center] (root2) at (9.5,-7.15) {\SuccessProb{Exact break-even equation}\par $\Money{-30}+\Money{220}\SuccessProb{p}=\ExpectedValue{0}$\par $\SuccessProb{p_2^*}=\SuccessProb{\frac{3}{22}}$};

\node[c3card,anchor=north,text width=5.633333333333333cm,align=center] (sign2) at (9.5,-9.45) {$\ExpectedValue{g_2}>\ExpectedValue{0}:\ \SuccessProb{p}\in \left(\SuccessProb{\frac{3}{22}},\SuccessProb{1}\right]$\par $\ExpectedValue{g_2}=\ExpectedValue{0}:\ \SuccessProb{p}=\SuccessProb{\frac{3}{22}}$\par $\ExpectedValue{g_2}<\ExpectedValue{0}:\ \SuccessProb{p}\in \left[\SuccessProb{0},\SuccessProb{\frac{3}{22}}\right)$};

\node[inner sep=0pt,anchor=north,text width=5.633333333333333cm,align=center] (eg2) at (9.5,-12.0) {\CThreeExpectationGraph{2}{-30}{190}{3/22}{1/2}};

\draw[c3flow] (pay2.south) -- (initial2.north);
\draw[c3flow] (initial2.south) -- (root2.north);
\draw[c3flow] (root2.south) -- (sign2.north);
\draw[c3flow] (sign2.south) -- (eg2.north);
\node[c3card,anchor=north,text width=5.633333333333333cm,align=center] (pay3) at (15.833333333333332,-2.7) {\textbf{\Money{$Y_3$}: circuit sessions}\par $\Money{f_3(0)}=\Money{-150},\quad\Money{f_3(1)}=\Money{150}\ \Money{\mathrm{USD}}$\par $\ExpectedValue{g_3}(\SuccessProb{p})=\Money{-150}+\Money{300}\SuccessProb{p}$};

\draw[c3flow] (formula.south) -- ++(0,-.25) -| (pay3.north);
\node[c3card,anchor=north,text width=5.633333333333333cm,align=center] (initial3) at (15.833333333333332,-5.0) {\ExpectedValue{Initial expectation}\par $\ExpectedValue{g_3}\!\left(\SuccessProb{\frac{1}{2}}\right)=\ExpectedValue{0\ \mathrm{USD}}$\par \ExpectedValue{break-even} per stated period};

\node[c3card,anchor=north,text width=5.633333333333333cm,align=center] (root3) at (15.833333333333332,-7.15) {\SuccessProb{Exact break-even equation}\par $\Money{-150}+\Money{300}\SuccessProb{p}=\ExpectedValue{0}$\par $\SuccessProb{p_3^*}=\SuccessProb{\frac{1}{2}}$};

\node[c3card,anchor=north,text width=5.633333333333333cm,align=center] (sign3) at (15.833333333333332,-9.45) {$\ExpectedValue{g_3}>\ExpectedValue{0}:\ \SuccessProb{p}\in \left(\SuccessProb{\frac{1}{2}},\SuccessProb{1}\right]$\par $\ExpectedValue{g_3}=\ExpectedValue{0}:\ \SuccessProb{p}=\SuccessProb{\frac{1}{2}}$\par $\ExpectedValue{g_3}<\ExpectedValue{0}:\ \SuccessProb{p}\in \left[\SuccessProb{0},\SuccessProb{\frac{1}{2}}\right)$};

\node[inner sep=0pt,anchor=north,text width=5.633333333333333cm,align=center] (eg3) at (15.833333333333332,-12.0) {\CThreeExpectationGraph{3}{-150}{150}{1/2}{1/2}};

\draw[c3flow] (pay3.south) -- (initial3.north);
\draw[c3flow] (initial3.south) -- (root3.north);
\draw[c3flow] (root3.south) -- (sign3.north);
\draw[c3flow] (sign3.south) -- (eg3.north);
\end{tikzpicture}\end{center}
The dashed horizontal line is \ExpectedValue{zero}; filled circles show endpoint payoffs and the exact root. The square marks the \ExpectedValue{initial expectation}. The $+$ and $-$ labels connect the graph to its strict sign intervals.
\par \ExpectedValue{Positive expectation} means \ExpectedValue{expected net gain}; \ExpectedValue{negative expectation} means \ExpectedValue{expected net loss}; \ExpectedValue{zero expectation} means break-even. These are long-run averages per stated period, not guarantees of realized profit. Each branch uses its own payoffs and root.
\endgroup
}
\begin{document}
\maketitle
\setcounter{tocdepth}{1}
\setcounter{secnumdepth}{2}
\phantomsection\label{toc}
\tableofcontents
\section*{Evaluated criterion}
\addcontentsline{toc}{section}{Evaluated criterion}
\textbf{C3:} Models binary economic outcomes with Bernoulli random variables, interpreting success and failure in contexts such as default, gain, loss, or claim occurrence.

\textbf{How to practise.} Attempt each C3 task before opening its linked worked solution.
Use the return arrows for the contents, the problem, and its problem list.
All data are hypothetical teaching models. Keep exact expressions where useful;
use exact break-even probabilities and decimals only as approximations. Monetary amounts are in USD.
Each monetary variable has one strictly positive and one strictly negative payoff.
Analyse the sign of each \ExpectedValue{expected value} separately, using its own zero threshold.
Each scenario compares hypothetical alternatives; choose one rather than running them together.
All monetary payoffs are net gains or losses after the relevant costs for the stated period.
All variables in a problem are functions of the same $X$; separate analyses do
not imply statistical independence.


\clearpage\section*{Highlight legend}
\phantomsection\label{tab:highlight-legend}
\PracticeHighlightLegend{G10-T2-C3-coordinated}

The two probability treatments distinguish an \HLMathEvent{event} from its \HLMathEvent{complement} within
the same model. Labels and units preserve meaning without color.
See the \hyperref[tab:notation-reference]{detailed notation reference}.

\clearpage\section{Introduction}\label{sec:t2-c3-introduction}
\subsection{Introductory example 1: campus art-print stall}\label{ex:intro-one}
\CThreePlain{\CThreeStatementIntroOne}
\CThreePlain{\CThreeIntroTaskOne}
\CThreePlain{\CThreeIntroTaskTwo}

\paragraph{Worked explanation --- highlighted statement.}
\CThreeStatementIntroOne
\CThreeIntroTaskOne
\paragraph{Outcomes, events and monetary transformations.}
An \HLMathEvent{outcome} $\HLMathEvent{\omega}$ is one possible outcome of observing the stall over the weekend.
The \HLMathEvent{sample space} $\HLMathEvent{\Omega}$ contains all such outcomes, so $\HLMathEvent{\omega\in\Omega}$.
The \HLMathEvent{event} $\HLMathEvent{E}\subseteq\HLMathEvent{\Omega}$ is the set of high-demand outcomes, and $\HLMathEvent{E^c}$ is its complement.
Define the \Indicator{indicator} by
\[
\Indicator{X(\omega)}=
\begin{cases}
\Indicator{1},&\text{if }\HLMathEvent{\omega\in E},\\
\Indicator{0},&\text{if }\HLMathEvent{\omega\in E^c}.
\end{cases}
\]
Now let $\Money{Y}=\TransformationFunction{f}(\Indicator{X})$.
This means $\Money{Y(\omega)}=\TransformationFunction{f}(\Indicator{X(\omega)})$ for every outcome:
\[
\Money{Y(\omega)}=
\begin{cases}
\Money{f(1)},&\text{if }\HLMathEvent{\omega\in E},\\
\Money{f(0)},&\text{if }\HLMathEvent{\omega\in E^c}.
\end{cases}
\]
If $\HLMathEvent{\omega\in E}$, then $\Indicator{X(\omega)}=\Indicator{1}$,
$\Money{Y_1(\omega)}=\Money{60\ \mathrm{USD}}$ and $\Money{Y_2(\omega)}=\Money{120\ \mathrm{USD}}$.
If $\HLMathEvent{\omega\in E^c}$, then $\Indicator{X(\omega)}=\Indicator{0}$,
$\Money{Y_1(\omega)}=\Money{-20\ \mathrm{USD}}$ and $\Money{Y_2(\omega)}=\Money{-40\ \mathrm{USD}}$.
The \Indicator{indicator} records which \HLMathEvent{event} occurred;
$\TransformationFunction{f_i}$ converts that indicator value into the \Money{monetary consequence} for alternative $i$.
All alternatives use the same observed $\HLMathEvent{\omega}$ and the same $\Indicator{X(\omega)}$.
Since the indicator is either \Indicator{0} or \Indicator{1}, the two state evaluations are equivalent to the affine rule
\[
\Money{Y_i}=\Money{f_i(0)}+[\Money{f_i(1)}-\Money{f_i(0)}]\Indicator{X},
\qquad
\Money{Y_i(\omega)}=\Money{f_i(0)}+[\Money{f_i(1)}-\Money{f_i(0)}]\Indicator{X(\omega)}.
\]
Write $\Indicator{X(\omega)}$ for $\HLMathEvent{\omega\in E}$: an event is a set of outcomes, not an individual input to the indicator.

Both monetary variables use this same $\Indicator{X}$:
\[
\begin{aligned}
\Money{Y_1}&=\Money{f_1(0)}+[\Money{f_1(1)}-\Money{f_1(0)}]\Indicator{X}=-20+80\Indicator{X}\quad\text{USD},\\
\Money{Y_2}&=\Money{f_2(0)}+[\Money{f_2(1)}-\Money{f_2(0)}]\Indicator{X}=-40+160\Indicator{X}\quad\text{USD}.
\end{aligned}
\]
Their supports are $\{-20,60\}$ and $\{-40,120\}$ USD; each puts mass 0.6
on its negative value and 0.4 on its positive value. Each CDF is 0 below its
negative value, 0.6 from that value up to (excluding) its positive value,
and 1 from its positive value onward.
For $i=1,2$, $\HLMathEvent{\{Y_i<0\}}=\HLMathEvent{\{X=0\}}=\HLMathEvent{E^c}$ has probability $1-p$, initially 0.6.
\CThreeModelIntroOne
\clearpage
\CThreeIntroTaskTwo
\paragraph{Initial expected values.}
\textbf{$Y_1$: initial expectation.}
Initially, $\ExpectedValue{\mathbb E}[\Money{Y_1}]=\Money{60}(\SuccessProb{\frac{2}{5}})+\Money{(-20)}(\FailureProb{\frac{3}{5}})=\ExpectedValue{12\,\text{USD}}$:
an \ExpectedValue{expected net gain} of \ExpectedValue{12 USD} per stated period.

\textbf{$Y_2$: initial expectation.}
Initially, $\ExpectedValue{\mathbb E}[\Money{Y_2}]=\Money{120}(\SuccessProb{\frac{2}{5}})+\Money{(-40)}(\FailureProb{\frac{3}{5}})=\ExpectedValue{24\,\text{USD}}$:
an \ExpectedValue{expected net gain} of \ExpectedValue{24 USD} per stated period.
\paragraph{Expectation formulas and sign intervals.}
\textbf{$Y_1$: its own break-even threshold.}

With both payoffs fixed and only $p\in\left[0,1\right]$ varying,
\[
\begin{aligned}
\ExpectedValue{\mathbb E}[\Money{Y_1}]&=\SuccessProb{p}(\Money{60})+(\FailureProb{1-p})(\Money{-20})=\Money{-20}+\Money{80}\SuccessProb{p}\ \text{USD},\\
\ExpectedValue{\mathbb E}[\Money{Y_1}]&=\ExpectedValue{0}\\
 &\iff \Money{-20}+\Money{80}\SuccessProb{p}=\ExpectedValue{0}\iff 80\SuccessProb{p}=20\iff \SuccessProb{p_1^*}=\frac{1}{4}.
\end{aligned}
\]
Solving $\ExpectedValue{\mathbb E}[\Money{Y_1}]>0$ gives $p\in\left(\frac{1}{4},1\right]$ (\ExpectedValue{expected net gain}).
Solving $\ExpectedValue{\mathbb E}[\Money{Y_1}]<0$ gives $p\in\left[0,\frac{1}{4}\right)$ (\ExpectedValue{expected net loss}).
At $p=\frac{1}{4}$ the \ExpectedValue{expectation} is \ExpectedValue{zero} (break-even).
\textbf{$Y_2$: its own break-even threshold.}

With both payoffs fixed and only $p\in\left[0,1\right]$ varying,
\[
\begin{aligned}
\ExpectedValue{\mathbb E}[\Money{Y_2}]&=\SuccessProb{p}(\Money{120})+(\FailureProb{1-p})(\Money{-40})=\Money{-40}+\Money{160}\SuccessProb{p}\ \text{USD},\\
\ExpectedValue{\mathbb E}[\Money{Y_2}]&=\ExpectedValue{0}\\
 &\iff \Money{-40}+\Money{160}\SuccessProb{p}=\ExpectedValue{0}\iff 160\SuccessProb{p}=40\iff \SuccessProb{p_2^*}=\frac{1}{4}.
\end{aligned}
\]
Solving $\ExpectedValue{\mathbb E}[\Money{Y_2}]>0$ gives $p\in\left(\frac{1}{4},1\right]$ (\ExpectedValue{expected net gain}).
Solving $\ExpectedValue{\mathbb E}[\Money{Y_2}]<0$ gives $p\in\left[0,\frac{1}{4}\right)$ (\ExpectedValue{expected net loss}).
At $p=\frac{1}{4}$ the \ExpectedValue{expectation} is \ExpectedValue{zero} (break-even).
\CThreeExpectIntroOne

\clearpage
\subsection{Introductory example 2: freelance livestream work}\label{ex:intro-two}
\CThreePlain{\CThreeStatementIntroTwo}
\CThreePlain{\CThreeIntroTaskOne}
\CThreePlain{\CThreeIntroTaskTwo}

\paragraph{Worked explanation --- highlighted statement.}
\CThreeStatementIntroTwo
\CThreeIntroTaskOne

For the same observed $\HLMathEvent{\omega\in\Omega}$, $\Indicator{X(\omega)}=\Indicator{1}$ represents \HLMathEvent{disruption}: then $\Money{Y_1(\omega)}=\Money{-20\ \mathrm{USD}}$ and $\Money{Y_2(\omega)}=\Money{-10\ \mathrm{USD}}$.
For $\HLMathEvent{\omega\in E^c}$, $\Indicator{X(\omega)}=\Indicator{0}$, giving $\Money{100\ \mathrm{USD}}$ and $\Money{60\ \mathrm{USD}}$.
Each $\Money{Y_i(\omega)}=\TransformationFunction{f_i}(\Indicator{X(\omega)})$ uses this same outcome and indicator; here the monetary rules are decreasing.

\[
\begin{aligned}
\Money{Y_1}&=\Money{f_1(0)}+[\Money{f_1(1)}-\Money{f_1(0)}]\Indicator{X}=100-120\Indicator{X}\quad\text{USD},\\
\Money{Y_2}&=\Money{f_2(0)}+[\Money{f_2(1)}-\Money{f_2(0)}]\Indicator{X}=60-70\Indicator{X}\quad\text{USD}.
\end{aligned}
\]
Their supports are $\{-20,100\}$ and $\{-10,60\}$ USD. Each PMF places
mass $p=0.1$ on the negative payoff and $1-p=0.9$ on the positive payoff.
Each CDF is 0 below its negative value, $p$ between its ordered values,
and 1 at or above its positive value. For each variable,
$\HLMathEvent{\{Y_i<0\}}=\HLMathEvent{\{X=1\}}=\HLMathEvent{E}$, so $P(Y_i<0)=p=0.1$ initially.
\CThreeModelIntroTwo
\clearpage
\CThreeIntroTaskTwo
\paragraph{Initial expected values.}
\textbf{$Y_1$: initial expectation.}
Initially, $\ExpectedValue{\mathbb E}[\Money{Y_1}]=\Money{-20}(\SuccessProb{\frac{1}{10}})+\Money{(100)}(\FailureProb{\frac{9}{10}})=\ExpectedValue{88\,\text{USD}}$:
an \ExpectedValue{expected net gain} of \ExpectedValue{88 USD} per stated period.

\textbf{$Y_2$: initial expectation.}
Initially, $\ExpectedValue{\mathbb E}[\Money{Y_2}]=\Money{-10}(\SuccessProb{\frac{1}{10}})+\Money{(60)}(\FailureProb{\frac{9}{10}})=\ExpectedValue{53\,\text{USD}}$:
an \ExpectedValue{expected net gain} of \ExpectedValue{53 USD} per stated period.
\paragraph{Expectation formulas and sign intervals.}
\textbf{$Y_1$: its own break-even threshold.}

With both payoffs fixed and only $p\in\left[0,1\right]$ varying,
\[
\begin{aligned}
\ExpectedValue{\mathbb E}[\Money{Y_1}]&=\SuccessProb{p}(\Money{-20})+(\FailureProb{1-p})(\Money{100})=\Money{100}-\Money{120}\SuccessProb{p}\ \text{USD},\\
\ExpectedValue{\mathbb E}[\Money{Y_1}]&=\ExpectedValue{0}\\
 &\iff \Money{100}-\Money{120}\SuccessProb{p}=\ExpectedValue{0}\iff -120\SuccessProb{p}=-100\iff \SuccessProb{p_1^*}=\frac{5}{6}.
\end{aligned}
\]
Solving $\ExpectedValue{\mathbb E}[\Money{Y_1}]>0$ gives $p\in\left[0,\frac{5}{6}\right)$ (\ExpectedValue{expected net gain}).
Solving $\ExpectedValue{\mathbb E}[\Money{Y_1}]<0$ gives $p\in\left(\frac{5}{6},1\right]$ (\ExpectedValue{expected net loss}).
At $p=\frac{5}{6}$ the \ExpectedValue{expectation} is \ExpectedValue{zero} (break-even).
\textbf{$Y_2$: its own break-even threshold.}

With both payoffs fixed and only $p\in\left[0,1\right]$ varying,
\[
\begin{aligned}
\ExpectedValue{\mathbb E}[\Money{Y_2}]&=\SuccessProb{p}(\Money{-10})+(\FailureProb{1-p})(\Money{60})=\Money{60}-\Money{70}\SuccessProb{p}\ \text{USD},\\
\ExpectedValue{\mathbb E}[\Money{Y_2}]&=\ExpectedValue{0}\\
 &\iff \Money{60}-\Money{70}\SuccessProb{p}=\ExpectedValue{0}\iff -70\SuccessProb{p}=-60\iff \SuccessProb{p_2^*}=\frac{6}{7}.
\end{aligned}
\]
Solving $\ExpectedValue{\mathbb E}[\Money{Y_2}]>0$ gives $p\in\left[0,\frac{6}{7}\right)$ (\ExpectedValue{expected net gain}).
Solving $\ExpectedValue{\mathbb E}[\Money{Y_2}]<0$ gives $p\in\left(\frac{6}{7},1\right]$ (\ExpectedValue{expected net loss}).
At $p=\frac{6}{7}$ the \ExpectedValue{expectation} is \ExpectedValue{zero} (break-even).
\CThreeExpectIntroTwo

\clearpage\subsection{Binary models and distributions}
One observation in a stated period gives two mutually exclusive and exhaustive
states. Write $X\sim\operatorname{Bernoulli}(p)$, $0\leq p\leq1$.
The support is a set of numerical values with \emph{positive} probability:
\[
 S_X=\{x\in\mathbb R:P(X=x)>0\}
 =\begin{cases}\{0,1\},&0<p<1,\\\{0\},&p=0,\\\{1\},&p=1.\end{cases}
\]
For a monetary consequence, $S_Y=\{y\in\mathbb R:P(Y=y)>0\}$.
An \HLMathEvent{event} such as $\HLMathEvent{\{Y_i<0\}}$ is a condition on the observation; it is not a support.
These discrete variables have probability mass functions (PMFs), not ordinary
probability density functions. Their cumulative distribution functions (CDFs)
add masses up to and including a threshold. No separate density is required.
All monetary variables in a problem share the same indicator $X$.
Analysing them separately does not make them statistically independent.
Strictly positive and negative describe monetary payoffs, not $\HLMathEvent{E}$ or the labels 0 and 1.

\clearpage\section{Notation Reference and Practice Procedure}\label{sec:t2-c3-notation}
\subsection*{Detailed notation reference}
\phantomsection\label{tab:notation-reference}
% Separate rows use notation contexts where available; introductory meanings stay local.
% Reuse table infrastructure and applicable fields without concept-union rendering.
\begingroup\small
\HighlightDetailedTable{%
$\HLMathEvent{\omega\in\Omega}$ & \HLMathEvent{Outcome and sample space} & An outcome $\HLMathEvent{\omega}$ is one possible observation; the sample space $\HLMathEvent{\Omega}$ is the set of all possible outcomes. An event satisfies $\HLMathEvent{E}\subseteq\HLMathEvent{\Omega}$. &
\textbullet\ \hyperref[ex:intro-one]{Example 1}: one possible weekend observation of the campus art-print stall.\par
\textbullet\ \hyperref[ex:intro-two]{Example 2}: one possible day of freelance livestream work.\\[4pt]
$\Indicator{X(\omega)}$ & \Indicator{Outcome evaluation} & Evaluate $\Indicator{X}$ at the individual outcome: $\Indicator{1}$ if $\HLMathEvent{\omega\in E}$, $\Indicator{0}$ if $\HLMathEvent{\omega\in E^c}$. The event is a set, not the argument of $X$. &
\textbullet\ \hyperref[ex:intro-one]{Example 1}: $\Indicator{1}$ on a high-demand weekend; $\Indicator{0}$ otherwise.\par
\textbullet\ \hyperref[ex:intro-two]{Example 2}: $\Indicator{1}$ on a disrupted day, even though disruption is financially undesirable.\\[4pt]
$\Money{Y_i(\omega)}$ & \Money{Monetary evaluation} & $\Money{Y_i(\omega)}=\TransformationFunction{f_i}(\Indicator{X(\omega)})$: compose the \Indicator{indicator} with alternative $i$'s \TransformationFunction{rule}. The evaluated amounts $\Money{f_i(0)}$ and $\Money{f_i(1)}$ are in \Money{USD}. &
\textbullet\ \hyperref[ex:intro-one]{Example 1}: at the same high-demand outcome, \Money{60 USD} for the small batch and \Money{120 USD} for the large batch.\par
\textbullet\ \hyperref[ex:intro-two]{Example 2}: at the same disrupted outcome, \Money{-20 USD} without backup and \Money{-10 USD} with backup.\\[4pt]
$\ExpectedValue{g_i}(\SuccessProb{p})$ & \ExpectedValue{Expectation} as a \ExpectedValue{function} of \SuccessProb{probability} & $\ExpectedValue{g_i}(\SuccessProb{p})=\ExpectedValue{\mathbb E}[\Money{Y_i}]$ for fixed state payoffs; varying $\SuccessProb{p}$ changes the distribution, not the payoffs. &
\textbullet\ \hyperref[ex:intro-one]{Example 1}: both expectation functions increase with $\SuccessProb{p}$.\par
\textbullet\ \hyperref[ex:intro-two]{Example 2}: both decrease because $\SuccessProb{p}$ measures disruption.\\[4pt]

$\HLMathEvent{E}$ & \HLMathEvent{Event} & A set of outcomes satisfying the stated condition in one observation period; distinct from the expectation operator $\ExpectedValue{\mathbb E}$. &
\textbullet\ \hyperref[ex:intro-one]{Example 1}: \HLMathEvent{high campus-market attendance} over one weekend.\par
\textbullet\ \hyperref[ex:intro-two]{Example 2}: \HLMathEvent{home-internet disruption} over one day; mathematical success can be financially undesirable.\\[4pt]
$\HLMathEvent{E^c}$ & \HLMathEvent{Complement} & Outcomes outside $\HLMathEvent{E}$; $\HLMathEvent{E}$ and $\HLMathEvent{E^c}$ are mutually exclusive and exhaustive. &
\textbullet\ \hyperref[ex:intro-one]{Example 1}: \HLMathEvent{campus-market attendance} \HLMathEvent{below its target} that weekend.\par
\textbullet\ \hyperref[ex:intro-two]{Example 2}: \HLMathEvent{no disruption} that day.\\[4pt]
$\Indicator{X}$ & \Indicator{Indicator} & $X=1$ on $\HLMathEvent{E}$ and $X=0$ on $\HLMathEvent{E^c}$; $X\sim\operatorname{Bernoulli}(p)$. Unitless, rather than a monetary payoff. All monetary variables in a problem use this same $X$; no independence is implied. &
\textbullet\ \hyperref[ex:intro-one]{Example 1}: $\Indicator{1}$ means \HLMathEvent{high demand}, $\Indicator{0}$ \HLMathEvent{low demand}.\par
\textbullet\ \hyperref[ex:intro-two]{Example 2}: $\Indicator{1}$ means \HLMathEvent{disruption}, $\Indicator{0}$ \HLMathEvent{no disruption}.\\[4pt]
$\SuccessProb{p}$ & \SuccessProb{Success probability} & \HighlightEffectiveField{g10-t2-c3-coordinated-10-p}{definition} Here $p=P(\HLMathEvent{E})$; allow $0\leq p\leq1$. &
\textbullet\ \hyperref[ex:intro-one]{Example 1}: $\SuccessProb{0.4}=\SuccessProb{2/5}$.\par
\textbullet\ \hyperref[ex:intro-two]{Example 2}: $\SuccessProb{0.1}=\SuccessProb{1/10}$.\\[4pt]
$\FailureProb{1-p}$ & \FailureProb{Complementary probability} & $1-p=P(\HLMathEvent{E^c})=P(X=0)$; the two masses total 1. Financial gain or loss is determined by the payoffs, not by the label ``failure''. &
\textbullet\ \hyperref[ex:intro-one]{Example 1}: $\FailureProb{0.6}=\FailureProb{3/5}$.\par
\textbullet\ \hyperref[ex:intro-two]{Example 2}: $\FailureProb{0.9}=\FailureProb{9/10}$.\\[4pt]
$\Money{Y}$ & \Money{Monetary consequence} & Generic monetary random variable: $Y=\TransformationFunction{f}(X)$, in USD. If $\Money{f(0)}=y_0$ and $\Money{f(1)}=y_1$, then $Y=y_0+(y_1-y_0)X$. Here the two payoffs have opposite strict signs. &
\textbullet\ \hyperref[ex:intro-one]{Example 1}: generic notation for a small- or large-batch net gain or loss.\par
\textbullet\ \hyperref[ex:intro-two]{Example 2}: generic notation for net gain or loss with or without mobile-data backup.\\[4pt]
$\Money{Y_i}$ & \Money{Indexed monetary consequence} & $Y_i=\TransformationFunction{f_i}(X)=\Money{f_i(0)}+[\Money{f_i(1)}-\Money{f_i(0)}]X$ in USD; $i$ identifies an alternative. Each variable has one strictly positive and one strictly negative payoff. &
\textbullet\ \hyperref[ex:intro-one]{Example 1}: $Y_1=-20+80X$, $Y_2=-40+160X$ USD (small and large print batches).\par
\textbullet\ \hyperref[ex:intro-two]{Example 2}: $Y_1=100-120X$, $Y_2=60-70X$ USD (home internet alone and mobile-data backup).\\[4pt]
$\TransformationFunction{f}$ & Transformation function & \HighlightEffectiveField{g10-t2-c3-notation-function}{definition} Here $\TransformationFunction{f}:\{0,1\}\to\mathbb R$ maps indicator states to money in USD. &
\textbullet\ \hyperref[ex:intro-one]{Example 1}: generic notation for either print batch's monetary rule.\par
\textbullet\ \hyperref[ex:intro-two]{Example 2}: generic notation for either internet choice's monetary rule.\\[4pt]
$\TransformationFunction{f_i}$ & Indexed transformation function & \HighlightEffectiveField{g10-t2-c3-notation-indexed-function}{definition} Here both state payoffs are fixed when only $p$ varies. &
\textbullet\ \hyperref[ex:intro-one]{Example 1}: $\TransformationFunction{f_1}(x)=-20+80x$, $\TransformationFunction{f_2}(x)=-40+160x$ USD for $x\in\{0,1\}$.\par
\textbullet\ \hyperref[ex:intro-two]{Example 2}: $\TransformationFunction{f_1}(x)=100-120x$, $\TransformationFunction{f_2}(x)=60-70x$ USD for $x\in\{0,1\}$.\\[4pt]
$\Money{f_i(0)}$ & \Money{State-0 payoff} & \HighlightEffectiveField{g10-t2-c3-notation-payoff-0}{definition} &
\textbullet\ \hyperref[ex:intro-one]{Example 1}: $\Money{f_1(0)}=\Money{-20}$, $\Money{f_2(0)}=\Money{-40}$ USD.\par
\textbullet\ \hyperref[ex:intro-two]{Example 2}: $\Money{f_1(0)}=\Money{100}$, $\Money{f_2(0)}=\Money{60}$ USD.\\[4pt]
$\Money{f_i(1)}$ & \Money{State-1 payoff} & \HighlightEffectiveField{g10-t2-c3-notation-payoff-1}{definition} &
\textbullet\ \hyperref[ex:intro-one]{Example 1}: $\Money{f_1(1)}=\Money{60}$, $\Money{f_2(1)}=\Money{120}$ USD.\par
\textbullet\ \hyperref[ex:intro-two]{Example 2}: $\Money{f_1(1)}=\Money{-20}$, $\Money{f_2(1)}=\Money{-10}$ USD.\\[4pt]
$\Indicator{S_X}$ & \Indicator{Indicator support} & $S_X=\{x\in\mathbb R:P(X=x)>0\}$. It is $\{0,1\}$ for $0<p<1$, $\{0\}$ at $p=0$, and $\{1\}$ at $p=1$; all values are unitless. &
\textbullet\ \hyperref[ex:intro-one]{Example 1} and \hyperref[ex:intro-two]{Example 2}: $S_X=\{0,1\}$ initially.\\[4pt]
$\Money{S_Y}$ & \Money{Monetary support} & $S_Y=\{y\in\mathbb R:P(Y=y)>0\}$. Only positive-mass monetary values belong; it is not an \HLMathEvent{event}. &
\textbullet\ \hyperref[ex:intro-one]{Example 1}: generic support notation; for the small print batch, $\{\Money{-20},\Money{60}\}$ USD initially.\\[4pt]
$\Money{S_{Y_i}}$ & \Money{Indexed monetary support} & $S_{Y_i}=\{y\in\mathbb R:P(Y_i=y)>0\}$. For $0<p<1$ it is $\{\Money{f_i(0)},\Money{f_i(1)}\}$; at $p=0$ only $\Money{f_i(0)}$ remains, and at $p=1$ only $\Money{f_i(1)}$ remains. &
\textbullet\ \hyperref[ex:intro-one]{Example 1}: $S_{Y_1}=\{\Money{-20},\Money{60}\}$, $S_{Y_2}=\{\Money{-40},\Money{120}\}$ USD initially.\par
\textbullet\ \hyperref[ex:intro-two]{Example 2}: $S_{Y_1}=\{\Money{-20},\Money{100}\}$, $S_{Y_2}=\{\Money{-10},\Money{60}\}$ USD initially.\\[4pt]
$\HLProbability{m_X(x)}$ & \HLProbability{Indicator PMF} & \HighlightEffectiveField{g10-t2-c3-notation-indicator-mass}{definition} &
\textbullet\ \hyperref[ex:intro-one]{Example 1}: masses 0.6 at 0 and 0.4 at 1.\par
\textbullet\ \hyperref[ex:intro-two]{Example 2}: masses 0.9 at 0 and 0.1 at 1.\\[4pt]
$\HLProbability{m_{Y_i}(y)}$ & \HLProbability{Monetary PMF} & \HighlightEffectiveField{g10-t2-c3-notation-transformed-mass}{definition} Here the distinct USD payoffs receive $1-p$ at $\Money{f_i(0)}$ and $p$ at $\Money{f_i(1)}$, with zero elsewhere. &
\textbullet\ \hyperref[ex:intro-one]{Example 1}: mass 0.6 on each negative payoff and 0.4 on each positive payoff.\par
\textbullet\ \hyperref[ex:intro-two]{Example 2}: $m_{Y_1}(-20)=0.1$, $m_{Y_2}(-10)=0.1$; each positive payoff has mass 0.9.\\[4pt]
$\HLCumulativeProbability{F_X(t)}$ & \HLCumulativeProbability{Indicator CDF} & \HighlightEffectiveField{g10-t2-c3-notation-indicator-cdf}{definition} &
\textbullet\ \hyperref[ex:intro-one]{Example 1}: middle plateau 0.6.\par
\textbullet\ \hyperref[ex:intro-two]{Example 2}: middle plateau 0.9.\\[4pt]
$\HLCumulativeProbability{F_{Y_i}(t)}$ & \HLCumulativeProbability{Monetary CDF} & \HighlightEffectiveField{g10-t2-c3-notation-transformed-cdf}{definition} Here the USD CDF is 0 below the smaller payoff, its mass up to the larger payoff, and 1 at or above the larger. &
\textbullet\ \hyperref[ex:intro-one]{Example 1}: middle plateau 0.6 for each variable.\par
\textbullet\ \hyperref[ex:intro-two]{Example 2}: middle plateau $p=0.1$; $F_{Y_1}(0)=0.1$.\\[4pt]
$\HLMathEvent{\{Y_i<0\}}$ & \HLMathEvent{Net-loss event} & The outcomes for which the monetary consequence is negative. Evaluate both payoffs to identify the corresponding indicator \HLMathEvent{event}; an \HLMathEvent{event} is not a support. &
\textbullet\ \hyperref[ex:intro-one]{Example 1}: $\HLMathEvent{\{Y_i<0\}}=\HLMathEvent{\{X=0\}}=\HLMathEvent{E^c}$, probability $1-p=0.6$ initially, for $i=1,2$.\par
\textbullet\ \hyperref[ex:intro-two]{Example 2}: $\HLMathEvent{\{Y_i<0\}}=\HLMathEvent{\{X=1\}}=\HLMathEvent{E}$, probability $p=0.1$ initially, for $i=1,2$.\\[4pt]
$\ExpectedValue{\mathbb E}$ & \ExpectedValue{Expectation operator} & \HighlightEffectiveField{g10-t2-c3-notation-expectation-operator}{definition} &
\textbullet\ \hyperref[ex:intro-one]{Example 1}: averages monetary payoffs over the weekend model.\par
\textbullet\ \hyperref[ex:intro-two]{Example 2}: averages monetary payoffs over the daily model; no realized profit is guaranteed.\\[4pt]
$\ExpectedValue{\mathbb E}[\Money{Y_i}]$ & \ExpectedValue{Expected monetary value} & \HighlightEffectiveField{g10-t2-c3-coordinated-10-expected}{definition} Here $\mathbb E[Y_i]=p \Money{f_i(1)}+(1-p)\Money{f_i(0)}$, in USD per stated period. &
\textbullet\ \hyperref[ex:intro-one]{Example 1}: $\mathbb E[Y_1]=-20+80p$, $\mathbb E[Y_2]=-40+160p$ USD; initially \ExpectedValue{12 USD} and \ExpectedValue{24 USD}, expected net gains per weekend.\par
\textbullet\ \hyperref[ex:intro-two]{Example 2}: $\mathbb E[Y_1]=100-120p$, $\mathbb E[Y_2]=60-70p$ USD; initially \ExpectedValue{88 USD} and \ExpectedValue{53 USD}, expected net gains per day.\\[4pt]
$\SuccessProb{p_i^*}$ & \SuccessProb{Individual break-even probability} & \HighlightEffectiveField{g10-t2-c3-coordinated-10-break}{definition} With fixed opposite-sign payoffs, $p_i^*=-\Money{f_i(0)}/[\Money{f_i(1)}-\Money{f_i(0)}]\in\left(0,1\right)$. Solve each strict sign inequality separately on $\left[0,1\right]$. &
\textbullet\ \hyperref[ex:intro-one]{Example 1}: $\SuccessProb{p_1^*}=\SuccessProb{1/4}$ and $\SuccessProb{p_2^*}=\SuccessProb{1/4}$; each expectation is positive above its root and negative below it.\par
\textbullet\ \hyperref[ex:intro-two]{Example 2}: $\SuccessProb{p_1^*}=\SuccessProb{5/6}$, $\SuccessProb{p_2^*}=\SuccessProb{6/7}$; each expectation is positive below its own root and negative above it. Equality means zero expected net profit.\\[4pt]
}\endgroup
Here $\TransformationFunction{f_i}$ transforms states into money; $m_{Y_i}$ assigns probabilities to
monetary values and $F_{Y_i}$ accumulates them. These are different functions.

\clearpage\needspace{20\baselineskip}\subsection{Reusable practice procedure}
\textbf{Task 1 --- Establish the binary model and monetary transformations.}
\begin{enumerate}
\item \textbf{Binary model and monetary transformations.} Present the shared observation, period, \HLMathEvent{event} $\HLMathEvent{E}$, \HLMathEvent{complement} $\HLMathEvent{E^c}$, \Indicator{indicator} $\Indicator{X}$, \Indicator{Bernoulli model} and initial $\SuccessProb{p}=\SuccessProb{P}(\HLMathEvent{E})$.\par Then give each \Money{monetary variable}'s \Money{state payoffs} in \Money{USD} and derive $\Money{Y_i}=\TransformationFunction{f_i}(\Indicator{X})$\allowbreak${}=\Money{f_i(0)}+[\Money{f_i(1)}-\Money{f_i(0)}]\Indicator{X}$.
\item \textbf{Supports, \HLProbability{PMFs}, and ordered \HLCumulativeProbability{CDFs}.} Give the \Indicator{indicator}'s \Indicator{set-builder support}, complete \HLProbability{PMF} and piecewise \HLCumulativeProbability{CDF}, including support cases at $\SuccessProb{p}=\SuccessProb{0},\SuccessProb{1}$. Then give the \Money{monetary supports}, complete \HLProbability{PMFs} (\HLProbability{zero elsewhere}) and numerically ordered piecewise \HLCumulativeProbability{CDFs}, with endpoint support cases.
\item \textbf{Initial probabilities and loss \HLMathEvent{events}.} State the initial \Indicator{indicator} probabilities. For each $\Money{Y_i}$, give both \HLProbability{initial masses}, identify $\HLMathEvent{\{Y_i<0\}}$ in terms of $\Indicator{X}$ and $\HLMathEvent{E}$ or $\HLMathEvent{E^c}$, calculate its probability and interpret it in context. All alternatives share $\Indicator{X}$; no independence is assumed.
\end{enumerate}
For $\SuccessProb{0}<\SuccessProb{p}<\SuccessProb{1}$, $\Indicator{S_X}=\{\Indicator{0},\Indicator{1}\}$; at $\SuccessProb{p}=\SuccessProb{0}$ only \Indicator{0} remains and at $\SuccessProb{p}=\SuccessProb{1}$ only \Indicator{1} remains.
\Money{Monetary supports} contain positive-mass values: $\Money{f_i(0)}$ only at $\SuccessProb{p}=\SuccessProb{0}$, $\Money{f_i(1)}$ only at $\SuccessProb{p}=\SuccessProb{1}$.
For ordered amounts $\Money{a}<\Money{b}$ and $\HLProbability{q}=\HLProbability{P}(\Money{Y_i}=\Money{a})$,
\[
\HLCumulativeProbability{F_{Y_i}(t)}=\begin{cases}\HLCumulativeProbability{0},&t<\Money{a},\\\HLProbability{q},&\Money{a}\leq t<\Money{b},\\\HLCumulativeProbability{1},&t\geq\Money{b}.\end{cases}
\]
Here $\HLProbability{q}=\FailureProb{1-p}$ if $\Money{a}=\Money{f_i(0)}$, and $\HLProbability{q}=\SuccessProb{p}$ if $\Money{a}=\Money{f_i(1)}$.
Check total mass \HLProbability{1}, \HLCumulativeProbability{right continuity}, units and opposite strict payoff signs.

\CThreeModelIntroOne
\clearpage
\textbf{Task 2 --- Individual \ExpectedValue{expected values}.}
\begin{enumerate}
\item \textbf{Initial \ExpectedValue{expected values}.} Calculate and interpret every \ExpectedValue{monetary expectation} at the supplied probability before sensitivity analysis.
\item \textbf{\ExpectedValue{Expectation formulas} and sign intervals.} Hold \Money{payoffs} fixed and vary only $\SuccessProb{p}\in\left[\SuccessProb{0},\SuccessProb{1}\right]$. For each $\Money{Y_i}$, derive $\ExpectedValue{\mathbb E}[\Money{Y_i}]=\SuccessProb{p} \Money{f_i(1)}+(\FailureProb{1-p})\Money{f_i(0)}$, solve its exact \SuccessProb{break-even equation}, and solve both strict sign inequalities with correct directions and endpoints. Explain \ExpectedValue{positive, negative and zero expectation}.
\item \textbf{Graphs.} Give a separate labelled graph for each variable, showing its \ExpectedValue{expectation line}, zero line, endpoint \Money{payoffs} and exact \SuccessProb{break-even probability}. Finish with endpoint and root checks, and distinguish \ExpectedValue{expected profit} from \ExpectedValue{realized profit}.
\end{enumerate}
Use $\SuccessProb{p_i^*}=-\Money{f_i(0)}/[\Money{f_i(1)}-\Money{f_i(0)}]$. For an increasing \ExpectedValue{expectation},
\ExpectedValue{gain} occurs on $\left(\SuccessProb{p_i^*},\SuccessProb{1}\right]$ and \ExpectedValue{loss} on $\left[\SuccessProb{0},\SuccessProb{p_i^*}\right)$; for a decreasing \ExpectedValue{expectation},
\ExpectedValue{gain} occurs on $\left[\SuccessProb{0},\SuccessProb{p_i^*}\right)$ and \ExpectedValue{loss} on $\left(\SuccessProb{p_i^*},\SuccessProb{1}\right]$. Dividing by a negative slope reverses inequalities.
Equality is \ExpectedValue{break-even}. Group corresponding steps across alternatives before moving to the next subtask.
Each analysis uses its own \Money{payoffs} and threshold; all alternatives share $\Indicator{X}$.

\CThreeExpectIntroOne
\enlargethispage{\baselineskip}
\secInicio[][sec:t2-c3-introduction]
\clearpage
\section{General Worked Problems}\label{sec:t2-c3-general}
Each problem applies C3 through one Bernoulli model and analyses each monetary \ExpectedValue{expectation} against zero separately.
\SectionProblemsTableOfContents{\SectionProblemLink{prob:c3-inventory}{Campus sticker sales for one weekend}\\
\SectionProblemLink{prob:c3-coffee}{Campus fair or online CV workshop}\\
\SectionProblemLink{prob:c3-energy}{Freelance digital projects for one month}\\
\SectionProblemLink{prob:c3-manufacturing}{Graduation photography for one quarter}}
\subsection{Campus sticker sales for one weekend}\label{prob:c3-inventory}
\CThreePlain{\CThreeStatementInventory}
\QuestionsTasks[sec:t2-c3-general]{
{3}/{sol:c3-inventory-model}/{\CThreePlain{\CThreeTaskOne}
},
{3}/{sol:c3-inventory-decision}/{\CThreePlain{\CThreeTaskTwo}
}
}
\needspace{16\baselineskip}\subsubsection*{Task 1 --- Worked solution}\label{sol:c3-inventory-model}
\paragraph{Highlighted problem statement.}
\CThreeStatementInventory
\CThreeTaskOne

\needspace{8\baselineskip}\paragraph{1. Binary model and monetary transformations.}
The observation is a student's campus sticker sales over one weekend. Let $\HLMathEvent{E}$ be \HLMathEvent{high demand}; $\HLMathEvent{E^c}$ is \HLMathEvent{low demand}. Exactly one state occurs. Define $\Indicator{X}=\Indicator{1}$ on $\HLMathEvent{E}$ and $\Indicator{X}=\Indicator{0}$ on $\HLMathEvent{E^c}$. This unitless indicator has $\Indicator{X}\sim\operatorname{Bernoulli}(\SuccessProb{p})$, initially $\SuccessProb{p}=\SuccessProb{\frac{2}{5}}$.
\PracticeColumnPair{
\textbf{$Y_1$: basic sticker pack.}
The \Money{monetary consequence} is $\Money{Y_1}=\TransformationFunction{f_1}(\Indicator{X})$ in USD.
\[
\begin{gathered}
f_1(0)=\Money{-40},\quad f_1(1)=\Money{80},\\
\Money{Y_1}=\Money{f_1(0)}+[\Money{f_1(1)}-\Money{f_1(0)}]\Indicator{X}\\
=\Money{-40}+[\Money{80}-(\Money{-40})]\Indicator{X}\\
=\Money{-40}+\Money{120}\Indicator{X}\quad\text{USD}.
\end{gathered}
\]
One payoff is a loss and the other a gain.
}{
\textbf{$Y_2$: premium sticker pack.}
The \Money{monetary consequence} is $\Money{Y_2}=\TransformationFunction{f_2}(\Indicator{X})$ in USD.
\[
\begin{gathered}
f_2(0)=\Money{-20},\quad f_2(1)=\Money{180},\\
\Money{Y_2}=\Money{f_2(0)}+[\Money{f_2(1)}-\Money{f_2(0)}]\Indicator{X}\\
=\Money{-20}+[\Money{180}-(\Money{-20})]\Indicator{X}\\
=\Money{-20}+\Money{200}\Indicator{X}\quad\text{USD}.
\end{gathered}
\]
One payoff is a loss and the other a gain.
}
\needspace{8\baselineskip}\paragraph{2. Supports, PMFs, and ordered CDFs.}

\[S_X=\{x\in\mathbb R:P(X=x)>0\}=\{0,1\}\quad(0<p<1).\]
\[m_X(x)=\begin{cases}\FailureProb{1-p},&x=0,\\\SuccessProb{p},&x=1,\\0,&\text{otherwise},\end{cases}\qquad
F_X(t)=\begin{cases}0,&t<0,\\\FailureProb{1-p},&0\leq t<1,\\1,&t\geq1.\end{cases}\]
At $p=0$, $S_X=\{0\}$; at $p=1$, $S_X=\{1\}$. The formulas remain valid at both endpoints.
\PracticeColumnPair{
\textbf{$Y_1$: basic sticker pack.}
\[
\begin{aligned}
S_{Y_1}&=\{y\in\mathbb R:P(Y_1=y)>0\}\\
&=\{\Money{-40},\Money{80}\}\quad(0<p<1).
\end{aligned}
\]
At $p=0$, $S_{Y_1}=\{\Money{-40}\}$; at $p=1$, $S_{Y_1}=\{\Money{80}\}$ (USD).
}{
\textbf{$Y_2$: premium sticker pack.}
\[
\begin{aligned}
S_{Y_2}&=\{y\in\mathbb R:P(Y_2=y)>0\}\\
&=\{\Money{-20},\Money{180}\}\quad(0<p<1).
\end{aligned}
\]
At $p=0$, $S_{Y_2}=\{\Money{-20}\}$; at $p=1$, $S_{Y_2}=\{\Money{180}\}$ (USD).
}
\PracticeColumnPair{
\textbf{$Y_1$: basic sticker pack.}
\[m_{Y_1}(y)=\begin{cases}\FailureProb{1-p},&y=\Money{-40},\\\SuccessProb{p},&y=\Money{80},\\0,&\text{otherwise}.\end{cases}\]
}{
\textbf{$Y_2$: premium sticker pack.}
\[m_{Y_2}(y)=\begin{cases}\FailureProb{1-p},&y=\Money{-20},\\\SuccessProb{p},&y=\Money{180},\\0,&\text{otherwise}.\end{cases}\]
}
\PracticeColumnPair{
\textbf{$Y_1$: basic sticker pack.}
\[F_{Y_1}(t)=\begin{cases}0,&t<-40,\\\FailureProb{1-p},&-40\leq t<80,\\1,&t\geq 80.\end{cases}\]
}{
\textbf{$Y_2$: premium sticker pack.}
\[F_{Y_2}(t)=\begin{cases}0,&t<-20,\\\FailureProb{1-p},&-20\leq t<180,\\1,&t\geq 180.\end{cases}\]
}
Each PMF totals $(1-p)+p=1$. CDF thresholds are numerically ordered,
right-continuous and include each jump at its threshold. At $p=0,1$, zero-mass jumps disappear.
Monetary supports are in USD; indicator values are unitless.
\needspace{8\baselineskip}\paragraph{3. Initial probabilities and loss \protect\HLMathEvent{events}.}
Initially $P(X=0)=\FailureProb{\frac{3}{5}}$ and $P(X=1)=\SuccessProb{\frac{2}{5}}$.
All monetary variables are functions of the same $X$; no independence is asserted.
\PracticeColumnPair{
\textbf{$Y_1$: basic sticker pack.}
\[
\begin{gathered}
m_{Y_1}(\Money{-40})=\FailureProb{\frac{3}{5}},\quad m_{Y_1}(\Money{80})=\SuccessProb{\frac{2}{5}},\\
\HLMathEvent{\{Y_1<0\}}=\HLMathEvent{\{X=0\}}=\HLMathEvent{E^c},\\
P(Y_1<0)=\FailureProb{1-p}= \frac{3}{5}\quad\text{initially}.
\end{gathered}
\]
For basic sticker pack, loss occurs on $\HLMathEvent{E^c}$: its initial probability per stated period is $\frac{3}{5}$.
}{
\textbf{$Y_2$: premium sticker pack.}
\[
\begin{gathered}
m_{Y_2}(\Money{-20})=\FailureProb{\frac{3}{5}},\quad m_{Y_2}(\Money{180})=\SuccessProb{\frac{2}{5}},\\
\HLMathEvent{\{Y_2<0\}}=\HLMathEvent{\{X=0\}}=\HLMathEvent{E^c},\\
P(Y_2<0)=\FailureProb{1-p}= \frac{3}{5}\quad\text{initially}.
\end{gathered}
\]
For premium sticker pack, loss occurs on $\HLMathEvent{E^c}$: its initial probability per stated period is $\frac{3}{5}$.
}
\CThreeModelInventory
\secInicio[prob:c3-inventory][sec:t2-c3-general]
\clearpage\subsubsection*{Task 2 --- Worked solution}\label{sol:c3-inventory-decision}
\CThreeTaskTwo

Keep the \HLMathEvent{event}, its \HLMathEvent{complement} and the indicator from Task 1.
\needspace{8\baselineskip}\paragraph{1. Initial expected values.}
Use the supplied $\SuccessProb{p}=\SuccessProb{\frac{2}{5}}$ for every alternative.
\PracticeColumnPair{
\textbf{$Y_1$: basic sticker pack.}
\[
\begin{aligned}
\ExpectedValue{\mathbb E}[\Money{Y_1}]&=\Money{80}\left(\SuccessProb{\frac{2}{5}}\right)\\
&\quad +(\Money{-40})\left(\FailureProb{\frac{3}{5}}\right)\\
&=\ExpectedValue{8\ \text{USD}}.
\end{aligned}
\]
This is \ExpectedValue{expected net gain} of $\ExpectedValue{8\ \text{USD}}$ per stated period as a long-run average.
}{
\textbf{$Y_2$: premium sticker pack.}
\[
\begin{aligned}
\ExpectedValue{\mathbb E}[\Money{Y_2}]&=\Money{180}\left(\SuccessProb{\frac{2}{5}}\right)\\
&\quad +(\Money{-20})\left(\FailureProb{\frac{3}{5}}\right)\\
&=\ExpectedValue{60\ \text{USD}}.
\end{aligned}
\]
This is \ExpectedValue{expected net gain} of $\ExpectedValue{60\ \text{USD}}$ per stated period as a long-run average.
}
\needspace{8\baselineskip}\paragraph{2. Expectation formulas and sign intervals.}
Hold all state payoffs fixed; vary only $p=P(\HLMathEvent{E})\in\left[0,1\right]$.
\PracticeColumnPair{
\textbf{$Y_1$: basic sticker pack.}
\[
\begin{aligned}
\ExpectedValue{\mathbb E}[\Money{Y_1}]&=\SuccessProb{p}\Money{f_1(1)}+(\FailureProb{1-p})\Money{f_1(0)}\\
&=\SuccessProb{p}(\Money{80})+(\FailureProb{1-p})(\Money{-40})\\
&=\Money{-40}+\Money{120}\SuccessProb{p}\quad\text{USD}.
\end{aligned}
\]
\[
\begin{aligned}
\ExpectedValue{\mathbb E}[\Money{Y_1}]=0&\iff -40+120p=0\\
&\iff 120p=40\\
&\iff p=p_1^*=\frac{1}{3}.
\end{aligned}
\]
}{
\textbf{$Y_2$: premium sticker pack.}
\[
\begin{aligned}
\ExpectedValue{\mathbb E}[\Money{Y_2}]&=\SuccessProb{p}\Money{f_2(1)}+(\FailureProb{1-p})\Money{f_2(0)}\\
&=\SuccessProb{p}(\Money{180})+(\FailureProb{1-p})(\Money{-20})\\
&=\Money{-20}+\Money{200}\SuccessProb{p}\quad\text{USD}.
\end{aligned}
\]
\[
\begin{aligned}
\ExpectedValue{\mathbb E}[\Money{Y_2}]=0&\iff -20+200p=0\\
&\iff 200p=20\\
&\iff p=p_2^*=\frac{1}{10}.
\end{aligned}
\]
}
\PracticeColumnPair{
\textbf{$Y_1$: basic sticker pack.}
\[
\begin{aligned}
\ExpectedValue{\mathbb E}[\Money{Y_1}]>0&\iff 120p>40\\
&\iff p>\frac{1}{3}\\
&\iff p\in \left(\frac{1}{3},1\right],\\
\ExpectedValue{\mathbb E}[\Money{Y_1}]<0&\iff 120p<40\\
&\iff p<\frac{1}{3}\\
&\iff p\in \left[0,\frac{1}{3}\right).
\end{aligned}
\]
At $\SuccessProb{p}=\frac{1}{3}$, $\ExpectedValue{\mathbb E}[\Money{Y_1}]$ is \ExpectedValue{zero}: break-even.
}{
\textbf{$Y_2$: premium sticker pack.}
\[
\begin{aligned}
\ExpectedValue{\mathbb E}[\Money{Y_2}]>0&\iff 200p>20\\
&\iff p>\frac{1}{10}\\
&\iff p\in \left(\frac{1}{10},1\right],\\
\ExpectedValue{\mathbb E}[\Money{Y_2}]<0&\iff 200p<20\\
&\iff p<\frac{1}{10}\\
&\iff p\in \left[0,\frac{1}{10}\right).
\end{aligned}
\]
At $\SuccessProb{p}=\frac{1}{10}$, $\ExpectedValue{\mathbb E}[\Money{Y_2}]$ is \ExpectedValue{zero}: break-even.
}
Positive \ExpectedValue{expectation} means \ExpectedValue{expected net gain}; negative
\ExpectedValue{expectation} means \ExpectedValue{expected net loss}; zero is break-even.
\needspace{8\baselineskip}\paragraph{3. Graphs.}
The line joins $(0,-40)$ and $(1,80)$; its zero is $p_1^*=\frac{1}{3}$.
The line joins $(0,-20)$ and $(1,180)$; its zero is $p_2^*=\frac{1}{10}$.

\textbf{Endpoint and root checks.} $Y_1$: $\mathbb E[Y_1](0)=-40$ and $\mathbb E[Y_1](1)=80$ USD; at $p_1^*=\frac{1}{3}$, $-40+(120)(\frac{1}{3})=0$.

$Y_2$: $\mathbb E[Y_2](0)=-20$ and $\mathbb E[Y_2](1)=180$ USD; at $p_2^*=\frac{1}{10}$, $-20+(200)(\frac{1}{10})=0$.

Each observation realizes one state payoff. An \ExpectedValue{expected profit} is a long-run average,
not a guarantee of realized profit. At break-even, zero expectation coexists with nonzero realized payoffs.
\CThreeExpectInventory
\secInicio[prob:c3-inventory][sec:t2-c3-general]
\clearpage
\subsection{Campus fair or online CV workshop}\label{prob:c3-coffee}
\CThreePlain{\CThreeStatementCoffee}
\QuestionsTasks[sec:t2-c3-general]{
{3}/{sol:c3-coffee-model}/{\CThreePlain{\CThreeTaskOne}
},
{3}/{sol:c3-coffee-decision}/{\CThreePlain{\CThreeTaskTwo}
}
}
\needspace{16\baselineskip}\subsubsection*{Task 1 --- Worked solution}\label{sol:c3-coffee-model}
\paragraph{Highlighted problem statement.}
\CThreeStatementCoffee
\CThreeTaskOne

\needspace{8\baselineskip}\paragraph{1. Binary model and monetary transformations.}
The observation is a graduate's paid sessions over one day. Let $\HLMathEvent{E}$ be \HLMathEvent{the fair going ahead}; $\HLMathEvent{E^c}$ is \HLMathEvent{its cancellation}. Exactly one state occurs. Define $\Indicator{X}=\Indicator{1}$ on $\HLMathEvent{E}$ and $\Indicator{X}=\Indicator{0}$ on $\HLMathEvent{E^c}$. This unitless indicator has $\Indicator{X}\sim\operatorname{Bernoulli}(\SuccessProb{p})$, initially $\SuccessProb{p}=\SuccessProb{\frac{1}{2}}$.
\PracticeColumnPair{
\textbf{$Y_1$: portrait sessions.}
The \Money{monetary consequence} is $\Money{Y_1}=\TransformationFunction{f_1}(\Indicator{X})$ in USD.
\[
\begin{gathered}
f_1(0)=\Money{-40},\quad f_1(1)=\Money{100},\\
\Money{Y_1}=\Money{f_1(0)}+[\Money{f_1(1)}-\Money{f_1(0)}]\Indicator{X}\\
=\Money{-40}+[\Money{100}-(\Money{-40})]\Indicator{X}\\
=\Money{-40}+\Money{140}\Indicator{X}\quad\text{USD}.
\end{gathered}
\]
One payoff is a loss and the other a gain.
}{
\textbf{$Y_2$: online CV workshop.}
The \Money{monetary consequence} is $\Money{Y_2}=\TransformationFunction{f_2}(\Indicator{X})$ in USD.
\[
\begin{gathered}
f_2(0)=\Money{100},\quad f_2(1)=\Money{-20},\\
\Money{Y_2}=\Money{f_2(0)}+[\Money{f_2(1)}-\Money{f_2(0)}]\Indicator{X}\\
=\Money{100}+[\Money{-20}-(\Money{100})]\Indicator{X}\\
=\Money{100}-\Money{120}\Indicator{X}\quad\text{USD}.
\end{gathered}
\]
One payoff is a loss and the other a gain.
}
\needspace{8\baselineskip}\paragraph{2. Supports, PMFs, and ordered CDFs.}

\[S_X=\{x\in\mathbb R:P(X=x)>0\}=\{0,1\}\quad(0<p<1).\]
\[m_X(x)=\begin{cases}\FailureProb{1-p},&x=0,\\\SuccessProb{p},&x=1,\\0,&\text{otherwise},\end{cases}\qquad
F_X(t)=\begin{cases}0,&t<0,\\\FailureProb{1-p},&0\leq t<1,\\1,&t\geq1.\end{cases}\]
At $p=0$, $S_X=\{0\}$; at $p=1$, $S_X=\{1\}$. The formulas remain valid at both endpoints.
\PracticeColumnPair{
\textbf{$Y_1$: portrait sessions.}
\[
\begin{aligned}
S_{Y_1}&=\{y\in\mathbb R:P(Y_1=y)>0\}\\
&=\{\Money{-40},\Money{100}\}\quad(0<p<1).
\end{aligned}
\]
At $p=0$, $S_{Y_1}=\{\Money{-40}\}$; at $p=1$, $S_{Y_1}=\{\Money{100}\}$ (USD).
}{
\textbf{$Y_2$: online CV workshop.}
\[
\begin{aligned}
S_{Y_2}&=\{y\in\mathbb R:P(Y_2=y)>0\}\\
&=\{\Money{-20},\Money{100}\}\quad(0<p<1).
\end{aligned}
\]
At $p=0$, $S_{Y_2}=\{\Money{100}\}$; at $p=1$, $S_{Y_2}=\{\Money{-20}\}$ (USD).
}
\PracticeColumnPair{
\textbf{$Y_1$: portrait sessions.}
\[m_{Y_1}(y)=\begin{cases}\FailureProb{1-p},&y=\Money{-40},\\\SuccessProb{p},&y=\Money{100},\\0,&\text{otherwise}.\end{cases}\]
}{
\textbf{$Y_2$: online CV workshop.}
\[m_{Y_2}(y)=\begin{cases}\FailureProb{1-p},&y=\Money{100},\\\SuccessProb{p},&y=\Money{-20},\\0,&\text{otherwise}.\end{cases}\]
}
\PracticeColumnPair{
\textbf{$Y_1$: portrait sessions.}
\[F_{Y_1}(t)=\begin{cases}0,&t<-40,\\\FailureProb{1-p},&-40\leq t<100,\\1,&t\geq 100.\end{cases}\]
}{
\textbf{$Y_2$: online CV workshop.}
\[F_{Y_2}(t)=\begin{cases}0,&t<-20,\\\SuccessProb{p},&-20\leq t<100,\\1,&t\geq 100.\end{cases}\]
}
Each PMF totals $(1-p)+p=1$. CDF thresholds are numerically ordered,
right-continuous and include each jump at its threshold. At $p=0,1$, zero-mass jumps disappear.
Monetary supports are in USD; indicator values are unitless.
\needspace{8\baselineskip}\paragraph{3. Initial probabilities and loss \protect\HLMathEvent{events}.}
Initially $P(X=0)=\FailureProb{\frac{1}{2}}$ and $P(X=1)=\SuccessProb{\frac{1}{2}}$.
All monetary variables are functions of the same $X$; no independence is asserted.
\PracticeColumnPair{
\textbf{$Y_1$: portrait sessions.}
\[
\begin{gathered}
m_{Y_1}(\Money{-40})=\FailureProb{\frac{1}{2}},\quad m_{Y_1}(\Money{100})=\SuccessProb{\frac{1}{2}},\\
\HLMathEvent{\{Y_1<0\}}=\HLMathEvent{\{X=0\}}=\HLMathEvent{E^c},\\
P(Y_1<0)=\FailureProb{1-p}= \frac{1}{2}\quad\text{initially}.
\end{gathered}
\]
For portrait sessions, loss occurs on $\HLMathEvent{E^c}$: its initial probability per stated period is $\frac{1}{2}$.
}{
\textbf{$Y_2$: online CV workshop.}
\[
\begin{gathered}
m_{Y_2}(\Money{100})=\FailureProb{\frac{1}{2}},\quad m_{Y_2}(\Money{-20})=\SuccessProb{\frac{1}{2}},\\
\HLMathEvent{\{Y_2<0\}}=\HLMathEvent{\{X=1\}}=\HLMathEvent{E},\\
P(Y_2<0)=\SuccessProb{p}= \frac{1}{2}\quad\text{initially}.
\end{gathered}
\]
For online CV workshop, loss occurs on $\HLMathEvent{E}$: its initial probability per stated period is $\frac{1}{2}$.
}
\CThreeModelCoffee
\secInicio[prob:c3-coffee][sec:t2-c3-general]
\clearpage\subsubsection*{Task 2 --- Worked solution}\label{sol:c3-coffee-decision}
\CThreeTaskTwo

Keep the \HLMathEvent{event}, its \HLMathEvent{complement} and the indicator from Task 1.
\needspace{8\baselineskip}\paragraph{1. Initial expected values.}
Use the supplied $\SuccessProb{p}=\SuccessProb{\frac{1}{2}}$ for every alternative.
\PracticeColumnPair{
\textbf{$Y_1$: portrait sessions.}
\[
\begin{aligned}
\ExpectedValue{\mathbb E}[\Money{Y_1}]&=\Money{100}\left(\SuccessProb{\frac{1}{2}}\right)\\
&\quad +(\Money{-40})\left(\FailureProb{\frac{1}{2}}\right)\\
&=\ExpectedValue{30\ \text{USD}}.
\end{aligned}
\]
This is \ExpectedValue{expected net gain} of $\ExpectedValue{30\ \text{USD}}$ per stated period as a long-run average.
}{
\textbf{$Y_2$: online CV workshop.}
\[
\begin{aligned}
\ExpectedValue{\mathbb E}[\Money{Y_2}]&=\Money{-20}\left(\SuccessProb{\frac{1}{2}}\right)\\
&\quad +(\Money{100})\left(\FailureProb{\frac{1}{2}}\right)\\
&=\ExpectedValue{40\ \text{USD}}.
\end{aligned}
\]
This is \ExpectedValue{expected net gain} of $\ExpectedValue{40\ \text{USD}}$ per stated period as a long-run average.
}
\needspace{8\baselineskip}\paragraph{2. Expectation formulas and sign intervals.}
Hold all state payoffs fixed; vary only $p=P(\HLMathEvent{E})\in\left[0,1\right]$.
\PracticeColumnPair{
\textbf{$Y_1$: portrait sessions.}
\[
\begin{aligned}
\ExpectedValue{\mathbb E}[\Money{Y_1}]&=\SuccessProb{p}\Money{f_1(1)}+(\FailureProb{1-p})\Money{f_1(0)}\\
&=\SuccessProb{p}(\Money{100})+(\FailureProb{1-p})(\Money{-40})\\
&=\Money{-40}+\Money{140}\SuccessProb{p}\quad\text{USD}.
\end{aligned}
\]
\[
\begin{aligned}
\ExpectedValue{\mathbb E}[\Money{Y_1}]=0&\iff -40+140p=0\\
&\iff 140p=40\\
&\iff p=p_1^*=\frac{2}{7}.
\end{aligned}
\]
}{
\textbf{$Y_2$: online CV workshop.}
\[
\begin{aligned}
\ExpectedValue{\mathbb E}[\Money{Y_2}]&=\SuccessProb{p}\Money{f_2(1)}+(\FailureProb{1-p})\Money{f_2(0)}\\
&=\SuccessProb{p}(\Money{-20})+(\FailureProb{1-p})(\Money{100})\\
&=\Money{100}-\Money{120}\SuccessProb{p}\quad\text{USD}.
\end{aligned}
\]
\[
\begin{aligned}
\ExpectedValue{\mathbb E}[\Money{Y_2}]=0&\iff 100-120p=0\\
&\iff -120p=-100\\
&\iff p=p_2^*=\frac{5}{6}.
\end{aligned}
\]
}
\PracticeColumnPair{
\textbf{$Y_1$: portrait sessions.}
\[
\begin{aligned}
\ExpectedValue{\mathbb E}[\Money{Y_1}]>0&\iff 140p>40\\
&\iff p>\frac{2}{7}\\
&\iff p\in \left(\frac{2}{7},1\right],\\
\ExpectedValue{\mathbb E}[\Money{Y_1}]<0&\iff 140p<40\\
&\iff p<\frac{2}{7}\\
&\iff p\in \left[0,\frac{2}{7}\right).
\end{aligned}
\]
At $\SuccessProb{p}=\frac{2}{7}$, $\ExpectedValue{\mathbb E}[\Money{Y_1}]$ is \ExpectedValue{zero}: break-even.
}{
\textbf{$Y_2$: online CV workshop.}
\[
\begin{aligned}
\ExpectedValue{\mathbb E}[\Money{Y_2}]>0&\iff -120p>-100\\
&\iff p<\frac{5}{6}\\
&\iff p\in \left[0,\frac{5}{6}\right),\\
\ExpectedValue{\mathbb E}[\Money{Y_2}]<0&\iff -120p<-100\\
&\iff p>\frac{5}{6}\\
&\iff p\in \left(\frac{5}{6},1\right].
\end{aligned}
\]
Dividing by the negative slope reverses both inequalities.
At $\SuccessProb{p}=\frac{5}{6}$, $\ExpectedValue{\mathbb E}[\Money{Y_2}]$ is \ExpectedValue{zero}: break-even.
}
Positive \ExpectedValue{expectation} means \ExpectedValue{expected net gain}; negative
\ExpectedValue{expectation} means \ExpectedValue{expected net loss}; zero is break-even.
\needspace{8\baselineskip}\paragraph{3. Graphs.}
The line joins $(0,-40)$ and $(1,100)$; its zero is $p_1^*=\frac{2}{7}$.
The line joins $(0,100)$ and $(1,-20)$; its zero is $p_2^*=\frac{5}{6}$.

\textbf{Endpoint and root checks.} $Y_1$: $\mathbb E[Y_1](0)=-40$ and $\mathbb E[Y_1](1)=100$ USD; at $p_1^*=\frac{2}{7}$, $-40+(140)(\frac{2}{7})=0$.

$Y_2$: $\mathbb E[Y_2](0)=100$ and $\mathbb E[Y_2](1)=-20$ USD; at $p_2^*=\frac{5}{6}$, $100+(-120)(\frac{5}{6})=0$.

Each observation realizes one state payoff. An \ExpectedValue{expected profit} is a long-run average,
not a guarantee of realized profit. At break-even, zero expectation coexists with nonzero realized payoffs.
\CThreeExpectCoffee
\secInicio[prob:c3-coffee][sec:t2-c3-general]
\clearpage
\subsection{Freelance digital projects for one month}\label{prob:c3-energy}
\CThreePlain{\CThreeStatementEnergy}
\QuestionsTasks[sec:t2-c3-general]{
{3}/{sol:c3-energy-model}/{\CThreePlain{\CThreeTaskOne}
},
{3}/{sol:c3-energy-decision}/{\CThreePlain{\CThreeTaskTwo}
}
}
\needspace{16\baselineskip}\subsubsection*{Task 1 --- Worked solution}\label{sol:c3-energy-model}
\paragraph{Highlighted problem statement.}
\CThreeStatementEnergy
\CThreeTaskOne

\needspace{8\baselineskip}\paragraph{1. Binary model and monetary transformations.}
The observation is a freelancer's digital projects over one month. Let $\HLMathEvent{E}$ be \HLMathEvent{high inquiry volume}; $\HLMathEvent{E^c}$ is \HLMathEvent{low inquiry volume}. Exactly one state occurs. Define $\Indicator{X}=\Indicator{1}$ on $\HLMathEvent{E}$ and $\Indicator{X}=\Indicator{0}$ on $\HLMathEvent{E^c}$. This unitless indicator has $\Indicator{X}\sim\operatorname{Bernoulli}(\SuccessProb{p})$, initially $\SuccessProb{p}=\SuccessProb{\frac{1}{4}}$.
\PracticeColumnPair{
\textbf{$Y_1$: social-media graphics.}
The \Money{monetary consequence} is $\Money{Y_1}=\TransformationFunction{f_1}(\Indicator{X})$ in USD.
\[
\begin{gathered}
f_1(0)=\Money{-30},\quad f_1(1)=\Money{90},\\
\Money{Y_1}=\Money{f_1(0)}+[\Money{f_1(1)}-\Money{f_1(0)}]\Indicator{X}\\
=\Money{-30}+[\Money{90}-(\Money{-30})]\Indicator{X}\\
=\Money{-30}+\Money{120}\Indicator{X}\quad\text{USD}.
\end{gathered}
\]
One payoff is a loss and the other a gain.
}{
\textbf{$Y_2$: short video edits.}
The \Money{monetary consequence} is $\Money{Y_2}=\TransformationFunction{f_2}(\Indicator{X})$ in USD.
\[
\begin{gathered}
f_2(0)=\Money{-60},\quad f_2(1)=\Money{60},\\
\Money{Y_2}=\Money{f_2(0)}+[\Money{f_2(1)}-\Money{f_2(0)}]\Indicator{X}\\
=\Money{-60}+[\Money{60}-(\Money{-60})]\Indicator{X}\\
=\Money{-60}+\Money{120}\Indicator{X}\quad\text{USD}.
\end{gathered}
\]
One payoff is a loss and the other a gain.
}
\needspace{8\baselineskip}\paragraph{2. Supports, PMFs, and ordered CDFs.}

\[S_X=\{x\in\mathbb R:P(X=x)>0\}=\{0,1\}\quad(0<p<1).\]
\[m_X(x)=\begin{cases}\FailureProb{1-p},&x=0,\\\SuccessProb{p},&x=1,\\0,&\text{otherwise},\end{cases}\qquad
F_X(t)=\begin{cases}0,&t<0,\\\FailureProb{1-p},&0\leq t<1,\\1,&t\geq1.\end{cases}\]
At $p=0$, $S_X=\{0\}$; at $p=1$, $S_X=\{1\}$. The formulas remain valid at both endpoints.
\PracticeColumnPair{
\textbf{$Y_1$: social-media graphics.}
\[
\begin{aligned}
S_{Y_1}&=\{y\in\mathbb R:P(Y_1=y)>0\}\\
&=\{\Money{-30},\Money{90}\}\quad(0<p<1).
\end{aligned}
\]
At $p=0$, $S_{Y_1}=\{\Money{-30}\}$; at $p=1$, $S_{Y_1}=\{\Money{90}\}$ (USD).
}{
\textbf{$Y_2$: short video edits.}
\[
\begin{aligned}
S_{Y_2}&=\{y\in\mathbb R:P(Y_2=y)>0\}\\
&=\{\Money{-60},\Money{60}\}\quad(0<p<1).
\end{aligned}
\]
At $p=0$, $S_{Y_2}=\{\Money{-60}\}$; at $p=1$, $S_{Y_2}=\{\Money{60}\}$ (USD).
}
\PracticeColumnPair{
\textbf{$Y_1$: social-media graphics.}
\[m_{Y_1}(y)=\begin{cases}\FailureProb{1-p},&y=\Money{-30},\\\SuccessProb{p},&y=\Money{90},\\0,&\text{otherwise}.\end{cases}\]
}{
\textbf{$Y_2$: short video edits.}
\[m_{Y_2}(y)=\begin{cases}\FailureProb{1-p},&y=\Money{-60},\\\SuccessProb{p},&y=\Money{60},\\0,&\text{otherwise}.\end{cases}\]
}
\PracticeColumnPair{
\textbf{$Y_1$: social-media graphics.}
\[F_{Y_1}(t)=\begin{cases}0,&t<-30,\\\FailureProb{1-p},&-30\leq t<90,\\1,&t\geq 90.\end{cases}\]
}{
\textbf{$Y_2$: short video edits.}
\[F_{Y_2}(t)=\begin{cases}0,&t<-60,\\\FailureProb{1-p},&-60\leq t<60,\\1,&t\geq 60.\end{cases}\]
}
Each PMF totals $(1-p)+p=1$. CDF thresholds are numerically ordered,
right-continuous and include each jump at its threshold. At $p=0,1$, zero-mass jumps disappear.
Monetary supports are in USD; indicator values are unitless.
\needspace{8\baselineskip}\paragraph{3. Initial probabilities and loss \protect\HLMathEvent{events}.}
Initially $P(X=0)=\FailureProb{\frac{3}{4}}$ and $P(X=1)=\SuccessProb{\frac{1}{4}}$.
All monetary variables are functions of the same $X$; no independence is asserted.
\PracticeColumnPair{
\textbf{$Y_1$: social-media graphics.}
\[
\begin{gathered}
m_{Y_1}(\Money{-30})=\FailureProb{\frac{3}{4}},\quad m_{Y_1}(\Money{90})=\SuccessProb{\frac{1}{4}},\\
\HLMathEvent{\{Y_1<0\}}=\HLMathEvent{\{X=0\}}=\HLMathEvent{E^c},\\
P(Y_1<0)=\FailureProb{1-p}= \frac{3}{4}\quad\text{initially}.
\end{gathered}
\]
For social-media graphics, loss occurs on $\HLMathEvent{E^c}$: its initial probability per stated period is $\frac{3}{4}$.
}{
\textbf{$Y_2$: short video edits.}
\[
\begin{gathered}
m_{Y_2}(\Money{-60})=\FailureProb{\frac{3}{4}},\quad m_{Y_2}(\Money{60})=\SuccessProb{\frac{1}{4}},\\
\HLMathEvent{\{Y_2<0\}}=\HLMathEvent{\{X=0\}}=\HLMathEvent{E^c},\\
P(Y_2<0)=\FailureProb{1-p}= \frac{3}{4}\quad\text{initially}.
\end{gathered}
\]
For short video edits, loss occurs on $\HLMathEvent{E^c}$: its initial probability per stated period is $\frac{3}{4}$.
}
\CThreeModelEnergy
\secInicio[prob:c3-energy][sec:t2-c3-general]
\clearpage\subsubsection*{Task 2 --- Worked solution}\label{sol:c3-energy-decision}
\CThreeTaskTwo

Keep the \HLMathEvent{event}, its \HLMathEvent{complement} and the indicator from Task 1.
\needspace{8\baselineskip}\paragraph{1. Initial expected values.}
Use the supplied $\SuccessProb{p}=\SuccessProb{\frac{1}{4}}$ for every alternative.
\PracticeColumnPair{
\textbf{$Y_1$: social-media graphics.}
\[
\begin{aligned}
\ExpectedValue{\mathbb E}[\Money{Y_1}]&=\Money{90}\left(\SuccessProb{\frac{1}{4}}\right)\\
&\quad +(\Money{-30})\left(\FailureProb{\frac{3}{4}}\right)\\
&=\ExpectedValue{0\ \text{USD}}.
\end{aligned}
\]
This is \ExpectedValue{break-even}: zero expected net profit per stated period.
}{
\textbf{$Y_2$: short video edits.}
\[
\begin{aligned}
\ExpectedValue{\mathbb E}[\Money{Y_2}]&=\Money{60}\left(\SuccessProb{\frac{1}{4}}\right)\\
&\quad +(\Money{-60})\left(\FailureProb{\frac{3}{4}}\right)\\
&=\ExpectedValue{-30\ \text{USD}}.
\end{aligned}
\]
This is \ExpectedValue{expected net loss} of $\ExpectedValue{30\ \text{USD}}$ per stated period as a long-run average.
}
\needspace{8\baselineskip}\paragraph{2. Expectation formulas and sign intervals.}
Hold all state payoffs fixed; vary only $p=P(\HLMathEvent{E})\in\left[0,1\right]$.
\PracticeColumnPair{
\textbf{$Y_1$: social-media graphics.}
\[
\begin{aligned}
\ExpectedValue{\mathbb E}[\Money{Y_1}]&=\SuccessProb{p}\Money{f_1(1)}+(\FailureProb{1-p})\Money{f_1(0)}\\
&=\SuccessProb{p}(\Money{90})+(\FailureProb{1-p})(\Money{-30})\\
&=\Money{-30}+\Money{120}\SuccessProb{p}\quad\text{USD}.
\end{aligned}
\]
\[
\begin{aligned}
\ExpectedValue{\mathbb E}[\Money{Y_1}]=0&\iff -30+120p=0\\
&\iff 120p=30\\
&\iff p=p_1^*=\frac{1}{4}.
\end{aligned}
\]
}{
\textbf{$Y_2$: short video edits.}
\[
\begin{aligned}
\ExpectedValue{\mathbb E}[\Money{Y_2}]&=\SuccessProb{p}\Money{f_2(1)}+(\FailureProb{1-p})\Money{f_2(0)}\\
&=\SuccessProb{p}(\Money{60})+(\FailureProb{1-p})(\Money{-60})\\
&=\Money{-60}+\Money{120}\SuccessProb{p}\quad\text{USD}.
\end{aligned}
\]
\[
\begin{aligned}
\ExpectedValue{\mathbb E}[\Money{Y_2}]=0&\iff -60+120p=0\\
&\iff 120p=60\\
&\iff p=p_2^*=\frac{1}{2}.
\end{aligned}
\]
}
\PracticeColumnPair{
\textbf{$Y_1$: social-media graphics.}
\[
\begin{aligned}
\ExpectedValue{\mathbb E}[\Money{Y_1}]>0&\iff 120p>30\\
&\iff p>\frac{1}{4}\\
&\iff p\in \left(\frac{1}{4},1\right],\\
\ExpectedValue{\mathbb E}[\Money{Y_1}]<0&\iff 120p<30\\
&\iff p<\frac{1}{4}\\
&\iff p\in \left[0,\frac{1}{4}\right).
\end{aligned}
\]
At $\SuccessProb{p}=\frac{1}{4}$, $\ExpectedValue{\mathbb E}[\Money{Y_1}]$ is \ExpectedValue{zero}: break-even.
}{
\textbf{$Y_2$: short video edits.}
\[
\begin{aligned}
\ExpectedValue{\mathbb E}[\Money{Y_2}]>0&\iff 120p>60\\
&\iff p>\frac{1}{2}\\
&\iff p\in \left(\frac{1}{2},1\right],\\
\ExpectedValue{\mathbb E}[\Money{Y_2}]<0&\iff 120p<60\\
&\iff p<\frac{1}{2}\\
&\iff p\in \left[0,\frac{1}{2}\right).
\end{aligned}
\]
At $\SuccessProb{p}=\frac{1}{2}$, $\ExpectedValue{\mathbb E}[\Money{Y_2}]$ is \ExpectedValue{zero}: break-even.
}
Positive \ExpectedValue{expectation} means \ExpectedValue{expected net gain}; negative
\ExpectedValue{expectation} means \ExpectedValue{expected net loss}; zero is break-even.
\needspace{8\baselineskip}\paragraph{3. Graphs.}
The line joins $(0,-30)$ and $(1,90)$; its zero is $p_1^*=\frac{1}{4}$.
The line joins $(0,-60)$ and $(1,60)$; its zero is $p_2^*=\frac{1}{2}$.

\textbf{Endpoint and root checks.} $Y_1$: $\mathbb E[Y_1](0)=-30$ and $\mathbb E[Y_1](1)=90$ USD; at $p_1^*=\frac{1}{4}$, $-30+(120)(\frac{1}{4})=0$.

$Y_2$: $\mathbb E[Y_2](0)=-60$ and $\mathbb E[Y_2](1)=60$ USD; at $p_2^*=\frac{1}{2}$, $-60+(120)(\frac{1}{2})=0$.

Each observation realizes one state payoff. An \ExpectedValue{expected profit} is a long-run average,
not a guarantee of realized profit. At break-even, zero expectation coexists with nonzero realized payoffs.
\CThreeExpectEnergy
\secInicio[prob:c3-energy][sec:t2-c3-general]
\clearpage
\subsection{Graduation photography for one quarter}\label{prob:c3-manufacturing}
\CThreePlain{\CThreeStatementManufacturing}
\QuestionsTasks[sec:t2-c3-general]{
{3}/{sol:c3-manufacturing-model}/{\CThreePlain{\CThreeTaskOne}
},
{3}/{sol:c3-manufacturing-decision}/{\CThreePlain{\CThreeTaskTwo}
}
}
\needspace{16\baselineskip}\subsubsection*{Task 1 --- Worked solution}\label{sol:c3-manufacturing-model}
\paragraph{Highlighted problem statement.}
\CThreeStatementManufacturing
\CThreeTaskOne

\needspace{8\baselineskip}\paragraph{1. Binary model and monetary transformations.}
The observation is a graduate's photography work over one quarter. Let $\HLMathEvent{E}$ be \HLMathEvent{the ceremony going ahead}; $\HLMathEvent{E^c}$ is \HLMathEvent{its cancellation}. Exactly one state occurs. Define $\Indicator{X}=\Indicator{1}$ on $\HLMathEvent{E}$ and $\Indicator{X}=\Indicator{0}$ on $\HLMathEvent{E^c}$. This unitless indicator has $\Indicator{X}\sim\operatorname{Bernoulli}(\SuccessProb{p})$, initially $\SuccessProb{p}=\SuccessProb{\frac{3}{5}}$.
\PracticeColumnPair{
\textbf{$Y_1$: basic photo package.}
The \Money{monetary consequence} is $\Money{Y_1}=\TransformationFunction{f_1}(\Indicator{X})$ in USD.
\[
\begin{gathered}
f_1(0)=\Money{-100},\quad f_1(1)=\Money{200},\\
\Money{Y_1}=\Money{f_1(0)}+[\Money{f_1(1)}-\Money{f_1(0)}]\Indicator{X}\\
=\Money{-100}+[\Money{200}-(\Money{-100})]\Indicator{X}\\
=\Money{-100}+\Money{300}\Indicator{X}\quad\text{USD}.
\end{gathered}
\]
One payoff is a loss and the other a gain.
}{
\textbf{$Y_2$: edited photo package.}
The \Money{monetary consequence} is $\Money{Y_2}=\TransformationFunction{f_2}(\Indicator{X})$ in USD.
\[
\begin{gathered}
f_2(0)=\Money{-100},\quad f_2(1)=\Money{500},\\
\Money{Y_2}=\Money{f_2(0)}+[\Money{f_2(1)}-\Money{f_2(0)}]\Indicator{X}\\
=\Money{-100}+[\Money{500}-(\Money{-100})]\Indicator{X}\\
=\Money{-100}+\Money{600}\Indicator{X}\quad\text{USD}.
\end{gathered}
\]
One payoff is a loss and the other a gain.
}
\needspace{8\baselineskip}\paragraph{2. Supports, PMFs, and ordered CDFs.}

\[S_X=\{x\in\mathbb R:P(X=x)>0\}=\{0,1\}\quad(0<p<1).\]
\[m_X(x)=\begin{cases}\FailureProb{1-p},&x=0,\\\SuccessProb{p},&x=1,\\0,&\text{otherwise},\end{cases}\qquad
F_X(t)=\begin{cases}0,&t<0,\\\FailureProb{1-p},&0\leq t<1,\\1,&t\geq1.\end{cases}\]
At $p=0$, $S_X=\{0\}$; at $p=1$, $S_X=\{1\}$. The formulas remain valid at both endpoints.
\PracticeColumnPair{
\textbf{$Y_1$: basic photo package.}
\[
\begin{aligned}
S_{Y_1}&=\{y\in\mathbb R:P(Y_1=y)>0\}\\
&=\{\Money{-100},\Money{200}\}\quad(0<p<1).
\end{aligned}
\]
At $p=0$, $S_{Y_1}=\{\Money{-100}\}$; at $p=1$, $S_{Y_1}=\{\Money{200}\}$ (USD).
}{
\textbf{$Y_2$: edited photo package.}
\[
\begin{aligned}
S_{Y_2}&=\{y\in\mathbb R:P(Y_2=y)>0\}\\
&=\{\Money{-100},\Money{500}\}\quad(0<p<1).
\end{aligned}
\]
At $p=0$, $S_{Y_2}=\{\Money{-100}\}$; at $p=1$, $S_{Y_2}=\{\Money{500}\}$ (USD).
}
\PracticeColumnPair{
\textbf{$Y_1$: basic photo package.}
\[m_{Y_1}(y)=\begin{cases}\FailureProb{1-p},&y=\Money{-100},\\\SuccessProb{p},&y=\Money{200},\\0,&\text{otherwise}.\end{cases}\]
}{
\textbf{$Y_2$: edited photo package.}
\[m_{Y_2}(y)=\begin{cases}\FailureProb{1-p},&y=\Money{-100},\\\SuccessProb{p},&y=\Money{500},\\0,&\text{otherwise}.\end{cases}\]
}
\PracticeColumnPair{
\textbf{$Y_1$: basic photo package.}
\[F_{Y_1}(t)=\begin{cases}0,&t<-100,\\\FailureProb{1-p},&-100\leq t<200,\\1,&t\geq 200.\end{cases}\]
}{
\textbf{$Y_2$: edited photo package.}
\[F_{Y_2}(t)=\begin{cases}0,&t<-100,\\\FailureProb{1-p},&-100\leq t<500,\\1,&t\geq 500.\end{cases}\]
}
Each PMF totals $(1-p)+p=1$. CDF thresholds are numerically ordered,
right-continuous and include each jump at its threshold. At $p=0,1$, zero-mass jumps disappear.
Monetary supports are in USD; indicator values are unitless.
\needspace{8\baselineskip}\paragraph{3. Initial probabilities and loss \protect\HLMathEvent{events}.}
Initially $P(X=0)=\FailureProb{\frac{2}{5}}$ and $P(X=1)=\SuccessProb{\frac{3}{5}}$.
All monetary variables are functions of the same $X$; no independence is asserted.
\PracticeColumnPair{
\textbf{$Y_1$: basic photo package.}
\[
\begin{gathered}
m_{Y_1}(\Money{-100})=\FailureProb{\frac{2}{5}},\quad m_{Y_1}(\Money{200})=\SuccessProb{\frac{3}{5}},\\
\HLMathEvent{\{Y_1<0\}}=\HLMathEvent{\{X=0\}}=\HLMathEvent{E^c},\\
P(Y_1<0)=\FailureProb{1-p}= \frac{2}{5}\quad\text{initially}.
\end{gathered}
\]
For basic photo package, loss occurs on $\HLMathEvent{E^c}$: its initial probability per stated period is $\frac{2}{5}$.
}{
\textbf{$Y_2$: edited photo package.}
\[
\begin{gathered}
m_{Y_2}(\Money{-100})=\FailureProb{\frac{2}{5}},\quad m_{Y_2}(\Money{500})=\SuccessProb{\frac{3}{5}},\\
\HLMathEvent{\{Y_2<0\}}=\HLMathEvent{\{X=0\}}=\HLMathEvent{E^c},\\
P(Y_2<0)=\FailureProb{1-p}= \frac{2}{5}\quad\text{initially}.
\end{gathered}
\]
For edited photo package, loss occurs on $\HLMathEvent{E^c}$: its initial probability per stated period is $\frac{2}{5}$.
}
\CThreeModelManufacturing
\secInicio[prob:c3-manufacturing][sec:t2-c3-general]
\clearpage\subsubsection*{Task 2 --- Worked solution}\label{sol:c3-manufacturing-decision}
\CThreeTaskTwo

Keep the \HLMathEvent{event}, its \HLMathEvent{complement} and the indicator from Task 1.
\needspace{8\baselineskip}\paragraph{1. Initial expected values.}
Use the supplied $\SuccessProb{p}=\SuccessProb{\frac{3}{5}}$ for every alternative.
\PracticeColumnPair{
\textbf{$Y_1$: basic photo package.}
\[
\begin{aligned}
\ExpectedValue{\mathbb E}[\Money{Y_1}]&=\Money{200}\left(\SuccessProb{\frac{3}{5}}\right)\\
&\quad +(\Money{-100})\left(\FailureProb{\frac{2}{5}}\right)\\
&=\ExpectedValue{80\ \text{USD}}.
\end{aligned}
\]
This is \ExpectedValue{expected net gain} of $\ExpectedValue{80\ \text{USD}}$ per stated period as a long-run average.
}{
\textbf{$Y_2$: edited photo package.}
\[
\begin{aligned}
\ExpectedValue{\mathbb E}[\Money{Y_2}]&=\Money{500}\left(\SuccessProb{\frac{3}{5}}\right)\\
&\quad +(\Money{-100})\left(\FailureProb{\frac{2}{5}}\right)\\
&=\ExpectedValue{260\ \text{USD}}.
\end{aligned}
\]
This is \ExpectedValue{expected net gain} of $\ExpectedValue{260\ \text{USD}}$ per stated period as a long-run average.
}
\needspace{8\baselineskip}\paragraph{2. Expectation formulas and sign intervals.}
Hold all state payoffs fixed; vary only $p=P(\HLMathEvent{E})\in\left[0,1\right]$.
\PracticeColumnPair{
\textbf{$Y_1$: basic photo package.}
\[
\begin{aligned}
\ExpectedValue{\mathbb E}[\Money{Y_1}]&=\SuccessProb{p}\Money{f_1(1)}+(\FailureProb{1-p})\Money{f_1(0)}\\
&=\SuccessProb{p}(\Money{200})+(\FailureProb{1-p})(\Money{-100})\\
&=\Money{-100}+\Money{300}\SuccessProb{p}\quad\text{USD}.
\end{aligned}
\]
\[
\begin{aligned}
\ExpectedValue{\mathbb E}[\Money{Y_1}]=0&\iff -100+300p=0\\
&\iff 300p=100\\
&\iff p=p_1^*=\frac{1}{3}.
\end{aligned}
\]
}{
\textbf{$Y_2$: edited photo package.}
\[
\begin{aligned}
\ExpectedValue{\mathbb E}[\Money{Y_2}]&=\SuccessProb{p}\Money{f_2(1)}+(\FailureProb{1-p})\Money{f_2(0)}\\
&=\SuccessProb{p}(\Money{500})+(\FailureProb{1-p})(\Money{-100})\\
&=\Money{-100}+\Money{600}\SuccessProb{p}\quad\text{USD}.
\end{aligned}
\]
\[
\begin{aligned}
\ExpectedValue{\mathbb E}[\Money{Y_2}]=0&\iff -100+600p=0\\
&\iff 600p=100\\
&\iff p=p_2^*=\frac{1}{6}.
\end{aligned}
\]
}
\PracticeColumnPair{
\textbf{$Y_1$: basic photo package.}
\[
\begin{aligned}
\ExpectedValue{\mathbb E}[\Money{Y_1}]>0&\iff 300p>100\\
&\iff p>\frac{1}{3}\\
&\iff p\in \left(\frac{1}{3},1\right],\\
\ExpectedValue{\mathbb E}[\Money{Y_1}]<0&\iff 300p<100\\
&\iff p<\frac{1}{3}\\
&\iff p\in \left[0,\frac{1}{3}\right).
\end{aligned}
\]
At $\SuccessProb{p}=\frac{1}{3}$, $\ExpectedValue{\mathbb E}[\Money{Y_1}]$ is \ExpectedValue{zero}: break-even.
}{
\textbf{$Y_2$: edited photo package.}
\[
\begin{aligned}
\ExpectedValue{\mathbb E}[\Money{Y_2}]>0&\iff 600p>100\\
&\iff p>\frac{1}{6}\\
&\iff p\in \left(\frac{1}{6},1\right],\\
\ExpectedValue{\mathbb E}[\Money{Y_2}]<0&\iff 600p<100\\
&\iff p<\frac{1}{6}\\
&\iff p\in \left[0,\frac{1}{6}\right).
\end{aligned}
\]
At $\SuccessProb{p}=\frac{1}{6}$, $\ExpectedValue{\mathbb E}[\Money{Y_2}]$ is \ExpectedValue{zero}: break-even.
}
Positive \ExpectedValue{expectation} means \ExpectedValue{expected net gain}; negative
\ExpectedValue{expectation} means \ExpectedValue{expected net loss}; zero is break-even.
\needspace{8\baselineskip}\paragraph{3. Graphs.}
The line joins $(0,-100)$ and $(1,200)$; its zero is $p_1^*=\frac{1}{3}$.
The line joins $(0,-100)$ and $(1,500)$; its zero is $p_2^*=\frac{1}{6}$.

\textbf{Endpoint and root checks.} $Y_1$: $\mathbb E[Y_1](0)=-100$ and $\mathbb E[Y_1](1)=200$ USD; at $p_1^*=\frac{1}{3}$, $-100+(300)(\frac{1}{3})=0$.

$Y_2$: $\mathbb E[Y_2](0)=-100$ and $\mathbb E[Y_2](1)=500$ USD; at $p_2^*=\frac{1}{6}$, $-100+(600)(\frac{1}{6})=0$.

Each observation realizes one state payoff. An \ExpectedValue{expected profit} is a long-run average,
not a guarantee of realized profit. At break-even, zero expectation coexists with nonzero realized payoffs.
\CThreeExpectManufacturing
\secInicio[prob:c3-manufacturing][sec:t2-c3-general]
\clearpage
\section{Advanced Problems}\label{sec:t2-c3-advanced}
Each problem applies C3 through one Bernoulli model and analyses each monetary \ExpectedValue{expectation} against zero separately.
\SectionProblemsTableOfContents{\SectionProblemLink{prob:c3-restaurant}{Campus concert: three venue choices, two states}\\
\SectionProblemLink{prob:c3-agriculture}{Student fitness sessions: three break-even thresholds}}
\subsection{Campus concert: three venue choices, two states}\label{prob:c3-restaurant}
\CThreePlain{\CThreeStatementRestaurant}
\QuestionsTasks[sec:t2-c3-advanced]{
{3}/{sol:c3-restaurant-model}/{\CThreePlain{\CThreeTaskOne}
},
{3}/{sol:c3-restaurant-decision}/{\CThreePlain{\CThreeTaskTwo}
}
}
\needspace{16\baselineskip}\subsubsection*{Task 1 --- Worked solution}\label{sol:c3-restaurant-model}
\paragraph{Highlighted problem statement.}
\CThreeStatementRestaurant
\CThreeTaskOne

\needspace{8\baselineskip}\paragraph{1. Binary model and monetary transformations.}
The observation is a student group's ticketed concert over one evening. Let $\HLMathEvent{E}$ be \HLMathEvent{high demand}; $\HLMathEvent{E^c}$ is \HLMathEvent{low demand}. Exactly one state occurs. Define $\Indicator{X}=\Indicator{1}$ on $\HLMathEvent{E}$ and $\Indicator{X}=\Indicator{0}$ on $\HLMathEvent{E^c}$. This unitless indicator has $\Indicator{X}\sim\operatorname{Bernoulli}(\SuccessProb{p})$, initially $\SuccessProb{p}=\SuccessProb{\frac{9}{25}}$.
\PracticeColumnPair{
\textbf{$Y_1$: small hired room.}
The \Money{monetary consequence} is $\Money{Y_1}=\TransformationFunction{f_1}(\Indicator{X})$ in USD.
\[
\begin{gathered}
f_1(0)=\Money{-80},\quad f_1(1)=\Money{100},\\
\Money{Y_1}=\Money{f_1(0)}+[\Money{f_1(1)}-\Money{f_1(0)}]\Indicator{X}\\
=\Money{-80}+[\Money{100}-(\Money{-80})]\Indicator{X}\\
=\Money{-80}+\Money{180}\Indicator{X}\quad\text{USD}.
\end{gathered}
\]
One payoff is a loss and the other a gain.
}{
\textbf{$Y_2$: student-union room.}
The \Money{monetary consequence} is $\Money{Y_2}=\TransformationFunction{f_2}(\Indicator{X})$ in USD.
\[
\begin{gathered}
f_2(0)=\Money{-20},\quad f_2(1)=\Money{220},\\
\Money{Y_2}=\Money{f_2(0)}+[\Money{f_2(1)}-\Money{f_2(0)}]\Indicator{X}\\
=\Money{-20}+[\Money{220}-(\Money{-20})]\Indicator{X}\\
=\Money{-20}+\Money{240}\Indicator{X}\quad\text{USD}.
\end{gathered}
\]
One payoff is a loss and the other a gain.
}
\noindent\begin{minipage}[t]{.48\linewidth}
\textbf{$Y_3$: large hired room.}
The \Money{monetary consequence} is $\Money{Y_3}=\TransformationFunction{f_3}(\Indicator{X})$ in USD.
\[
\begin{gathered}
f_3(0)=\Money{-100},\quad f_3(1)=\Money{400},\\
\Money{Y_3}=\Money{f_3(0)}+[\Money{f_3(1)}-\Money{f_3(0)}]\Indicator{X}\\
=\Money{-100}+[\Money{400}-(\Money{-100})]\Indicator{X}\\
=\Money{-100}+\Money{500}\Indicator{X}\quad\text{USD}.
\end{gathered}
\]
One payoff is a loss and the other a gain.
\end{minipage}\par
\needspace{8\baselineskip}\paragraph{2. Supports, PMFs, and ordered CDFs.}

\[S_X=\{x\in\mathbb R:P(X=x)>0\}=\{0,1\}\quad(0<p<1).\]
\[m_X(x)=\begin{cases}\FailureProb{1-p},&x=0,\\\SuccessProb{p},&x=1,\\0,&\text{otherwise},\end{cases}\qquad
F_X(t)=\begin{cases}0,&t<0,\\\FailureProb{1-p},&0\leq t<1,\\1,&t\geq1.\end{cases}\]
At $p=0$, $S_X=\{0\}$; at $p=1$, $S_X=\{1\}$. The formulas remain valid at both endpoints.
\PracticeColumnPair{
\textbf{$Y_1$: small hired room.}
\[
\begin{aligned}
S_{Y_1}&=\{y\in\mathbb R:P(Y_1=y)>0\}\\
&=\{\Money{-80},\Money{100}\}\quad(0<p<1).
\end{aligned}
\]
At $p=0$, $S_{Y_1}=\{\Money{-80}\}$; at $p=1$, $S_{Y_1}=\{\Money{100}\}$ (USD).
}{
\textbf{$Y_2$: student-union room.}
\[
\begin{aligned}
S_{Y_2}&=\{y\in\mathbb R:P(Y_2=y)>0\}\\
&=\{\Money{-20},\Money{220}\}\quad(0<p<1).
\end{aligned}
\]
At $p=0$, $S_{Y_2}=\{\Money{-20}\}$; at $p=1$, $S_{Y_2}=\{\Money{220}\}$ (USD).
}
\noindent\begin{minipage}[t]{.48\linewidth}
\textbf{$Y_3$: large hired room.}
\[
\begin{aligned}
S_{Y_3}&=\{y\in\mathbb R:P(Y_3=y)>0\}\\
&=\{\Money{-100},\Money{400}\}\quad(0<p<1).
\end{aligned}
\]
At $p=0$, $S_{Y_3}=\{\Money{-100}\}$; at $p=1$, $S_{Y_3}=\{\Money{400}\}$ (USD).
\end{minipage}\par
\PracticeColumnPair{
\textbf{$Y_1$: small hired room.}
\[m_{Y_1}(y)=\begin{cases}\FailureProb{1-p},&y=\Money{-80},\\\SuccessProb{p},&y=\Money{100},\\0,&\text{otherwise}.\end{cases}\]
}{
\textbf{$Y_2$: student-union room.}
\[m_{Y_2}(y)=\begin{cases}\FailureProb{1-p},&y=\Money{-20},\\\SuccessProb{p},&y=\Money{220},\\0,&\text{otherwise}.\end{cases}\]
}
\noindent\begin{minipage}[t]{.48\linewidth}
\textbf{$Y_3$: large hired room.}
\[m_{Y_3}(y)=\begin{cases}\FailureProb{1-p},&y=\Money{-100},\\\SuccessProb{p},&y=\Money{400},\\0,&\text{otherwise}.\end{cases}\]
\end{minipage}\par
\PracticeColumnPair{
\textbf{$Y_1$: small hired room.}
\[F_{Y_1}(t)=\begin{cases}0,&t<-80,\\\FailureProb{1-p},&-80\leq t<100,\\1,&t\geq 100.\end{cases}\]
}{
\textbf{$Y_2$: student-union room.}
\[F_{Y_2}(t)=\begin{cases}0,&t<-20,\\\FailureProb{1-p},&-20\leq t<220,\\1,&t\geq 220.\end{cases}\]
}
\noindent\begin{minipage}[t]{.48\linewidth}
\textbf{$Y_3$: large hired room.}
\[F_{Y_3}(t)=\begin{cases}0,&t<-100,\\\FailureProb{1-p},&-100\leq t<400,\\1,&t\geq 400.\end{cases}\]
\end{minipage}\par
Each PMF totals $(1-p)+p=1$. CDF thresholds are numerically ordered,
right-continuous and include each jump at its threshold. At $p=0,1$, zero-mass jumps disappear.
Monetary supports are in USD; indicator values are unitless.
\clearpage
\needspace{8\baselineskip}\paragraph{3. Initial probabilities and loss \protect\HLMathEvent{events}.}
Initially $P(X=0)=\FailureProb{\frac{16}{25}}$ and $P(X=1)=\SuccessProb{\frac{9}{25}}$.
All monetary variables are functions of the same $X$; no independence is asserted.
\PracticeColumnPair{
\textbf{$Y_1$: small hired room.}
\[
\begin{gathered}
m_{Y_1}(\Money{-80})=\FailureProb{\frac{16}{25}},\quad m_{Y_1}(\Money{100})=\SuccessProb{\frac{9}{25}},\\
\HLMathEvent{\{Y_1<0\}}=\HLMathEvent{\{X=0\}}=\HLMathEvent{E^c},\\
P(Y_1<0)=\FailureProb{1-p}= \frac{16}{25}\quad\text{initially}.
\end{gathered}
\]
For small hired room, loss occurs on $\HLMathEvent{E^c}$: its initial probability per stated period is $\frac{16}{25}$.
}{
\textbf{$Y_2$: student-union room.}
\[
\begin{gathered}
m_{Y_2}(\Money{-20})=\FailureProb{\frac{16}{25}},\quad m_{Y_2}(\Money{220})=\SuccessProb{\frac{9}{25}},\\
\HLMathEvent{\{Y_2<0\}}=\HLMathEvent{\{X=0\}}=\HLMathEvent{E^c},\\
P(Y_2<0)=\FailureProb{1-p}= \frac{16}{25}\quad\text{initially}.
\end{gathered}
\]
For student-union room, loss occurs on $\HLMathEvent{E^c}$: its initial probability per stated period is $\frac{16}{25}$.
}
\noindent\begin{minipage}[t]{.48\linewidth}
\textbf{$Y_3$: large hired room.}
\[
\begin{gathered}
m_{Y_3}(\Money{-100})=\FailureProb{\frac{16}{25}},\quad m_{Y_3}(\Money{400})=\SuccessProb{\frac{9}{25}},\\
\HLMathEvent{\{Y_3<0\}}=\HLMathEvent{\{X=0\}}=\HLMathEvent{E^c},\\
P(Y_3<0)=\FailureProb{1-p}= \frac{16}{25}\quad\text{initially}.
\end{gathered}
\]
For large hired room, loss occurs on $\HLMathEvent{E^c}$: its initial probability per stated period is $\frac{16}{25}$.
\end{minipage}\par
\CThreeModelRestaurant
\secInicio[prob:c3-restaurant][sec:t2-c3-advanced]
\clearpage\subsubsection*{Task 2 --- Worked solution}\label{sol:c3-restaurant-decision}
\CThreeTaskTwo

Keep the \HLMathEvent{event}, its \HLMathEvent{complement} and the indicator from Task 1.
\needspace{8\baselineskip}\paragraph{1. Initial expected values.}
Use the supplied $\SuccessProb{p}=\SuccessProb{\frac{9}{25}}$ for every alternative.
\PracticeColumnPair{
\textbf{$Y_1$: small hired room.}
\[
\begin{aligned}
\ExpectedValue{\mathbb E}[\Money{Y_1}]&=\Money{100}\left(\SuccessProb{\frac{9}{25}}\right)\\
&\quad +(\Money{-80})\left(\FailureProb{\frac{16}{25}}\right)\\
&=\ExpectedValue{\frac{-76}{5}\ \text{USD}}.
\end{aligned}
\]
This is \ExpectedValue{expected net loss} of $\ExpectedValue{\frac{76}{5}\ \text{USD}}$ per stated period as a long-run average.
}{
\textbf{$Y_2$: student-union room.}
\[
\begin{aligned}
\ExpectedValue{\mathbb E}[\Money{Y_2}]&=\Money{220}\left(\SuccessProb{\frac{9}{25}}\right)\\
&\quad +(\Money{-20})\left(\FailureProb{\frac{16}{25}}\right)\\
&=\ExpectedValue{\frac{332}{5}\ \text{USD}}.
\end{aligned}
\]
This is \ExpectedValue{expected net gain} of $\ExpectedValue{\frac{332}{5}\ \text{USD}}$ per stated period as a long-run average.
}
\noindent\begin{minipage}[t]{.48\linewidth}
\textbf{$Y_3$: large hired room.}
\[
\begin{aligned}
\ExpectedValue{\mathbb E}[\Money{Y_3}]&=\Money{400}\left(\SuccessProb{\frac{9}{25}}\right)\\
&\quad +(\Money{-100})\left(\FailureProb{\frac{16}{25}}\right)\\
&=\ExpectedValue{80\ \text{USD}}.
\end{aligned}
\]
This is \ExpectedValue{expected net gain} of $\ExpectedValue{80\ \text{USD}}$ per stated period as a long-run average.
\end{minipage}\par
\needspace{8\baselineskip}\paragraph{2. Expectation formulas and sign intervals.}
Hold all state payoffs fixed; vary only $p=P(\HLMathEvent{E})\in\left[0,1\right]$.
\PracticeColumnPair{
\textbf{$Y_1$: small hired room.}
\[
\begin{aligned}
\ExpectedValue{\mathbb E}[\Money{Y_1}]&=\SuccessProb{p}\Money{f_1(1)}+(\FailureProb{1-p})\Money{f_1(0)}\\
&=\SuccessProb{p}(\Money{100})+(\FailureProb{1-p})(\Money{-80})\\
&=\Money{-80}+\Money{180}\SuccessProb{p}\quad\text{USD}.
\end{aligned}
\]
\[
\begin{aligned}
\ExpectedValue{\mathbb E}[\Money{Y_1}]=0&\iff -80+180p=0\\
&\iff 180p=80\\
&\iff p=p_1^*=\frac{4}{9}.
\end{aligned}
\]
}{
\textbf{$Y_2$: student-union room.}
\[
\begin{aligned}
\ExpectedValue{\mathbb E}[\Money{Y_2}]&=\SuccessProb{p}\Money{f_2(1)}+(\FailureProb{1-p})\Money{f_2(0)}\\
&=\SuccessProb{p}(\Money{220})+(\FailureProb{1-p})(\Money{-20})\\
&=\Money{-20}+\Money{240}\SuccessProb{p}\quad\text{USD}.
\end{aligned}
\]
\[
\begin{aligned}
\ExpectedValue{\mathbb E}[\Money{Y_2}]=0&\iff -20+240p=0\\
&\iff 240p=20\\
&\iff p=p_2^*=\frac{1}{12}.
\end{aligned}
\]
}
\noindent\begin{minipage}[t]{.48\linewidth}
\textbf{$Y_3$: large hired room.}
\[
\begin{aligned}
\ExpectedValue{\mathbb E}[\Money{Y_3}]&=\SuccessProb{p}\Money{f_3(1)}+(\FailureProb{1-p})\Money{f_3(0)}\\
&=\SuccessProb{p}(\Money{400})+(\FailureProb{1-p})(\Money{-100})\\
&=\Money{-100}+\Money{500}\SuccessProb{p}\quad\text{USD}.
\end{aligned}
\]
\[
\begin{aligned}
\ExpectedValue{\mathbb E}[\Money{Y_3}]=0&\iff -100+500p=0\\
&\iff 500p=100\\
&\iff p=p_3^*=\frac{1}{5}.
\end{aligned}
\]
\end{minipage}\par
\clearpage
\textbf{Sign intervals.}\par
\PracticeColumnPair{
\textbf{$Y_1$: small hired room.}
\[
\begin{aligned}
\ExpectedValue{\mathbb E}[\Money{Y_1}]>0&\iff 180p>80\\
&\iff p>\frac{4}{9}\\
&\iff p\in \left(\frac{4}{9},1\right],\\
\ExpectedValue{\mathbb E}[\Money{Y_1}]<0&\iff 180p<80\\
&\iff p<\frac{4}{9}\\
&\iff p\in \left[0,\frac{4}{9}\right).
\end{aligned}
\]
At $\SuccessProb{p}=\frac{4}{9}$, $\ExpectedValue{\mathbb E}[\Money{Y_1}]$ is \ExpectedValue{zero}: break-even.
}{
\textbf{$Y_2$: student-union room.}
\[
\begin{aligned}
\ExpectedValue{\mathbb E}[\Money{Y_2}]>0&\iff 240p>20\\
&\iff p>\frac{1}{12}\\
&\iff p\in \left(\frac{1}{12},1\right],\\
\ExpectedValue{\mathbb E}[\Money{Y_2}]<0&\iff 240p<20\\
&\iff p<\frac{1}{12}\\
&\iff p\in \left[0,\frac{1}{12}\right).
\end{aligned}
\]
At $\SuccessProb{p}=\frac{1}{12}$, $\ExpectedValue{\mathbb E}[\Money{Y_2}]$ is \ExpectedValue{zero}: break-even.
}
\noindent\begin{minipage}[t]{.48\linewidth}
\textbf{$Y_3$: large hired room.}
\[
\begin{aligned}
\ExpectedValue{\mathbb E}[\Money{Y_3}]>0&\iff 500p>100\\
&\iff p>\frac{1}{5}\\
&\iff p\in \left(\frac{1}{5},1\right],\\
\ExpectedValue{\mathbb E}[\Money{Y_3}]<0&\iff 500p<100\\
&\iff p<\frac{1}{5}\\
&\iff p\in \left[0,\frac{1}{5}\right).
\end{aligned}
\]
At $\SuccessProb{p}=\frac{1}{5}$, $\ExpectedValue{\mathbb E}[\Money{Y_3}]$ is \ExpectedValue{zero}: break-even.
\end{minipage}\par
Positive \ExpectedValue{expectation} means \ExpectedValue{expected net gain}; negative
\ExpectedValue{expectation} means \ExpectedValue{expected net loss}; zero is break-even.
\needspace{8\baselineskip}\paragraph{3. Graphs.}
The line joins $(0,-80)$ and $(1,100)$; its zero is $p_1^*=\frac{4}{9}$.
The line joins $(0,-20)$ and $(1,220)$; its zero is $p_2^*=\frac{1}{12}$.
The line joins $(0,-100)$ and $(1,400)$; its zero is $p_3^*=\frac{1}{5}$.

\textbf{Endpoint and root checks.} $Y_1$: $\mathbb E[Y_1](0)=-80$ and $\mathbb E[Y_1](1)=100$ USD; at $p_1^*=\frac{4}{9}$, $-80+(180)(\frac{4}{9})=0$.

$Y_2$: $\mathbb E[Y_2](0)=-20$ and $\mathbb E[Y_2](1)=220$ USD; at $p_2^*=\frac{1}{12}$, $-20+(240)(\frac{1}{12})=0$.

$Y_3$: $\mathbb E[Y_3](0)=-100$ and $\mathbb E[Y_3](1)=400$ USD; at $p_3^*=\frac{1}{5}$, $-100+(500)(\frac{1}{5})=0$.

Each observation realizes one state payoff. An \ExpectedValue{expected profit} is a long-run average,
not a guarantee of realized profit. At break-even, zero expectation coexists with nonzero realized payoffs.
\CThreeExpectRestaurant
\secInicio[prob:c3-restaurant][sec:t2-c3-advanced]
\clearpage
\subsection{Student fitness sessions: three break-even thresholds}\label{prob:c3-agriculture}
\CThreePlain{\CThreeStatementAgriculture}
\QuestionsTasks[sec:t2-c3-advanced]{
{3}/{sol:c3-agriculture-model}/{\CThreePlain{\CThreeTaskOne}
},
{3}/{sol:c3-agriculture-decision}/{\CThreePlain{\CThreeTaskTwo}
}
}
\needspace{16\baselineskip}\subsubsection*{Task 1 --- Worked solution}\label{sol:c3-agriculture-model}
\paragraph{Highlighted problem statement.}
\CThreeStatementAgriculture
\CThreeTaskOne

\needspace{8\baselineskip}\paragraph{1. Binary model and monetary transformations.}
The observation is a student's paid fitness sessions over one university term. Let $\HLMathEvent{E}$ be \HLMathEvent{high enrollment}; $\HLMathEvent{E^c}$ is \HLMathEvent{low enrollment}. Exactly one state occurs. Define $\Indicator{X}=\Indicator{1}$ on $\HLMathEvent{E}$ and $\Indicator{X}=\Indicator{0}$ on $\HLMathEvent{E^c}$. This unitless indicator has $\Indicator{X}\sim\operatorname{Bernoulli}(\SuccessProb{p})$, initially $\SuccessProb{p}=\SuccessProb{\frac{1}{2}}$.
\PracticeColumnPair{
\textbf{$Y_1$: walking group.}
The \Money{monetary consequence} is $\Money{Y_1}=\TransformationFunction{f_1}(\Indicator{X})$ in USD.
\[
\begin{gathered}
f_1(0)=\Money{-70},\quad f_1(1)=\Money{90},\\
\Money{Y_1}=\Money{f_1(0)}+[\Money{f_1(1)}-\Money{f_1(0)}]\Indicator{X}\\
=\Money{-70}+[\Money{90}-(\Money{-70})]\Indicator{X}\\
=\Money{-70}+\Money{160}\Indicator{X}\quad\text{USD}.
\end{gathered}
\]
One payoff is a loss and the other a gain.
}{
\textbf{$Y_2$: dance sessions.}
The \Money{monetary consequence} is $\Money{Y_2}=\TransformationFunction{f_2}(\Indicator{X})$ in USD.
\[
\begin{gathered}
f_2(0)=\Money{-30},\quad f_2(1)=\Money{190},\\
\Money{Y_2}=\Money{f_2(0)}+[\Money{f_2(1)}-\Money{f_2(0)}]\Indicator{X}\\
=\Money{-30}+[\Money{190}-(\Money{-30})]\Indicator{X}\\
=\Money{-30}+\Money{220}\Indicator{X}\quad\text{USD}.
\end{gathered}
\]
One payoff is a loss and the other a gain.
}
\noindent\begin{minipage}[t]{.48\linewidth}
\textbf{$Y_3$: circuit sessions.}
The \Money{monetary consequence} is $\Money{Y_3}=\TransformationFunction{f_3}(\Indicator{X})$ in USD.
\[
\begin{gathered}
f_3(0)=\Money{-150},\quad f_3(1)=\Money{150},\\
\Money{Y_3}=\Money{f_3(0)}+[\Money{f_3(1)}-\Money{f_3(0)}]\Indicator{X}\\
=\Money{-150}+[\Money{150}-(\Money{-150})]\Indicator{X}\\
=\Money{-150}+\Money{300}\Indicator{X}\quad\text{USD}.
\end{gathered}
\]
One payoff is a loss and the other a gain.
\end{minipage}\par
\needspace{8\baselineskip}\paragraph{2. Supports, PMFs, and ordered CDFs.}

\[S_X=\{x\in\mathbb R:P(X=x)>0\}=\{0,1\}\quad(0<p<1).\]
\[m_X(x)=\begin{cases}\FailureProb{1-p},&x=0,\\\SuccessProb{p},&x=1,\\0,&\text{otherwise},\end{cases}\qquad
F_X(t)=\begin{cases}0,&t<0,\\\FailureProb{1-p},&0\leq t<1,\\1,&t\geq1.\end{cases}\]
At $p=0$, $S_X=\{0\}$; at $p=1$, $S_X=\{1\}$. The formulas remain valid at both endpoints.
\PracticeColumnPair{
\textbf{$Y_1$: walking group.}
\[
\begin{aligned}
S_{Y_1}&=\{y\in\mathbb R:P(Y_1=y)>0\}\\
&=\{\Money{-70},\Money{90}\}\quad(0<p<1).
\end{aligned}
\]
At $p=0$, $S_{Y_1}=\{\Money{-70}\}$; at $p=1$, $S_{Y_1}=\{\Money{90}\}$ (USD).
}{
\textbf{$Y_2$: dance sessions.}
\[
\begin{aligned}
S_{Y_2}&=\{y\in\mathbb R:P(Y_2=y)>0\}\\
&=\{\Money{-30},\Money{190}\}\quad(0<p<1).
\end{aligned}
\]
At $p=0$, $S_{Y_2}=\{\Money{-30}\}$; at $p=1$, $S_{Y_2}=\{\Money{190}\}$ (USD).
}
\noindent\begin{minipage}[t]{.48\linewidth}
\textbf{$Y_3$: circuit sessions.}
\[
\begin{aligned}
S_{Y_3}&=\{y\in\mathbb R:P(Y_3=y)>0\}\\
&=\{\Money{-150},\Money{150}\}\quad(0<p<1).
\end{aligned}
\]
At $p=0$, $S_{Y_3}=\{\Money{-150}\}$; at $p=1$, $S_{Y_3}=\{\Money{150}\}$ (USD).
\end{minipage}\par
\PracticeColumnPair{
\textbf{$Y_1$: walking group.}
\[m_{Y_1}(y)=\begin{cases}\FailureProb{1-p},&y=\Money{-70},\\\SuccessProb{p},&y=\Money{90},\\0,&\text{otherwise}.\end{cases}\]
}{
\textbf{$Y_2$: dance sessions.}
\[m_{Y_2}(y)=\begin{cases}\FailureProb{1-p},&y=\Money{-30},\\\SuccessProb{p},&y=\Money{190},\\0,&\text{otherwise}.\end{cases}\]
}
\noindent\begin{minipage}[t]{.48\linewidth}
\textbf{$Y_3$: circuit sessions.}
\[m_{Y_3}(y)=\begin{cases}\FailureProb{1-p},&y=\Money{-150},\\\SuccessProb{p},&y=\Money{150},\\0,&\text{otherwise}.\end{cases}\]
\end{minipage}\par
\PracticeColumnPair{
\textbf{$Y_1$: walking group.}
\[F_{Y_1}(t)=\begin{cases}0,&t<-70,\\\FailureProb{1-p},&-70\leq t<90,\\1,&t\geq 90.\end{cases}\]
}{
\textbf{$Y_2$: dance sessions.}
\[F_{Y_2}(t)=\begin{cases}0,&t<-30,\\\FailureProb{1-p},&-30\leq t<190,\\1,&t\geq 190.\end{cases}\]
}
\noindent\begin{minipage}[t]{.48\linewidth}
\textbf{$Y_3$: circuit sessions.}
\[F_{Y_3}(t)=\begin{cases}0,&t<-150,\\\FailureProb{1-p},&-150\leq t<150,\\1,&t\geq 150.\end{cases}\]
\end{minipage}\par
Each PMF totals $(1-p)+p=1$. CDF thresholds are numerically ordered,
right-continuous and include each jump at its threshold. At $p=0,1$, zero-mass jumps disappear.
Monetary supports are in USD; indicator values are unitless.
\clearpage
\needspace{8\baselineskip}\paragraph{3. Initial probabilities and loss \protect\HLMathEvent{events}.}
Initially $P(X=0)=\FailureProb{\frac{1}{2}}$ and $P(X=1)=\SuccessProb{\frac{1}{2}}$.
All monetary variables are functions of the same $X$; no independence is asserted.
\PracticeColumnPair{
\textbf{$Y_1$: walking group.}
\[
\begin{gathered}
m_{Y_1}(\Money{-70})=\FailureProb{\frac{1}{2}},\quad m_{Y_1}(\Money{90})=\SuccessProb{\frac{1}{2}},\\
\HLMathEvent{\{Y_1<0\}}=\HLMathEvent{\{X=0\}}=\HLMathEvent{E^c},\\
P(Y_1<0)=\FailureProb{1-p}= \frac{1}{2}\quad\text{initially}.
\end{gathered}
\]
For walking group, loss occurs on $\HLMathEvent{E^c}$: its initial probability per stated period is $\frac{1}{2}$.
}{
\textbf{$Y_2$: dance sessions.}
\[
\begin{gathered}
m_{Y_2}(\Money{-30})=\FailureProb{\frac{1}{2}},\quad m_{Y_2}(\Money{190})=\SuccessProb{\frac{1}{2}},\\
\HLMathEvent{\{Y_2<0\}}=\HLMathEvent{\{X=0\}}=\HLMathEvent{E^c},\\
P(Y_2<0)=\FailureProb{1-p}= \frac{1}{2}\quad\text{initially}.
\end{gathered}
\]
For dance sessions, loss occurs on $\HLMathEvent{E^c}$: its initial probability per stated period is $\frac{1}{2}$.
}
\noindent\begin{minipage}[t]{.48\linewidth}
\textbf{$Y_3$: circuit sessions.}
\[
\begin{gathered}
m_{Y_3}(\Money{-150})=\FailureProb{\frac{1}{2}},\quad m_{Y_3}(\Money{150})=\SuccessProb{\frac{1}{2}},\\
\HLMathEvent{\{Y_3<0\}}=\HLMathEvent{\{X=0\}}=\HLMathEvent{E^c},\\
P(Y_3<0)=\FailureProb{1-p}= \frac{1}{2}\quad\text{initially}.
\end{gathered}
\]
For circuit sessions, loss occurs on $\HLMathEvent{E^c}$: its initial probability per stated period is $\frac{1}{2}$.
\end{minipage}\par
\CThreeModelAgriculture
\secInicio[prob:c3-agriculture][sec:t2-c3-advanced]
\clearpage\subsubsection*{Task 2 --- Worked solution}\label{sol:c3-agriculture-decision}
\CThreeTaskTwo

Keep the \HLMathEvent{event}, its \HLMathEvent{complement} and the indicator from Task 1.
\needspace{8\baselineskip}\paragraph{1. Initial expected values.}
Use the supplied $\SuccessProb{p}=\SuccessProb{\frac{1}{2}}$ for every alternative.
\PracticeColumnPair{
\textbf{$Y_1$: walking group.}
\[
\begin{aligned}
\ExpectedValue{\mathbb E}[\Money{Y_1}]&=\Money{90}\left(\SuccessProb{\frac{1}{2}}\right)\\
&\quad +(\Money{-70})\left(\FailureProb{\frac{1}{2}}\right)\\
&=\ExpectedValue{10\ \text{USD}}.
\end{aligned}
\]
This is \ExpectedValue{expected net gain} of $\ExpectedValue{10\ \text{USD}}$ per stated period as a long-run average.
}{
\textbf{$Y_2$: dance sessions.}
\[
\begin{aligned}
\ExpectedValue{\mathbb E}[\Money{Y_2}]&=\Money{190}\left(\SuccessProb{\frac{1}{2}}\right)\\
&\quad +(\Money{-30})\left(\FailureProb{\frac{1}{2}}\right)\\
&=\ExpectedValue{80\ \text{USD}}.
\end{aligned}
\]
This is \ExpectedValue{expected net gain} of $\ExpectedValue{80\ \text{USD}}$ per stated period as a long-run average.
}
\noindent\begin{minipage}[t]{.48\linewidth}
\textbf{$Y_3$: circuit sessions.}
\[
\begin{aligned}
\ExpectedValue{\mathbb E}[\Money{Y_3}]&=\Money{150}\left(\SuccessProb{\frac{1}{2}}\right)\\
&\quad +(\Money{-150})\left(\FailureProb{\frac{1}{2}}\right)\\
&=\ExpectedValue{0\ \text{USD}}.
\end{aligned}
\]
This is \ExpectedValue{break-even}: zero expected net profit per stated period.
\end{minipage}\par
\needspace{8\baselineskip}\paragraph{2. Expectation formulas and sign intervals.}
Hold all state payoffs fixed; vary only $p=P(\HLMathEvent{E})\in\left[0,1\right]$.
\PracticeColumnPair{
\textbf{$Y_1$: walking group.}
\[
\begin{aligned}
\ExpectedValue{\mathbb E}[\Money{Y_1}]&=\SuccessProb{p}\Money{f_1(1)}+(\FailureProb{1-p})\Money{f_1(0)}\\
&=\SuccessProb{p}(\Money{90})+(\FailureProb{1-p})(\Money{-70})\\
&=\Money{-70}+\Money{160}\SuccessProb{p}\quad\text{USD}.
\end{aligned}
\]
\[
\begin{aligned}
\ExpectedValue{\mathbb E}[\Money{Y_1}]=0&\iff -70+160p=0\\
&\iff 160p=70\\
&\iff p=p_1^*=\frac{7}{16}.
\end{aligned}
\]
}{
\textbf{$Y_2$: dance sessions.}
\[
\begin{aligned}
\ExpectedValue{\mathbb E}[\Money{Y_2}]&=\SuccessProb{p}\Money{f_2(1)}+(\FailureProb{1-p})\Money{f_2(0)}\\
&=\SuccessProb{p}(\Money{190})+(\FailureProb{1-p})(\Money{-30})\\
&=\Money{-30}+\Money{220}\SuccessProb{p}\quad\text{USD}.
\end{aligned}
\]
\[
\begin{aligned}
\ExpectedValue{\mathbb E}[\Money{Y_2}]=0&\iff -30+220p=0\\
&\iff 220p=30\\
&\iff p=p_2^*=\frac{3}{22}.
\end{aligned}
\]
}
\noindent\begin{minipage}[t]{.48\linewidth}
\textbf{$Y_3$: circuit sessions.}
\[
\begin{aligned}
\ExpectedValue{\mathbb E}[\Money{Y_3}]&=\SuccessProb{p}\Money{f_3(1)}+(\FailureProb{1-p})\Money{f_3(0)}\\
&=\SuccessProb{p}(\Money{150})+(\FailureProb{1-p})(\Money{-150})\\
&=\Money{-150}+\Money{300}\SuccessProb{p}\quad\text{USD}.
\end{aligned}
\]
\[
\begin{aligned}
\ExpectedValue{\mathbb E}[\Money{Y_3}]=0&\iff -150+300p=0\\
&\iff 300p=150\\
&\iff p=p_3^*=\frac{1}{2}.
\end{aligned}
\]
\end{minipage}\par
\clearpage
\textbf{Sign intervals.}\par
\PracticeColumnPair{
\textbf{$Y_1$: walking group.}
\[
\begin{aligned}
\ExpectedValue{\mathbb E}[\Money{Y_1}]>0&\iff 160p>70\\
&\iff p>\frac{7}{16}\\
&\iff p\in \left(\frac{7}{16},1\right],\\
\ExpectedValue{\mathbb E}[\Money{Y_1}]<0&\iff 160p<70\\
&\iff p<\frac{7}{16}\\
&\iff p\in \left[0,\frac{7}{16}\right).
\end{aligned}
\]
At $\SuccessProb{p}=\frac{7}{16}$, $\ExpectedValue{\mathbb E}[\Money{Y_1}]$ is \ExpectedValue{zero}: break-even.
}{
\textbf{$Y_2$: dance sessions.}
\[
\begin{aligned}
\ExpectedValue{\mathbb E}[\Money{Y_2}]>0&\iff 220p>30\\
&\iff p>\frac{3}{22}\\
&\iff p\in \left(\frac{3}{22},1\right],\\
\ExpectedValue{\mathbb E}[\Money{Y_2}]<0&\iff 220p<30\\
&\iff p<\frac{3}{22}\\
&\iff p\in \left[0,\frac{3}{22}\right).
\end{aligned}
\]
At $\SuccessProb{p}=\frac{3}{22}$, $\ExpectedValue{\mathbb E}[\Money{Y_2}]$ is \ExpectedValue{zero}: break-even.
}
\noindent\begin{minipage}[t]{.48\linewidth}
\textbf{$Y_3$: circuit sessions.}
\[
\begin{aligned}
\ExpectedValue{\mathbb E}[\Money{Y_3}]>0&\iff 300p>150\\
&\iff p>\frac{1}{2}\\
&\iff p\in \left(\frac{1}{2},1\right],\\
\ExpectedValue{\mathbb E}[\Money{Y_3}]<0&\iff 300p<150\\
&\iff p<\frac{1}{2}\\
&\iff p\in \left[0,\frac{1}{2}\right).
\end{aligned}
\]
At $\SuccessProb{p}=\frac{1}{2}$, $\ExpectedValue{\mathbb E}[\Money{Y_3}]$ is \ExpectedValue{zero}: break-even.
\end{minipage}\par
Positive \ExpectedValue{expectation} means \ExpectedValue{expected net gain}; negative
\ExpectedValue{expectation} means \ExpectedValue{expected net loss}; zero is break-even.
\needspace{8\baselineskip}\paragraph{3. Graphs.}
The line joins $(0,-70)$ and $(1,90)$; its zero is $p_1^*=\frac{7}{16}$.
The line joins $(0,-30)$ and $(1,190)$; its zero is $p_2^*=\frac{3}{22}$.
The line joins $(0,-150)$ and $(1,150)$; its zero is $p_3^*=\frac{1}{2}$.

\textbf{Endpoint and root checks.} $Y_1$: $\mathbb E[Y_1](0)=-70$ and $\mathbb E[Y_1](1)=90$ USD; at $p_1^*=\frac{7}{16}$, $-70+(160)(\frac{7}{16})=0$.

$Y_2$: $\mathbb E[Y_2](0)=-30$ and $\mathbb E[Y_2](1)=190$ USD; at $p_2^*=\frac{3}{22}$, $-30+(220)(\frac{3}{22})=0$.

$Y_3$: $\mathbb E[Y_3](0)=-150$ and $\mathbb E[Y_3](1)=150$ USD; at $p_3^*=\frac{1}{2}$, $-150+(300)(\frac{1}{2})=0$.

Each observation realizes one state payoff. An \ExpectedValue{expected profit} is a long-run average,
not a guarantee of realized profit. At break-even, zero expectation coexists with nonzero realized payoffs.
\CThreeExpectAgriculture
\secInicio[prob:c3-agriculture][sec:t2-c3-advanced]

\clearpage\subsection*{Closing reminders: Bernoulli modelling and individual break-even analysis}
Define one \HLMathEvent{event} $\HLMathEvent{E}$, its \HLMathEvent{complement} $\HLMathEvent{E^c}$, and keep $X=1$ on $\HLMathEvent{E}$, $X=0$ on $\HLMathEvent{E^c}$.
The indicator is unitless; $p=P(\HLMathEvent{E})$. Each $Y_i=f_i(X)$ is in USD, with one
strictly positive and one strictly negative monetary payoff. These signs refer
to money, not to the \HLMathEvent{event} or its indicator. Show both state mappings and
$Y_i=f_i(0)+[f_i(1)-f_i(0)]X$.
Use supports with positive probability; omit zero-mass values at $p=0,1$.
Distinguish the transformation, PMF and CDF; sort amounts numerically for CDFs.
For each $Y_i$ separately, keep both payoffs fixed, vary only $p\in\left[0,1\right]$,
derive its \ExpectedValue{expectation}, solve its own zero equation and both strict inequalities,
and include the applicable interval endpoints. Positive \ExpectedValue{expectation} is expected
net gain, negative is \ExpectedValue{expected net loss} and zero is break-even. Evaluate the initial
probability and provide a separate labelled graph with the zero line and exact root.
Separate analyses do not imply independence: all variables share $X$.
An \ExpectedValue{expected value} is a long-run average, not a guarantee of realized profit.
\secInicio[][sec:t2-c3-advanced]
\end{document}
